% Application de la biotechnologie dans la production de biens et services — SVT 1ère A — Cours signé OnBuch+
% Programme officiel MINESEC. Compiler avec : tectonic NA08-application-biotechnologie-biens-services.tex
\newcommand{\DOCMATIERE}{SVT}
\newcommand{\DOCNIVEAU}{1ère A}
\newcommand{\DOCMODULE}{Module 4 : La Biotechnologie}
\newcommand{\DOCLECON}{NA08}
\newcommand{\DOCTITRE}{Application de la biotechnologie dans la production de biens et services}
% OnBuch+ — préambule « cours signés » ENRICHI (compilé avec tectonic / XeLaTeX)
% Système visuel « Pop » d'OnBuch+ : fond crème (jamais blanc), bordures encre
% épaisses, ombres « dures » décalées, titres Archivo Black en capitales.
% Version MATHÉMATIQUES : graphes (pgfplots/tikz), boîtes pédagogiques
% (définition, propriété, méthode, attention, le savais-tu, exercice résolu,
% exercice, TP, bilan), page de garde et signature OnBuch+ ancrée en bas.
%
% Macros attendues dans main.tex AVANT \input{preamble} :
%   \DOCMATIERE  \DOCNIVEAU  \DOCMODULE  \DOCLECON (numéro)  \DOCTITRE
\documentclass[11pt]{article}

\usepackage{fontspec}
\usepackage[french]{babel}
\usepackage[a4paper,margin=2cm,top=3.4cm,bottom=2.8cm,headheight=1.4cm,footskip=1.3cm]{geometry}
\usepackage{amsmath,amssymb,amsfonts}
\usepackage{mathtools} % \xmapsto, \coloneqq... souvent utilisés par les modèles
\usepackage{cancel} % simplifications barrées (bilans de forces)
% \abs{z} (module, valeur absolue) et \norm{u} (norme), utilisés par les modèles
\providecommand{\abs}{}\let\abs\relax\DeclarePairedDelimiter{\abs}{\lvert}{\rvert}
\providecommand{\norm}{}\let\norm\relax\DeclarePairedDelimiter{\norm}{\lVert}{\rVert}
\usepackage[table]{xcolor} % [table] : fournit aussi \rowcolors (alternance de lignes)
\usepackage{pagecolor}
\usepackage{tikz}
\usetikzlibrary{arrows.meta,positioning,calc,shapes.geometric,shapes.misc,shapes.arrows,decorations.pathmorphing,decorations.pathreplacing,decorations.markings,patterns,shadows,mindmap,trees,fit,backgrounds,matrix,intersections,angles,quotes}
\usepackage{pgfplots}
\usepackage[version=4]{mhchem} % équations nucléaires \ce{^{235}_{92}U}
\usepackage[european,siunitx]{circuitikz} % schémas électriques
\pgfplotsset{compat=1.18}
\usepgfplotslibrary{fillbetween,groupplots} % aires entre courbes (calcul intégral)
\usepackage{enumitem}
\usepackage{fancyhdr}
% [most] charge toutes les bibliothèques tcolorbox (listings, minted, raster,
% documentation...) et gonfle inutilement la mémoire TeX sur un document long
% (~300 boîtes) : on ne charge que ce qui est réellement utilisé.
\usepackage[skins,breakable]{tcolorbox}
\usepackage{graphicx}
\usepackage{colortbl}
\usepackage{booktabs}
\usepackage{tabularx}
\usepackage{diagbox} % case d'en-tête diagonale des tableaux à double entrée
% Colonnes extensibles centrée / gauche / droite (C, L, R), souvent utilisées par les modèles
\newcolumntype{C}{>{\centering\arraybackslash}X}
\newcolumntype{L}{>{\raggedright\arraybackslash}X}
\newcolumntype{R}{>{\raggedleft\arraybackslash}X}
\usepackage{array}
\usepackage{multicol}
\usepackage{siunitx}
% parse-numbers=false : accepte aussi \SI{2,5 \times 10^{-3}}{...}
\sisetup{locale=FR,output-decimal-marker={,},group-minimum-digits=4,parse-numbers=false}
\DeclareSIUnit\ohm{\ensuremath{\symup{\Omega}}} % U+2126 absent de la police
\DeclareSIUnit\division{div} % réglages d'oscilloscope (ms/div, V/div)
\DeclareSIUnit\tr{tr} % tours (vitesse de rotation, ex. \SI{1500}{\tr\per\minute})
\DeclareSIUnit\ch{ch} % cheval-vapeur (puissance moteur, unité usuelle hors SI)
\DeclareSIUnit\franccfa{FCFA} % franc CFA (coûts, contexte camerounais)
\DeclareSIUnit\dioptre{\ensuremath{\delta}} % vergence d'une lentille (dioptrie)
\usepackage{qrcode}
% \degree : commande fréquemment utilisée par les modèles pour un angle ou une
% température en dehors de \SI{}{} (non fournie par les paquets chargés).
\providecommand{\degree}{^{\circ}}
% \qed (fin de démonstration), utilisé par les modèles sans amsthm
\providecommand{\qed}{\hfill$\square$}
% \barre : double barre des valeurs interdites dans un tableau de variations (tkz-tab)
\providecommand{\barre}{\ensuremath{\|}}
% environnement proof (amsthm non chargé) utilisé par les modèles
\ifdefined\proof\else\newenvironment{proof}[1][Démonstration]{\par\noindent\textit{#1.}\ }{\qed\par}\fi
\usepackage{lastpage}
\usepackage{hyperref}
\hypersetup{hidelinks}

% ============================================================
% Palette « Pop » OnBuch+ — mêmes hex que la classe P (plus_ui.dart)
% ============================================================
\definecolor{poporange}{HTML}{F59321}
\definecolor{popdark}{HTML}{E07A0C}
\definecolor{popcream}{HTML}{FFF6E8}
\definecolor{popink}{HTML}{1C1714}
\definecolor{poppurple}{HTML}{7A5AE0}
\definecolor{popgreen}{HTML}{2E9E6A}
\definecolor{poppink}{HTML}{E0457B}
\definecolor{popblue}{HTML}{2F7DE1}
\definecolor{popgold}{HTML}{C98A12}
\definecolor{popmuted}{HTML}{8A7F76}
\definecolor{popline}{HTML}{EADFD2}
\definecolor{popgray}{HTML}{8A7F76}
\colorlet{popgrey}{popgray}
% Alias tolérants (le modèle nomme parfois les couleurs différemment)
\colorlet{popwhite}{white}
\colorlet{popblack}{popink}
\colorlet{popred}{poppink}
\colorlet{popyellow}{popgold}
% Teintes claires pour fonds de tableaux / zones
\colorlet{popblueL}{popblue!10!white}
\colorlet{poporangeL}{poporange!14!white}
\colorlet{popgreenL}{popgreen!12!white}
\colorlet{poppurpleL}{poppurple!12!white}
\colorlet{poppinkL}{poppink!12!white}
\colorlet{popgoldL}{popgold!14!white}

\pagecolor{popcream}

% ============================================================
% Typographie
% ============================================================
\defaultfontfeatures{Ligatures=TeX}
\setmainfont{PlusJakartaSans-Medium.ttf}[Path=../../../fonts/,BoldFont=PlusJakartaSans-Bold.ttf]
% Maths en sans-serif (Fira Math) avec chiffres et lettres droites de Plus
% Jakarta, pour que les formules se fondent dans le texte.
\usepackage[math-style=ISO,bold-style=ISO]{unicode-math}
\setmathfont{FiraMath-Regular.otf}[Path=../../../fonts/]
\setmathfont{PlusJakartaSans-Medium.ttf}[Path=../../../fonts/,range={up/num,up/{Latin,latin},\mathup}]
\usepackage{icomma} % virgule décimale sans espace en mode math
% Caractères mathématiques Unicode absents des polices de texte (Plus Jakarta,
% Archivo Black) : rendus par leur équivalent mathématique, en texte comme en
% formule — sinon ils s'affichent en carré vide (ex. « Arithmétique dans ℤ »).
\usepackage{newunicodechar}
\newunicodechar{ℝ}{\ensuremath{\mathbb{R}}}
\newunicodechar{ℤ}{\ensuremath{\mathbb{Z}}}
\newunicodechar{ℂ}{\ensuremath{\mathbb{C}}}
\newunicodechar{ℕ}{\ensuremath{\mathbb{N}}}
\newunicodechar{ℚ}{\ensuremath{\mathbb{Q}}}
\newunicodechar{𝔻}{\ensuremath{\mathbb{D}}}
\newunicodechar{α}{\ensuremath{\alpha}}
\newunicodechar{β}{\ensuremath{\beta}}
\newunicodechar{γ}{\ensuremath{\gamma}}
\newunicodechar{δ}{\ensuremath{\delta}}
\newunicodechar{ε}{\ensuremath{\varepsilon}}
\newunicodechar{θ}{\ensuremath{\theta}}
\newunicodechar{λ}{\ensuremath{\lambda}}
\newunicodechar{μ}{\ensuremath{\mu}}
\newunicodechar{σ}{\ensuremath{\sigma}}
\newunicodechar{τ}{\ensuremath{\tau}}
\newunicodechar{φ}{\ensuremath{\varphi}}
\newunicodechar{ψ}{\ensuremath{\psi}}
\newunicodechar{ω}{\ensuremath{\omega}}
\newunicodechar{Σ}{\ensuremath{\Sigma}}
% majuscules produites par \MakeUppercase dans les titres (α-aminés -> Α)
\newunicodechar{Α}{\ensuremath{\alpha}}
\newunicodechar{Β}{\ensuremath{\beta}}
\newunicodechar{Γ}{\ensuremath{\Gamma}}
\newunicodechar{Δ}{\ensuremath{\Delta}}
\newunicodechar{Ε}{\ensuremath{\varepsilon}}
\newunicodechar{Θ}{\ensuremath{\Theta}}
\newunicodechar{Λ}{\ensuremath{\Lambda}}
\newunicodechar{Μ}{\ensuremath{\mu}}
\newunicodechar{Τ}{\ensuremath{\tau}}
\newunicodechar{Φ}{\ensuremath{\Phi}}
\newunicodechar{Ψ}{\ensuremath{\Psi}}
\newunicodechar{Ω}{\ensuremath{\Omega}}
\newunicodechar{↦}{\ensuremath{\mapsto}}
\newunicodechar{∈}{\ensuremath{\in}}
\newunicodechar{∉}{\ensuremath{\notin}}
\newunicodechar{∧}{\ensuremath{\wedge}}
\newunicodechar{⇔}{\ensuremath{\Leftrightarrow}}
\newunicodechar{⟺}{\ensuremath{\Longleftrightarrow}}
\newunicodechar{⇒}{\ensuremath{\Rightarrow}}
\newunicodechar{⟹}{\ensuremath{\Longrightarrow}}
\newunicodechar{∘}{\ensuremath{\circ}}
\newunicodechar{∪}{\ensuremath{\cup}}
\newunicodechar{∩}{\ensuremath{\cap}}
\newunicodechar{≡}{\ensuremath{\equiv}}
\newunicodechar{⊂}{\ensuremath{\subset}}
\newunicodechar{⊥}{\ensuremath{\perp}}
\newunicodechar{∃}{\ensuremath{\exists}}
\newunicodechar{∀}{\ensuremath{\forall}}
\newunicodechar{∛}{\ensuremath{\sqrt[3]{\,}}}
\newunicodechar{∜}{\ensuremath{\sqrt[4]{\,}}}
\newunicodechar{⃗}{\ensuremath{{}^{\to}}}
\newunicodechar{ⁿ}{\ifmmode^{n}\else\textsuperscript{\ensuremath{n}}\fi}
\newunicodechar{ˣ}{\ifmmode^{x}\else\textsuperscript{\ensuremath{x}}\fi}
\newunicodechar{ᵇ}{\ifmmode^{b}\else\textsuperscript{\ensuremath{b}}\fi}
\newunicodechar{ʳ}{\ifmmode^{r}\else\textsuperscript{\ensuremath{r}}\fi}
\newunicodechar{ᵘ}{\ifmmode^{u}\else\textsuperscript{\ensuremath{u}}\fi}
\newunicodechar{ʸ}{\ifmmode^{y}\else\textsuperscript{\ensuremath{y}}\fi}
\newunicodechar{ᵃ}{\ifmmode^{a}\else\textsuperscript{\ensuremath{a}}\fi}
\newunicodechar{ᵉ}{\ifmmode^{e}\else\textsuperscript{\ensuremath{e}}\fi}
\newunicodechar{ᵏ}{\ifmmode^{k}\else\textsuperscript{\ensuremath{k}}\fi}
\newunicodechar{ᵖ}{\ifmmode^{p}\else\textsuperscript{\ensuremath{p}}\fi}
\newunicodechar{ᴺ}{\ifmmode^{N}\else\textsuperscript{\ensuremath{N}}\fi}
\newunicodechar{ᶜ}{\ifmmode^{c}\else\textsuperscript{\ensuremath{c}}\fi}
\newunicodechar{ʰ}{\ifmmode^{h}\else\textsuperscript{\ensuremath{h}}\fi}
\newunicodechar{ᵠ}{\ifmmode^{q}\else\textsuperscript{\ensuremath{q}}\fi}
\newunicodechar{ₙ}{\ifmmode_{n}\else\textsubscript{\ensuremath{n}}\fi}
\newunicodechar{ₐ}{\ifmmode_{a}\else\textsubscript{\ensuremath{a}}\fi}
\newunicodechar{ᵢ}{\ifmmode_{i}\else\textsubscript{\ensuremath{i}}\fi}
\newunicodechar{ᵣ}{\ifmmode_{r}\else\textsubscript{\ensuremath{r}}\fi}
\newunicodechar{ₖ}{\ifmmode_{k}\else\textsubscript{\ensuremath{k}}\fi}
\newunicodechar{ᵦ}{\ifmmode_{\beta}\else\textsubscript{\ensuremath{\beta}}\fi}
\newunicodechar{ₓ}{\ifmmode_{x}\else\textsubscript{\ensuremath{x}}\fi}
\newfontfamily\popdisplay{ArchivoBlack-Regular.ttf}[Path=../../../fonts/]
\newfontfamily\poplogo{SpaceGrotesk-Bold.ttf}[Path=../../../fonts/]
\newfontfamily\popsemi{PlusJakartaSans-SemiBold.ttf}[Path=../../../fonts/]
\newfontfamily\popxbold{PlusJakartaSans-ExtraBold.ttf}[Path=../../../fonts/]
\newfontfamily\popxboldls{PlusJakartaSans-ExtraBold.ttf}[Path=../../../fonts/,LetterSpace=3.0]

\setlist[enumerate]{leftmargin=1.7em,itemsep=5pt,topsep=4pt}
\setlist[itemize]{leftmargin=1.7em,itemsep=4pt,topsep=4pt,label={\color{poporange}\textbullet}}
\setlength{\parindent}{0pt}
\setlength{\parskip}{5pt}
\renewcommand{\arraystretch}{1.25}

% pgfplots : style Pop par défaut
\pgfplotsset{
  every axis/.append style={axis line style={popink,line width=1pt},tick style={popink},
    grid style={popline,line width=0.5pt},label style={font=\small},tick label style={font=\footnotesize}},
  popcurve/.style={poporange,line width=1.8pt,no markers,smooth},
}
% Même style disponible dans un \draw TikZ simple (hors pgfplots)
\tikzset{popcurve/.style={poporange,line width=1.8pt}}
\tikzset{popgrid/.style={popline,very thin}}
\colorlet{popcurve}{poporange} % le nom du style est parfois utilisé comme couleur

% ============================================================
% Pill / squiggle
% ============================================================
\newcommand{\poppill}[3]{%
  \tikz[baseline=-0.55ex]\node[draw=popink,line width=1.2pt,rounded corners=2.6pt,%
    fill=#2,inner xsep=6.5pt,inner ysep=3pt]{%
    \popxboldls\fontsize{7}{7}\selectfont\color{#3}\MakeUppercase{#1}};%
}
\newcommand{\popsquiggle}[1]{%
  \tikz[baseline=-0.4ex]{
    \draw[#1,line width=2.6pt,line cap=round]
      plot [smooth] coordinates {(0,0.09) (0.34,0.01) (0.68,0.21) (1.02,0.01) (1.36,0.10)};
  }%
}
% Mot-clé surligné (marqueur orange sous le texte)
\newcommand{\cle}[1]{{\popxbold\color{popdark}#1}}

% ============================================================
% En-tête / pied
% ============================================================
\pagestyle{fancy}
\fancyhf{}
\renewcommand{\headrulewidth}{0pt}
\renewcommand{\footrulewidth}{0pt}
\lhead{\raisebox{-0.15ex}{\poplogo\fontsize{13}{13}\selectfont\color{popink}OnBuch\color{poporange}+}}
\rhead{\poppill{\DOCMATIERE\ \textbullet\ \DOCNIVEAU\ \textbullet\ Leçon \DOCLECON}{white}{popink}}
\lfoot{\popxbold\fontsize{7.6}{7.6}\selectfont\color{popmuted}\MakeUppercase{Cours signé OnBuch+}\ \textbullet\ {\popsemi\DOCTITRE}}
\rfoot{\tikz[baseline=-0.6ex]\node[rounded corners=8.5pt,fill=popink,minimum height=17pt,minimum width=17pt,inner xsep=5pt,inner sep=0]{\popxbold\fontsize{8}{8}\selectfont\color{popcream}\thepage\,/\,\pageref*{LastPage}};}

% ============================================================
% Titres
% ============================================================
\newcommand{\chaptitre}[2]{%
  \begin{tcolorbox}[enhanced,colback=poporange,colframe=popink,boxrule=2.6pt,arc=11pt,
      drop shadow=popink,left=15pt,right=15pt,top=12pt,bottom=11pt,before skip=3pt,after skip=18pt]
    \poppill{#2}{popink}{popcream}\par\vspace{9pt}
    {\raggedright\hyphenpenalty=10000\exhyphenpenalty=10000\popdisplay\fontsize{22}{24}\selectfont\color{popcream}\MakeUppercase{#1}\par}
  \end{tcolorbox}
}
% Section numérotée : pastille orange avec le numéro + titre Archivo
\newcounter{popsec}
\newcommand{\coursec}[1]{%
  \stepcounter{popsec}\setcounter{popsub}{0}%
  \par\vspace{16pt}\noindent
  \begin{minipage}{\linewidth}
  \tikz[baseline=-0.7ex]\node[draw=popink,line width=1.6pt,fill=poporange,rounded corners=4pt,minimum size=22pt,inner sep=0]{\popdisplay\fontsize{12}{12}\selectfont\color{popcream}\thepopsec};%
  \hspace{8pt}{\popdisplay\fontsize{15}{17}\selectfont\color{popink}\MakeUppercase{#1}}\par
  \vspace{4pt}\noindent{\color{poporange}\rule{36pt}{3.6pt}}
  \end{minipage}\par\nopagebreak\vspace{8pt}
}
\newcounter{popsub}
\newcommand{\courssub}[1]{%
  \stepcounter{popsub}%
  \par\vspace{9pt}\noindent{\popxbold\fontsize{11.5}{13}\selectfont\color{popink}{\color{poporange}\thepopsec.\thepopsub}\hspace{6pt}#1}\par\nopagebreak\vspace{4pt}
}

% ============================================================
% Boîtes pédagogiques — même recette : fond blanc, bordure épaisse colorée,
% ombre dure encre, badge plein. Titre optionnel [#1].
% ============================================================
\tcbset{popbase/.style={breakable,enhanced,colback=white,boxrule=2.3pt,arc=9pt,
  shadow={3pt}{-3pt}{0pt}{popink},left=14pt,right=14pt,top=11pt,bottom=11pt,before skip=11pt,after skip=15pt,
  fontupper=\popsemi\fontsize{10.6}{14.5}\selectfont\color{popink},
  fontlower=\popsemi\fontsize{10.6}{14.5}\selectfont\color{popink}}}
% Coche dessinée (le glyphe ✓ n'existe pas dans les polices du document)
\newcommand{\popcheck}[1]{\tikz[baseline=-0.5ex]\draw[#1,line width=1.6pt,line cap=round,line join=round](-0.13,0.0)--(-0.04,-0.09)--(0.13,0.1);}
\providecommand{\coche}{\popcheck{popgreen}}
\providecommand{\resultat}[1]{\par\smallskip\noindent\fcolorbox{popgreen}{popgreenL}{\parbox{\dimexpr\linewidth-2\fboxsep-2\fboxrule\relax}{\textbf{Résultat.} #1}}\par\smallskip}
\newcommand{\popboxhead}[3]{\poppill{#1}{#2}{white}\ifx\relax#3\relax\else\hspace{6pt}{\popxbold\fontsize{10.6}{12}\selectfont\color{popink}#3}\fi\par\vspace{7pt}}

\newtcolorbox{aretenir}[1][]{popbase,colframe=popgreen,before upper={\popboxhead{À retenir}{popgreen}{#1}}}
\newtcolorbox{exemplebox}[1][]{popbase,colframe=poporange,before upper={\popboxhead{Exemple}{poporange}{#1}}}
\newtcolorbox{definition}[1][]{popbase,colframe=popblue,before upper={\popboxhead{Définition}{popblue}{#1}}}
\newtcolorbox{propriete}[1][]{popbase,colframe=poppurple,before upper={\popboxhead{Propriété}{poppurple}{#1}}}
\newtcolorbox{methode}[1][]{popbase,colframe=popgold,before upper={\popboxhead{Méthode}{popgold}{#1}}}
\newtcolorbox{attention}[1][]{popbase,colframe=poppink,before upper={\popboxhead{Attention}{poppink}{#1}}}
\newtcolorbox{experience}[1][]{popbase,colframe=popblue,colback=popblueL,before upper={\popboxhead{Expérience}{popblue}{#1}}}
% « Le savais-tu ? » : carte violette pleine (contexte, culture, Cameroun)
\newtcolorbox{savaistu}[1][]{popbase,colback=poppurple,colframe=popink,
  fontupper=\popsemi\fontsize{10.4}{14.2}\selectfont\color{white},
  before upper={\poppill{Le savais-tu\,?}{white}{poppurple}\ifx\relax#1\relax\else\hspace{6pt}{\popxbold\color{white}#1}\fi\par\vspace{7pt}}}
% Exercice résolu : énoncé en haut, \tcblower puis la solution sur fond crème
\newcounter{popexo}\newcounter{popexr}
\newtcolorbox{exoresolu}[1][]{popbase,colframe=poporange,colbacklower=popcream,
  segmentation style={popink,line width=1.2pt,dashed},
  before upper={\stepcounter{popexr}\popboxhead{Exercice résolu \thepopexr}{poporange}{#1}},
  before lower={\poppill{Solution}{popink}{popcream}\par\vspace{6pt}}}
% Exercice (non résolu en place) — la correction est dans la partie Corrigés
\newtcolorbox{exercice}[1][]{popbase,colframe=popink,
  before upper={\stepcounter{popexo}\popboxhead{Exercice \thepopexo}{popink}{#1}}}
% Corrigé
\newtcolorbox{corrige}[1][]{popbase,colframe=popgreen,colback=popgreenL,
  before upper={\popboxhead{Corrigé}{popgreen}{#1}}}
% Objectifs / prérequis / bilan
\newtcolorbox{objectifs}{popbase,colframe=popink,colback=poporangeL,
  before upper={\popboxhead{Objectifs de la leçon}{popink}{}}}
\newtcolorbox{prerequis}{popbase,colframe=popink,colback=popblueL,
  before upper={\popboxhead{Prérequis}{popblue}{}}}
\newtcolorbox{bilan}{popbase,colframe=popink,colback=white,boxrule=2.8pt,
  before upper={\popboxhead{Fiche bilan}{poporange}{L'essentiel à connaître pour l'examen}}}
% Figure encadrée avec légende
\newtcolorbox{popfigure}[1][]{enhanced,breakable=false,colback=white,colframe=popline,boxrule=1.4pt,arc=8pt,
  left=10pt,right=10pt,top=10pt,bottom=8pt,before skip=10pt,after skip=12pt,halign=center,
  lower separated=false,halign lower=center,
  fontlower=\popsemi\fontsize{9}{11.5}\selectfont\color{popmuted},
  #1}
\newcommand{\legende}[1]{\par\vspace{6pt}{\popsemi\fontsize{9}{11.5}\selectfont\color{popmuted}#1}}

% ============================================================
% Signature OnBuch+
% ============================================================
\newcommand{\poplogoinline}[1]{{\poplogo\fontsize{#1}{#1}\selectfont\color{popink}OnBuch\color{poporange}+}}
\newcommand{\signatureonbuch}{%
  \begin{tcolorbox}[enhanced,colback=popink,colframe=popink,boxrule=2.2pt,arc=10pt,
      drop shadow=poporange,left=14pt,right=14pt,top=10pt,bottom=10pt,before skip=0pt,after skip=0pt]
    \tikz[baseline=-0.7ex]\node[circle,fill=popgreen,draw=popcream,line width=1.3pt,minimum size=22pt,inner sep=0]{\popcheck{white}};%
    \hspace{9pt}\begin{minipage}[c]{0.62\linewidth}
      {\popxbold\fontsize{12}{13}\selectfont\color{popcream}Signé\ \textbullet\ {\poplogo OnBuch\color{poporange}+}}\par\vspace{2pt}
      {\raggedright\popsemi\fontsize{8}{10.5}\selectfont\color{popline}\MakeUppercase{Équipe pédagogique OnBuch+}\\\MakeUppercase{Conforme au programme officiel MINESEC}\par}
    \end{minipage}\hfill
    {\poplogo\fontsize{22}{22}\selectfont\color{popcream}OnBuch\color{poporange}+}
  \end{tcolorbox}%
}

% ============================================================
% Page de garde
% #1 titre · #2 matière · #3 niveau · #4 module · #5 n° leçon · #6 durée ·
% #7 description / objectifs (texte ou itemize)
% ============================================================
\newcommand{\pagegarde}[7]{%
  \thispagestyle{empty}
  \newgeometry{a4paper,margin=1.7cm,top=1.6cm,bottom=1.4cm}%
  \begin{tikzpicture}[remember picture,overlay]
    \fill[poporange!18] ([xshift=-3.2cm,yshift=-3.6cm]current page.north east) circle (4.6cm);
    \fill[poppurple!14] ([xshift=2.4cm,yshift=7.6cm]current page.south west) circle (3.8cm);
  \end{tikzpicture}%
  {\poplogo\fontsize{30}{30}\selectfont\color{popink}OnBuch\color{poporange}+}\hfill
  \raisebox{0.6ex}{\poppill{Cours signé \textbullet\ Premium}{popink}{popcream}}\par
  \vspace{1pt}\popsquiggle{poppurple}\par
  \vspace{8pt}
  \begin{tcolorbox}[enhanced,colback=poporange,colframe=popink,boxrule=2.8pt,arc=13pt,
      drop shadow=popink,left=18pt,right=18pt,top=12pt,bottom=13pt,after skip=10pt]
    \poppill{#2 \textbullet\ #3}{popink}{popcream}\hspace{5pt}\poppill{#4}{white}{popink}\par\vspace{8pt}
    {\popdisplay\fontsize{14}{14}\selectfont\color{popink}LEÇON #5}\par\vspace{4pt}
    {\raggedright\hyphenpenalty=10000\exhyphenpenalty=10000\popdisplay\fontsize{25}{27}\selectfont\color{popcream}\MakeUppercase{#1}\par}
    \vspace{8pt}
    {\popxbold\fontsize{9.5}{9.5}\selectfont\color{popink}\MakeUppercase{Durée conseillée : #6}}
  \end{tcolorbox}
  \begin{tcolorbox}[enhanced,colback=white,colframe=popink,boxrule=2.2pt,arc=10pt,
      drop shadow=popink,left=16pt,right=16pt,top=10pt,bottom=10pt,after skip=8pt]
    \poppill{Dans cette leçon}{poppurple}{white}\par\vspace{5pt}
    \popsemi\fontsize{9.3}{12.6}\selectfont\color{popink}#7
  \end{tcolorbox}
  \noindent\begin{minipage}[t]{0.487\linewidth}
    \begin{tcolorbox}[enhanced,colback=popgreen,colframe=popink,boxrule=2.2pt,arc=10pt,
        drop shadow=popink,left=11pt,right=11pt,top=10pt,bottom=10pt,height=3.1cm,valign=center]
      \tcbox[colback=white,colframe=popink,boxrule=1.3pt,arc=5pt,boxsep=0pt,nobeforeafter,
        left=3pt,right=3pt,top=3pt,bottom=3pt,valign=center]{\qrcode[height=1.9cm]{https://play.google.com/store/apps/details?id=com.luvvix.onbuch&hl=fr}}%
      \hspace{8pt}\begin{minipage}[c]{0.5\linewidth}
        {\popdisplay\fontsize{11}{12.5}\selectfont\color{white}\MakeUppercase{Télécharge OnBuch}}\par\vspace{4pt}
        {\raggedright\popsemi\fontsize{8.4}{11}\selectfont\color{white}Gratuit \textbullet\ Google Play\\Tuteur IA, annales, concours, résultats\par}
      \end{minipage}
    \end{tcolorbox}
  \end{minipage}\hfill
  \begin{minipage}[t]{0.487\linewidth}
    \begin{tcolorbox}[enhanced,colback=poppurple,colframe=popink,boxrule=2.2pt,arc=10pt,
        drop shadow=popink,left=11pt,right=11pt,top=10pt,bottom=10pt,height=3.1cm,valign=center]
      \tikz[baseline=-0.7ex]\node[circle,fill=white,draw=popink,line width=1.3pt,minimum size=1.6cm,inner sep=0]{\popdisplay\fontsize{24}{24}\selectfont\color{poppurple}+};%
      \hspace{8pt}\begin{minipage}[c]{0.6\linewidth}
        {\popdisplay\fontsize{11}{12.5}\selectfont\color{white}\MakeUppercase{Passe à OnBuch+}}\par\vspace{4pt}
        {\raggedright\popsemi\fontsize{8.4}{11}\selectfont\color{white}Tous les cours signés, Tuteur IA illimité, fiches PDF — onglet OnBuch+ de l'app\par}
      \end{minipage}
    \end{tcolorbox}
  \end{minipage}
  \vspace{8pt}
  \popcontenu
  \vfill
  \signatureonbuch
  \restoregeometry
  \newpage
}

% Rangée « Ce cours contient » (page de garde)
\newcommand{\poptile}[3]{% #1 couleur #2 gros texte #3 légende
  \begin{tcolorbox}[enhanced,colback=#1,colframe=popink,boxrule=2pt,arc=9pt,shadow={2.5pt}{-2.5pt}{0pt}{popink},
      left=6pt,right=6pt,top=9pt,bottom=9pt,halign=center,valign=center,height=1.75cm,width=\linewidth,nobeforeafter]
    {\popdisplay\fontsize{12.5}{13}\selectfont\color{white}\MakeUppercase{#2}}\par\vspace{3pt}
    {\centering\popsemi\fontsize{7.8}{9.5}\selectfont\color{white}#3\par}
  \end{tcolorbox}}
\newcommand{\popcontenu}{%
  \par\noindent{\popxboldls\fontsize{7.5}{7.5}\selectfont\color{popmuted}CE COURS SIGNÉ CONTIENT}\par\vspace{7pt}
  \noindent\begin{minipage}[t]{0.235\linewidth}\poptile{popblue}{Cours}{complet, illustré, pas à pas}\end{minipage}\hfill
  \begin{minipage}[t]{0.235\linewidth}\poptile{poporange}{Méthodes}{et exercices résolus}\end{minipage}\hfill
  \begin{minipage}[t]{0.235\linewidth}\poptile{poppink}{Exercices}{type examen + corrigés}\end{minipage}\hfill
  \begin{minipage}[t]{0.235\linewidth}\poptile{popgreen}{Fiche bilan}{l'essentiel à retenir}\end{minipage}\par}

% Sommaire
\newtcolorbox{sommaire}{enhanced,colback=white,colframe=popink,boxrule=2.2pt,arc=10pt,
  drop shadow=popink,left=18pt,right=18pt,top=13pt,bottom=13pt,before skip=6pt,after skip=20pt,
  before upper={\poppill{Sommaire}{poppurple}{white}\par\vspace{9pt}\popsemi\fontsize{10.6}{14.5}\selectfont\color{popink}}}

% Signature de fin de cours — poussée tout en bas de la dernière page
\newcommand{\findecours}{%
  \par\vfill\nopagebreak
  \begin{center}\popsquiggle{poporange}\end{center}\vspace{4pt}
  \signatureonbuch
}
\AtBeginDocument{\def\llbracket{\mathopen{[\mkern-3.5mu[}}\def\rrbracket{\mathclose{]\mkern-3.5mu]}}} % intervalles d'entiers (glyphe ⟦ absent de la police)
% Glyphes absents de Fira Math : redéfinis à partir de glyphes disponibles
\AtBeginDocument{%
  \def\setminus{\mathbin{\backslash}}\let\smallsetminus\setminus
  \def\vdots{\mathord{\vbox{\baselineskip3.2pt\lineskiplimit0pt\kern4pt\hbox{.}\hbox{.}\hbox{.}}}}%
  \def\ddots{\mathinner{\mkern1mu\raise6.5pt\hbox{.}\mkern2mu\raise3.6pt\hbox{.}\mkern2mu\raise0.7pt\hbox{.}\mkern1mu}}%
  \def\triangle{\mathord{\scalebox{0.9}{$\symup{\Delta}$}}}%
  \def\nabla{\mathord{\raisebox{\depth}{\scalebox{1}[-1]{$\symup{\Delta}$}}}}%
  \def\checkmark{\popcheck{popdark}}%
  \def\star{\mathbin{\ast}}\def\bigstar{\mathord{\ast}}%
  \def\diamond{\mathbin{\scalebox{0.6}{$\Diamond$}}}%
  \def\bigcup{\mathop{\mathchoice{\raisebox{-0.3em}{\scalebox{1.7}{$\cup$}}}{\scalebox{1.25}{$\cup$}}{\cup}{\cup}}\displaylimits}%
  \def\bigcap{\mathop{\mathchoice{\raisebox{-0.3em}{\scalebox{1.7}{$\cap$}}}{\scalebox{1.25}{$\cap$}}{\cap}{\cap}}\displaylimits}%
  \def\lesseqqgtr{\mathrel{\lessgtr}}\def\bigoplus{\mathop{\oplus}\displaylimits}%
}
\providecommand{\ssi}{si et seulement si} % abréviation fréquente
\AtBeginDocument{\providecommand{\wideparen}{\overparen}} % arc de cercle

%% FIGFIX-BEGIN v3 — correctifs globaux des figures (docs/onbuch-generation/tools/figfix)
\usepackage{adjustbox}\usepackage{etoolbox}
% toute figure TikZ/pgfplots plus large que la ligne est réduite à la largeur disponible
\BeforeBeginEnvironment{tikzpicture}{\begin{adjustbox}{max width=\linewidth}}
\AfterEndEnvironment{tikzpicture}{\end{adjustbox}}
% légendes pgfplots lisibles par-dessus les courbes
\pgfplotsset{every axis/.append style={legend style={fill=white,fill opacity=0.92,text opacity=1,draw=black!15,inner sep=3pt}}}
% étiquettes d'arêtes (syntaxe edge["..."]) avec fond blanc
\tikzset{every edge quotes/.append style={fill=white,inner sep=1.2pt}}
% bulles de cartes mentales : texte à la ligne dans la bulle, bulle un peu plus grande
\tikzset{every concept/.append style={text width=2.0cm,align=center,minimum size=2.3cm,inner sep=2pt}}
%% FIGFIX-END
\begin{document}

\pagegarde{Application de la biotechnologie dans la production de biens et services}{SVT}{1ère A}{Module 4 : La Biotechnologie}{NA08}{3 h}{Cette leçon introduit la notion de biotechnologie, ses domaines d'application et ses réalisations traditionnelles (fermentation : bière, vin, yaourt, pain) et modernes (agriculture, médecine). Elle analyse enfin les enjeux économiques et sociaux des biotechnologies dans le contexte camerounais.
\vspace{4pt}
\begin{multicols}{2}\small
\begin{itemize}
\item À la fin de cette leçon, je sais définir la biotechnologie et citer ses grands domaines d'application.
\item Je sais décrire le principe général d'un procédé biotechnologique simple.
\item Je sais expliquer les étapes de la fermentation alcoolique et de la fermentation lactique.
\item Je sais citer des exemples traditionnels de biotechnologies (bière, vin, yaourt, pain, bobolo) et les microorganismes impliqués.
\item Je sais citer des exemples modernes de biotechnologies agricoles et médicales.
\item Je sais identifier des exemples locaux d'application de la biotechnologie au Cameroun.
\end{itemize}\end{multicols}}

\begin{sommaire}
\begin{multicols}{2}
\begin{enumerate}
\item Notion de biotechnologie
\item Principe d'un procédé biotechnologique : la fermentation
\item Biotechnologies traditionnelles : exemples
\item Biotechnologies modernes : agriculture et médecine
\item Avantages, limites et débats
\item Enjeux économiques et sociaux au Cameroun
\item Activité d'intégration
\item Méthodes et exercices résolus
\item Exercices
\item Corrigés des exercices
\item Fiche bilan
\end{enumerate}
\end{multicols}
\end{sommaire}

\begin{objectifs}
\begin{itemize}
\item À la fin de cette leçon, je sais définir la biotechnologie et citer ses grands domaines d'application.
\item Je sais décrire le principe général d'un procédé biotechnologique simple.
\item Je sais expliquer les étapes de la fermentation alcoolique et de la fermentation lactique.
\item Je sais citer des exemples traditionnels de biotechnologies (bière, vin, yaourt, pain, bobolo) et les microorganismes impliqués.
\item Je sais citer des exemples modernes de biotechnologies agricoles et médicales.
\item Je sais identifier des exemples locaux d'application de la biotechnologie au Cameroun.
\item Je sais discuter les avantages et les limites d'une biotechnologie donnée.
\item Je sais analyser les enjeux économiques et sociaux des biotechnologies pour le Cameroun.
\end{itemize}
\end{objectifs}

\begin{prerequis}
\begin{itemize}
\item Notion de microorganisme (bactéries, levures) vue en classe de Seconde
\item Notion de métabolisme cellulaire et de respiration/fermentation (début du programme de Première)
\item Notion d'enzyme et de catalyse biologique
\item Notion de gène et d'ADN (rappel utile pour les biotechnologies modernes)
\end{itemize}
\end{prerequis}

\newpage

\coursec{Notion de biotechnologie}

\courssub{Une situation de la vie quotidienne : la biotechnologie sans le savoir}

Chaque matin, au marché Mokolo à Yaoundé, Mama Ngo Bell vend du \emph{miondo} et du \emph{bobolo}, ces bâtons de pâte de manioc fermentée enveloppés dans des feuilles. Pour obtenir cette pâte souple et légèrement acidulée, elle a laissé reposer les racines de manioc dans l'eau pendant plusieurs jours : des \cle{microorganismes} invisibles à l'œil nu ont transformé le manioc, détruit ses composés toxiques (les composés cyanogènes) et développé son goût caractéristique. Sans le savoir, Mama Ngo Bell maîtrise une \cle{biotechnologie} vieille de plusieurs millénaires : la \cle{fermentation}.

Quelques kilomètres plus loin, à Douala, une grande brasserie produit chaque jour des milliers de bouteilles de bière grâce à des \cle{levures}, champignons microscopiques qui transforment les sucres en alcool. Et dans les hôpitaux du pays, on vaccine les nourrissons avec des vaccins dont certains sont fabriqués par \cle{génie génétique}, une biotechnologie moderne.

\begin{center}
\textbf{Problématique} : comment des êtres vivants, souvent microscopiques, peuvent-ils produire des biens et des services utiles à l'homme ? Qu'est-ce qui distingue les techniques ancestrales des biotechnologies modernes ?
\end{center}

\courssub{Définition de la biotechnologie}

\textbf{Observation et hypothèse.} Le miondo, le vin de palme, le yaourt, le pain, les vaccins, les variétés de maïs améliorées : ces produits très différents semblent n'avoir rien en commun. Pourtant, tous sont obtenus grâce à l'action d'\emph{êtres vivants} ou de \emph{molécules issues du vivant}, utilisés de manière maîtrisée par l'homme. On peut donc émettre l'hypothèse qu'il existe une famille de techniques unies par un principe commun : \emph{mettre le vivant au service de la production}.

\textbf{Vérification par les faits.} Dans chaque cas, on identifie un agent biologique et une transformation utile :

\begin{center}
\begin{tabularx}{\linewidth}{|l|X|X|}
\hline
\rowcolor{popblueL} \textbf{Produit} & \textbf{Agent biologique} & \textbf{Transformation utile} \\
\hline
Miondo, bobolo & Bactéries lactiques du manioc & Ramollissement, détoxification, acidification de la pâte \\
\hline
Bière, vin de palme & Levures (\emph{Saccharomyces}) & Transformation des sucres en éthanol et en \ce{CO2} \\
\hline
Yaourt & Bactéries lactiques & Transformation du lait en lait caillé acidifié \\
\hline
Pain levé & Levures de boulanger & Dégagement de \ce{CO2} qui fait lever la pâte \\
\hline
Vaccin contre l'hépatite B & Levures génétiquement modifiées & Production d'une protéine du virus servant d'antigène \\
\hline
\end{tabularx}
\end{center}

\textbf{Conclusion.} L'hypothèse est confirmée : dans tous ces cas, l'homme exploite l'activité d'organismes vivants ou de leurs constituants pour obtenir un bien ou un service. C'est exactement ce qu'on appelle une biotechnologie.

\begin{definition}[Biotechnologie]
La \cle{biotechnologie} est l'ensemble des techniques qui utilisent des \textbf{organismes vivants} (microorganismes, cellules végétales ou animales) ou leurs \textbf{constituants} (enzymes, gènes) pour fabriquer ou modifier des produits, améliorer des plantes ou des animaux, ou développer des microorganismes pour des usages spécifiques utiles à l'homme.
\end{definition}

\begin{aretenir}[Les trois ingrédients d'une biotechnologie]
Toute biotechnologie associe : (1) un \textbf{agent biologique} (organisme vivant, cellule, enzyme ou gène) ; (2) une \textbf{transformation} ou une \textbf{production} utile ; (3) une \textbf{maîtrise humaine} du procédé (choix des conditions : température, milieu, durée).
\end{aretenir}

\begin{methode}[Reconnaître un procédé biotechnologique]
Pour décider si un procédé est biotechnologique, vérifie trois critères :
\begin{enumerate}
\item Un \textbf{agent biologique} vivant ou dérivé du vivant est utilisé (levure, bactérie, moisissure, enzyme, cellule en culture, gène) ;
\item Une \textbf{transformation utile} est réalisée (production d'un aliment, d'un médicament, épuration d'une eau, amélioration d'une plante) ;
\item Le procédé est \textbf{maîtrisé par l'homme} (conditions contrôlées, résultat recherché et reproductible).
\end{enumerate}
Si l'un des trois critères manque, ce n'est pas une biotechnologie : la pourriture spontanée d'un fruit n'est pas une biotechnologie (critère 3 absent), et la cuisson du manioc à l'eau bouillante n'en est pas une non plus (critère 1 absent).
\end{methode}

\begin{savaistu}[D'où vient le mot « biotechnologie » ?]
Le mot \emph{biotechnologie} a été utilisé pour la première fois en \textbf{1919} par l'ingénieur hongrois \textbf{Karl Ereky}, pour désigner la production de biens à partir de matières vivantes, à l'aide d'organismes vivants. Le mot est donc récent, mais la pratique est immémoriale : les Sumériens et les Égyptiens fabriquaient déjà de la bière il y a plus de 6 000 ans !
\end{savaistu}

\courssub{Biotechnologies traditionnelles et biotechnologies modernes}

\textbf{Observation.} Mama Ngo Bell n'a jamais vu les bactéries qui fermentent son manioc ; elle applique un savoir-faire transmis par sa grand-mère. En revanche, l'ingénieur qui produit un vaccin sait exactement quel \cle{gène} il a introduit dans la levure. Ces deux démarches sont des biotechnologies, mais elles ne reposent pas sur le même niveau de connaissance.

\begin{definition}[Biotechnologies traditionnelles et modernes]
\begin{itemize}
\item Les \cle{biotechnologies traditionnelles} sont des techniques \textbf{empiriques et ancestrales}, mises au point par essais et erreurs, sans connaissance des mécanismes biologiques. Elles reposent essentiellement sur la \textbf{fermentation} (pain, bière, vin, yaourt, miondo, bobolo, ndôlô...) et sur la sélection empirique des plantes et des animaux domestiques.
\item Les \cle{biotechnologies modernes} sont fondées sur la \textbf{connaissance de l'ADN} et du fonctionnement des cellules. Elles comprennent le \textbf{génie génétique} (transfert de gènes, ADN recombinant), les \textbf{cultures cellulaires}, le clonage et la production d'enzymes industrielles.
\end{itemize}
\end{definition}

\begin{attention}[Erreur fréquente au Probatoire]
La biotechnologie \textbf{n'est pas une invention récente}. Le pain et la bière sont des produits biotechnologiques vieux de plusieurs millénaires. Seules les \emph{biotechnologies modernes} (génie génétique, ADN recombinant) datent du XX\textsuperscript{e} siècle. Écrire « la biotechnologie est une science nouvelle » est donc une erreur : c'est la \emph{compréhension scientifique} des procédés qui est récente, pas leur utilisation.
\end{attention}

\begin{popfigure}
\begin{tikzpicture}[scale=1]
\draw[thick, popdark] (0,0) -- (10,0);
\foreach \x in {0,1,...,10} \draw[popdark] (\x,-0.08) -- (\x,0.08);
% Étape 1
\fill[popgreen] (1,0) circle (0.12);
\node[above, align=center, font=\small, text=popgreen!60!black] at (1,0.15) {\textbf{vers -6000}\\Biotechnologies\\traditionnelles};
\node[below, align=center, font=\footnotesize] at (1,-0.15) {Fermentation : bière,\\pain, vin (Mésopotamie,\\Égypte, Afrique)};
% Étape 2
\fill[poporange] (5.5,0) circle (0.12);
\node[above, align=center, font=\small, text=poporange!70!black] at (5.5,0.15) {\textbf{1857}\\Découverte des microbes};
\node[below, align=center, font=\footnotesize] at (5.5,-0.15) {Louis Pasteur montre que\\la fermentation est due\\à des microorganismes vivants};
% Étape 3
\fill[poppurple] (9,0) circle (0.12);
\node[above, align=center, font=\small, text=poppurple!70!black] at (9,0.15) {\textbf{1973}\\Biotechnologies modernes};
\node[below, align=center, font=\footnotesize] at (9,-0.15) {Premier ADN recombinant\\(Cohen et Boyer) : naissance\\du génie génétique};
\end{tikzpicture}
\legende{Frise chronologique : de la fermentation empirique aux biotechnologies modernes. La découverte de Pasteur (1857) marque le passage d'un savoir-faire empirique à une compréhension scientifique ; l'ADN recombinant (1973) ouvre l'ère du génie génétique.}
\end{popfigure}

\courssub{Les grands domaines d'application de la biotechnologie}

Les biotechnologies touchent presque tous les secteurs de la vie quotidienne. On les classe classiquement en cinq grands domaines, auxquels correspond un \textbf{code couleur international} qu'il est utile de connaître pour le Probatoire.

\begin{center}
\begin{tabularx}{\linewidth}{|l|X|X|}
\hline
\rowcolor{popblueL} \textbf{Couleur} & \textbf{Domaine} & \textbf{Exemples} \\
\hline
Rouge & Santé, médecine & Vaccins, insuline humaine produite par des bactéries, antibiotiques, tests de diagnostic \\
\hline
Verte & Agriculture & Variétés améliorées, plantes transgéniques, culture de tissus, biopesticides \\
\hline
Blanche & Industrie & Enzymes des lessives, production d'éthanol-carburant, plastiques biodégradables \\
\hline
Bleue & Milieu marin & Exploitation des algues et des organismes marins (médicaments, cosmétiques) \\
\hline
Jaune & Environnement & Épuration des eaux usées, traitement des sols pollués (bioremédiation) \\
\hline
\end{tabularx}
\end{center}

\begin{popfigure}
\begin{tikzpicture}
% Centre
\node[circle, draw=popblue, fill=popblueL, thick, minimum size=2.4cm, align=center, font=\bfseries] (C) at (0,0) {Biotech-\\nologie};
% Branches
\node[rectangle, draw=popgreen, fill=popgreenL, rounded corners, align=center, font=\small] (agro) at (90:3.4) {\textbf{Agroalimentaire}\\bière, yaourt, miondo};
\node[rectangle, draw=poppink, fill=poppinkL, rounded corners, align=center, font=\small] (sante) at (18:3.6) {\textbf{Santé (rouge)}\\vaccins, insuline};
\node[rectangle, draw=popgreen!70!black, fill=popgreenL, rounded corners, align=center, font=\small] (agri) at (-54:3.6) {\textbf{Agriculture (verte)}\\variétés améliorées};
\node[rectangle, draw=popgold, fill=popgoldL, rounded corners, align=center, font=\small] (env) at (-126:3.6) {\textbf{Environnement (jaune)}\\épuration des eaux};
\node[rectangle, draw=popmuted, fill=popcream, rounded corners, align=center, font=\small] (ind) at (162:3.6) {\textbf{Industrie (blanche)}\\enzymes, bioéthanol};
\draw[thick, popline] (C) -- (agro);
\draw[thick, popline] (C) -- (sante);
\draw[thick, popline] (C) -- (agri);
\draw[thick, popline] (C) -- (env);
\draw[thick, popline] (C) -- (ind);
\end{tikzpicture}
\legende{Schéma en étoile des cinq grands domaines d'application de la biotechnologie, avec un exemple pour chacun. Les couleurs renvoient au code international (le domaine agroalimentaire, au centre des usages traditionnels, est rattaché historiquement aux biotechnologies blanche et rouge selon les produits).}
\end{popfigure}

\courssub{Les acteurs biologiques des biotechnologies}

Qui sont les « ouvriers invisibles » des biotechnologies ? On distingue plusieurs catégories d'agents biologiques.

\begin{itemize}
\item Les \cle{levures} : champignons unicellulaires (ex. \emph{Saccharomyces cerevisiae}). Elles réalisent la fermentation alcoolique (bière, vin de palme) et font lever le pain. Certaines levures modifiées produisent des protéines médicamenteuses.
\item Les \cle{bactéries} : les bactéries lactiques transforment le lait en yaourt et le manioc en pâte fermentée ; la bactérie \emph{Escherichia coli}, modifiée par génie génétique, produit l'insuline humaine.
\item Les \cle{moisissures} : champignons filamenteux. \emph{Penicillium} produit la pénicilline, premier antibiotique ; d'autres moisissures servent à fabriquer des fromages ou des enzymes industrielles.
\item Les \cle{enzymes} : protéines catalysatrices extraites d'organismes vivants. Elles entrent dans la composition des lessives (pour dégrader les taches) et de nombreux procédés alimentaires.
\item Les \cle{cellules en culture} : cellules végétales (culture de tissus pour multiplier rapidement des plants de bananiers ou de palmiers sains) ou cellules animales (production de vaccins et d'anticorps).
\end{itemize}

\begin{exemplebox}[La levure, championne toutes catégories]
\emph{Saccharomyces cerevisiae} est utilisée depuis l'Antiquité pour le pain et la bière (biotechnologie traditionnelle). La même espèce, génétiquement modifiée, sert aujourd'hui à fabriquer l'antigène du vaccin contre l'hépatite B (biotechnologie moderne). Un même organisme peut donc être l'acteur des deux types de biotechnologies : c'est la \emph{maîtrise du procédé} et le \emph{niveau de connaissance} qui changent, pas nécessairement l'organisme.
\end{exemplebox}

\begin{exoresolu}[Reconnaître une biotechnologie]
\textbf{Énoncé.} Pour chacune des situations suivantes, dire s'il s'agit d'une biotechnologie et justifier : (a) la fabrication du bobolo par fermentation du manioc ; (b) la décomposition d'une carcasse dans la savane ; (c) la production d'insuline par une bactérie portant le gène humain de l'insuline ; (d) la cuisson des arachides au feu de bois.
\tcblower
\textbf{Solution détaillée.} On applique les trois critères de la méthode : agent biologique, transformation utile, maîtrise humaine.

(a) \textbf{Oui.} Des bactéries lactiques (agent biologique) transforment le manioc en pâte détoxifiée et acidulée (transformation utile), selon un procédé maîtrisé et transmis (maîtrise humaine). C'est une biotechnologie \emph{traditionnelle}.

(b) \textbf{Non.} La décomposition met bien en jeu des microorganismes, mais elle n'est ni recherchée ni maîtrisée par l'homme : le critère 3 fait défaut.

(c) \textbf{Oui.} Une bactérie modifiée (agent biologique) produit une hormone utile au traitement du diabète (transformation utile) dans des fermenteurs contrôlés (maîtrise). C'est une biotechnologie \emph{moderne} (génie génétique, domaine rouge).

(d) \textbf{Non.} La cuisson est une transformation physico-chimique par la chaleur : aucun agent biologique n'intervient (critère 1 absent).
\end{exoresolu}

\begin{exercice}[À toi de jouer]
Le vin de palme, très consommé dans la région du Littoral, devient vinaigre s'il est conservé trop longtemps à l'air libre. (1) Identifie l'agent biologique responsable de la première étape (production d'alcool). (2) Le passage au vinaigre est-il, dans ce cas, une biotechnologie ? Justifie à l'aide des trois critères.
\end{exercice}

\begin{corrige}[Exercice : vin de palme]
(1) L'agent biologique est une \textbf{levure} (\emph{Saccharomyces}), présente naturellement dans la sève de palmier, qui réalise la fermentation alcoolique : \ce{C6H12O6 -> 2 C2H5OH + 2 CO2}.

(2) Le passage au vinaigre est dû à des bactéries acétiques qui oxydent l'éthanol en acide acétique. Si ce phénomène est subi (vin laissé à l'air libre), ce n'est \textbf{pas} une biotechnologie car la maîtrise humaine fait défaut. En revanche, si un producteur contrôle volontairement cette acétification pour fabriquer du vinaigre de palme, les trois critères sont réunis : c'est alors une biotechnologie traditionnelle.
\end{corrige}

\begin{aretenir}[L'essentiel de la section]
\begin{itemize}
\item La biotechnologie utilise des organismes vivants ou leurs constituants pour produire des biens et des services utiles.
\item Elle est \textbf{ancienne} (fermentation, sélection empirique) mais ses formes \textbf{modernes} reposent sur la connaissance de l'ADN (génie génétique, cultures cellulaires).
\item Ses domaines d'application suivent un code couleur : rouge (santé), verte (agriculture), blanche (industrie), bleue (milieu marin), jaune (environnement).
\item Ses acteurs sont les levures, bactéries, moisissures, enzymes et cellules en culture.
\end{itemize}
\end{aretenir}

\coursec{Principe d'un procédé biotechnologique : la fermentation}

Dans la section précédente, nous avons défini la \cle{biotechnologie} comme l'ensemble des techniques qui utilisent des organismes vivants, ou leurs constituants, pour produire des biens et des services utiles à l'homme. Mais une question demeure : \emph{comment} un microorganisme invisible à l'œil nu peut-il transformer du jus de maïs en bière, du lait en yaourt ou de la farine en pain ? La réponse tient en un mot : la \cle{fermentation}. C'est le procédé biotechnologique le plus ancien et le plus fondamental ; le comprendre, c'est comprendre le principe général de toutes les biotechnologies traditionnelles.

\courssub{Une observation de la vie quotidienne : la pâte qui lève}

\begin{experience}[La pâte à beignets qui « gonfle »]
\textbf{Observation.} À Douala comme à Bafoussam, la vendeuse de beignets mélange farine, sucre, eau tiède et un peu de levure de boulanger. Une heure plus tard, la pâte a \emph{doublé de volume} et présente de nombreuses bulles. Pourtant, personne n'y a ajouté d'air !

\textbf{Hypothèse.} La levure, organisme vivant, produit un gaz en consommant le sucre de la pâte.

\textbf{Vérification expérimentale.} On réalise deux témoins :
\begin{itemize}
\item \textbf{Ballon A} : eau tiède + sucre + levure, fermé par un doigt de gant ;
\item \textbf{Ballon B} : eau tiède + sucre seul (sans levure).
\end{itemize}
Au bout d'une heure, le doigt de gant du ballon A se gonfle d'un gaz qui trouble l'eau de chaux (c'est du dioxyde de carbone, \ce{CO2}) ; le ballon B reste inchangé.

\textbf{Interprétation.} La levure transforme le sucre en dégageant du \ce{CO2} (et, on le verra, de l'alcool). Ce gaz, piégé dans la pâte, la fait lever.

\textbf{Conclusion.} Un microorganisme placé dans un milieu sucré, à température convenable, réalise une transformation chimique utile : c'est une \cle{fermentation}.
\end{experience}

\courssub{Définition de la fermentation}

\begin{definition}[La fermentation]
La \cle{fermentation} est une \textbf{dégradation partielle de la matière organique} (généralement des sucres) par des \textbf{microorganismes} (levures ou bactéries), en \textbf{anaérobiose} (c'est-à-dire en l'absence de dioxygène). Elle fournit de l'\textbf{énergie au microorganisme} pour vivre et se multiplier, et produit des \textbf{substances utiles à l'homme} (alcool, gaz carbonique, acide lactique, arômes\ldots).
\end{definition}

Trois idées sont à retenir dans cette définition :
\begin{itemize}
\item La dégradation est \emph{partielle} : les produits obtenus (éthanol, acide lactique) contiennent encore beaucoup d'énergie chimique, contrairement à la respiration qui dégrade le glucose complètement en \ce{CO2} et \ce{H2O}.
\item Elle se déroule \emph{sans oxygène} : c'est une voie de secours qui permet au microorganisme de produire de l'énergie (sous forme d'ATP) quand le dioxygène manque.
\item Elle est \emph{utile à deux acteurs} : le microorganisme y gagne de l'énergie ; l'homme y gagne un produit (boisson, aliment, conservateur naturel).
\end{itemize}

\begin{attention}[Fermentation n'est pas respiration]
Erreur très fréquente au Probatoire : confondre fermentation et respiration cellulaire.
\begin{itemize}
\item La \textbf{respiration} se déroule \emph{en présence} de dioxygène (aérobiose) et dégrade \emph{totalement} le glucose : \ce{C6H12O6 + 6O2 -> 6CO2 + 6H2O}. Elle libère \emph{beaucoup} d'énergie.
\item La \textbf{fermentation} se déroule \emph{en l'absence} de dioxygène (anaérobiose), dégrade le glucose \emph{partiellement} et libère \emph{peu} d'énergie.
\end{itemize}
Dire « la levure respire le sucre pour faire de l'alcool » est donc une faute scientifique : la levure \emph{fermente} le sucre.
\end{attention}

\courssub{Les deux grandes fermentations à connaître}

\textbf{La fermentation alcoolique.}

Elle est réalisée par des \cle{levures}, champignons unicellulaires, dont l'espèce de référence est \emph{Saccharomyces cerevisiae} (levure de bière et de boulanger). En l'absence de dioxygène, la levure transforme le glucose en \textbf{éthanol} (alcool éthylique) et en \textbf{dioxyde de carbone}.

\begin{propriete}[Bilan de la fermentation alcoolique (admis)]
La fermentation alcoolique du glucose suit le bilan :
\[ \ce{C6H12O6 -> 2 C2H5OH + 2 CO2} + \text{énergie} \]
Ce bilan, établi grâce aux travaux de Louis Pasteur (qui démontra en 1857 que la fermentation est un phénomène \emph{vivant}, dû aux levures) puis précisé par la biochimie du XX\textsuperscript{e} siècle, est \textbf{admis sans démonstration} à ce niveau. Il faut savoir l'écrire et l'exploiter quantitativement.
\end{propriete}

\textbf{La fermentation lactique.}

Elle est réalisée par des \cle{bactéries lactiques}, notamment des genres \emph{Lactobacillus} et \emph{Streptococcus} (par exemple \emph{Streptococcus thermophilus} et \emph{Lactobacillus bulgaricus} dans le yaourt). Ces bactéries transforment le glucose (ou le lactose du lait, après hydrolyse) en \textbf{acide lactique} :
\[ \ce{C6H12O6 -> 2 C3H6O3} + \text{énergie} \]
L'acide lactique acidifie le milieu (le pH diminue), ce qui fait \emph{cailler} le lait et empêche le développement des microbes de détérioration : c'est un procédé naturel de \textbf{conservation} des aliments.

\begin{popfigure}
\begin{center}
\begin{tikzpicture}[font=\small, >=Stealth]
% Glucose au centre
\node[draw=popdark, fill=popgoldL, rounded corners, thick, minimum width=3cm, minimum height=1cm] (glu) at (0,0) {\textbf{Glucose} \ce{C6H12O6}};
% Branche gauche : fermentation alcoolique
\node[draw=popblue, fill=popblueL, rounded corners, thick, minimum width=3.6cm, align=center] (alc) at (-5.5,0) {\textbf{Éthanol} $2\,\ce{C2H5OH}$\\ \textbf{+ dioxyde de carbone} $2\,\ce{CO2}$};
\node[draw=popblue, fill=white, rounded corners, align=center, font=\footnotesize] (lev) at (-2.9,1.5) {Levures\\ \emph{Saccharomyces cerevisiae}};
\draw[->, very thick, popblue] (glu) -- (alc);
% Branche droite : fermentation lactique
\node[draw=popgreen, fill=popgreenL, rounded corners, thick, minimum width=3.6cm, align=center] (lac) at (5.5,0) {\textbf{Acide lactique}\\ $2\,\ce{C3H6O3}$};
\node[draw=popgreen, fill=white, rounded corners, align=center, font=\footnotesize] (bac) at (2.9,1.5) {Bactéries lactiques\\ \emph{Lactobacillus}, \emph{Streptococcus}};
\draw[->, very thick, popgreen] (glu) -- (lac);
% Conditions communes
\node[draw=poporange, fill=poporangeL, rounded corners, align=center, font=\footnotesize] (cond) at (0,-2.2) {Conditions communes : absence de dioxygène (anaérobiose),\\ température adaptée, milieu sucré, hygiène};
\draw[->, poporange, thick] (cond) -- (glu);
% Applications
\node[font=\footnotesize\itshape, popblue, align=center] at (-5.5,-1.6) {Applications : bière, vin,\\ pain (le \ce{CO2} fait lever la pâte)};
\node[font=\footnotesize\itshape, popgreen, align=center] at (5.5,-1.6) {Applications : yaourt, fromage,\\ lait caillé, choucroute};
\end{tikzpicture}
\end{center}
\legende{Schéma comparatif des deux grandes fermentations : à partir du même substrat (le glucose), les levures produisent éthanol et \ce{CO2} (fermentation alcoolique, à gauche) tandis que les bactéries lactiques produisent de l'acide lactique (fermentation lactique, à droite).}
\end{popfigure}

\begin{exemplebox}[Un calcul de bilan : combien d'alcool et de gaz ?]
Les masses molaires sont : $M(\ce{C6H12O6}) = 180$ g/mol ; $M(\ce{C2H5OH}) = 46$ g/mol ; $M(\ce{CO2}) = 44$ g/mol. D'après l'équation bilan, une mole de glucose donne deux moles d'éthanol et deux moles de \ce{CO2}. Donc, à partir de \textbf{180 g de glucose}, la fermentation alcoolique produit théoriquement :
\[ 2 \times 46 = 92 \text{ g d'éthanol} \qquad \text{et} \qquad 2 \times 44 = 88 \text{ g de } \ce{CO2}. \]
C'est précisément ce \ce{CO2} qui, piégé dans le réseau de gluten de la pâte à pain, forme les bulles qui la font \textbf{lever}. L'éthanol, lui, s'évapore presque totalement à la cuisson : le pain n'est pas un aliment alcoolisé !
\end{exemplebox}

\courssub{Les conditions d'une fermentation maîtrisée}

Pour qu'une fermentation réussisse et donne un produit sain et de qualité, cinq conditions doivent être réunies :

\begin{enumerate}
\item \textbf{Un substrat} : la matière organique à transformer, généralement un sucre (glucose du moût, lactose du lait, saccharose de la canne, amidon des céréales préalablement transformé en sucres).
\item \textbf{Un microorganisme} (le \cle{levain} ou la culture) : levures ou bactéries sélectionnées, introduites en quantité suffisante.
\item \textbf{Une température adaptée} : les enzymes des microorganismes ne fonctionnent bien que dans une gamme de température (environ 25--30 \textdegree C pour les levures de boulanger, 40--45 \textdegree C pour les ferments du yaourt). Trop froid : la fermentation s'arrête ; trop chaud : les microorganismes meurent.
\item \textbf{L'absence de dioxygène} (anaérobiose) : indispensable pour orienter le métabolisme vers la fermentation et non vers la respiration.
\item \textbf{L'hygiène} : récipients et ustensiles propres, pour éviter la contamination par des microorganismes indésirables qui donneraient de mauvais goûts ou rendraient le produit dangereux.
\end{enumerate}

\courssub{Les étapes générales d'un procédé biotechnologique simple}

Tout procédé de fermentation artisanal ou industriel suit la même logique en quatre étapes, qu'il s'agisse de la fabrication du « matango » (vin de palme), de la bière de maïs (« bil-bil » au Nord-Cameroun), du yaourt ou du pain.

\begin{popfigure}
\begin{center}
\begin{tikzpicture}[font=\small, >=Stealth, node distance=0.9cm,
etape/.style={draw=popdark, fill=popblueL, rounded corners, thick, minimum width=3.1cm, minimum height=1.1cm, align=center},
cond/.style={draw=poporange, fill=poporangeL, rounded corners, font=\footnotesize, align=center, minimum width=4.6cm}]
\node[etape, align=center] (e1){\textbf{1. Préparation}\\ \textbf{du substrat}};
\node[etape, below=of e1, align=center] (e2){\textbf{2. Inoculation}\\ (ensemencement)};
\node[etape, below=of e2, align=center] (e3){\textbf{3. Incubation}\\ (fermentation)};
\node[etape, below=of e3, align=center] (e4){\textbf{4. Récolte et}\\ \textbf{traitement du produit}};
\node[cond, right=1.6cm of e1, align=center] (c1){Broyer, cuire, sucrer ;\\ milieu propre et nutritif};
\node[cond, right=1.6cm of e2, align=center] (c2){Ajouter le levain : levures\\ ou bactéries sélectionnées};
\node[cond, right=1.6cm of e3, align=center] (c3){Température adaptée,\\ anaérobiose, durée contrôlée};
\node[cond, right=1.6cm of e4, align=center] (c4){Filtrer, cuire, conditionner,\\ conserver le produit};
\draw[->, very thick, popdark] (e1) -- (e2);
\draw[->, very thick, popdark] (e2) -- (e3);
\draw[->, very thick, popdark] (e3) -- (e4);
\draw[->, poporange] (c1) -- (e1);
\draw[->, poporange] (c2) -- (e2);
\draw[->, poporange] (c3) -- (e3);
\draw[->, poporange] (c4) -- (e4);
\end{tikzpicture}
\end{center}
\legende{Organigramme des quatre étapes d'un procédé de fermentation, avec les conditions associées à chaque étape. Ce schéma s'applique aussi bien à la fabrication artisanale du vin de palme qu'à la production industrielle de yaourt.}
\end{popfigure}

\begin{methode}[Décrire un procédé biotechnologique simple]
Face à une question du type « Décrire le principe de fabrication du yaourt / de la bière / du pain », suis toujours le même ordre, en cinq points :
\begin{enumerate}
\item \textbf{Substrat} : quelle matière organique est transformée ? (lait, moût de céréales, pâte de farine\ldots)
\item \textbf{Microorganisme} : quel agent vivant réalise la transformation ? (levure, bactérie lactique\ldots)
\item \textbf{Conditions} : température, anaérobiose, hygiène, durée.
\item \textbf{Transformation} : quelle réaction ? (écris l'équation bilan si elle est connue).
\item \textbf{Produit obtenu} : quelles substances utiles ? (éthanol, \ce{CO2}, acide lactique\ldots) et quel usage ?
\end{enumerate}
Une réponse qui suit cet ordre est complète, logique et valorisée par le correcteur.
\end{methode}

\courssub{Le rôle des produits de la fermentation dans les aliments}

\textbf{Le \ce{CO2} et la panification.} Dans la pâte à pain, la farine apporte des sucres, la levure de boulanger (\emph{Saccharomyces cerevisiae}) réalise la fermentation alcoolique. Le dioxyde de carbone dégagé reste emprisonné dans le réseau élastique du gluten : la pâte gonfle, devient alvéolée. À la cuisson, le gaz se dilate encore, l'alcool s'évapore, et le pain prend sa texture légère. Sans fermentation, on obtient une galette dure (pain azyme).

\textbf{L'éthanol et les boissons fermentées.} Dans le vin de palme, la bière de maïs ou de mil, le vin de raisin, c'est l'éthanol produit par les levures qui fait la spécificité de la boisson. Le taux d'alcool dépend de la quantité de sucre initiale et de la durée de fermentation. Le \ce{CO2}, lui, est responsable de l'effervescence de certaines boissons (bière, champagne).

\textbf{L'acide lactique et la conservation.} Dans le yaourt et le lait caillé, l'acide lactique abaisse le pH, coagule les protéines du lait (d'où la consistance ferme) et bloque le développement des bactéries de putréfaction : l'aliment se conserve plus longtemps et devient plus digeste.

\begin{exoresolu}[Décrire la fabrication du yaourt]
\textbf{Énoncé.} En utilisant la méthode en cinq points, décris le principe de la fabrication du yaourt, puis écris l'équation bilan de la transformation mise en jeu.
\tcblower
\textbf{Solution détaillée.}
\begin{enumerate}
\item \textbf{Substrat} : le lait, dont le sucre principal est le lactose (hydrolysé en glucose et galactose par les ferments).
\item \textbf{Microorganismes} : deux bactéries lactiques associées, \emph{Streptococcus thermophilus} et \emph{Lactobacillus bulgaricus}, ajoutées sous forme de ferments.
\item \textbf{Conditions} : lait préalablement pasteurisé (chauffé vers 85--90 \textdegree C puis refroidi) vers 42--45 \textdegree C, ensemencé, puis maintenu à cette température en pots fermés (milieu pauvre en oxygène) pendant quelques heures, dans des conditions d'hygiène strictes.
\item \textbf{Transformation} : fermentation lactique du glucose issu du lactose. Équation bilan :
\[ \ce{C6H12O6 -> 2 C3H6O3} + \text{énergie} \]
\item \textbf{Produit obtenu} : l'acide lactique acidifie le lait, ce qui coagule ses protéines : on obtient le yaourt, aliment ferme, acidulé et mieux conservé que le lait frais. On le refroidit ensuite pour stopper la fermentation.
\end{enumerate}
\end{exoresolu}

\begin{attention}[Les pièges à éviter]
\begin{itemize}
\item Ne jamais écrire que la fermentation « consomme de l'oxygène » : elle se déroule \emph{en anaérobiose}.
\item Ne pas confondre les produits : fermentation \emph{alcoolique} $\rightarrow$ éthanol + \ce{CO2} (levures) ; fermentation \emph{lactique} $\rightarrow$ acide lactique seul, sans dégagement de gaz (bactéries lactiques).
\item Dans l'équation de la fermentation alcoolique, ne pas oublier les coefficients : \ce{C6H12O6 -> 2 C2H5OH + 2 CO2}. Vérifie l'équilibre : 6 C, 12 H et 6 O de chaque côté.
\item Le pain ne contient pratiquement pas d'alcool : l'éthanol s'évapore à la cuisson. C'est le \ce{CO2}, et non l'alcool, qui fait lever la pâte.
\end{itemize}
\end{attention}

\begin{savaistu}[Pasteur, père de la microbiologie]
Au XIX\textsuperscript{e} siècle, on croyait que la fermentation était une simple réaction chimique « spontanée ». Louis Pasteur démontra entre 1857 et 1860, par des expériences soigneusement contrôlées, que ce sont des \emph{êtres vivants microscopiques} — les levures — qui en sont responsables, et que les « maladies du vin » venaient de contaminations par d'autres microbes. Ces travaux fondèrent la microbiologie et conduisirent à la \emph{pasteurisation}, procédé encore utilisé aujourd'hui pour conserver le lait vendu dans les supermarchés de Yaoundé et de Douala.
\end{savaistu}

\begin{aretenir}[L'essentiel de la section]
\begin{itemize}
\item La \textbf{fermentation} est une dégradation partielle de sucres par des microorganismes, \textbf{sans oxygène}, qui fournit de l'énergie au microorganisme et des produits utiles à l'homme.
\item \textbf{Fermentation alcoolique} (levures, \emph{Saccharomyces cerevisiae}) : \ce{C6H12O6 -> 2 C2H5OH + 2 CO2} + énergie.
\item \textbf{Fermentation lactique} (bactéries lactiques) : \ce{C6H12O6 -> 2 C3H6O3} + énergie.
\item Conditions : substrat sucré, levain, température adaptée, anaérobiose, hygiène.
\item Étapes d'un procédé : préparation du substrat $\rightarrow$ inoculation $\rightarrow$ incubation $\rightarrow$ récolte et traitement du produit.
\item Le \ce{CO2} fait lever la pâte à pain ; l'éthanol caractérise les boissons fermentées ; l'acide lactique conserve et épaissit le lait.
\end{itemize}
\end{aretenir}

Maintenant que le principe de la fermentation est clair, nous pouvons passer en revue, dans la section suivante, les grandes \textbf{biotechnologies traditionnelles} qui en découlent : bière, vin, yaourt et pain.

\coursec{Biotechnologies traditionnelles : exemples}

\begin{definition}[Biotechnologie traditionnelle]
On appelle \cle{biotechnologie traditionnelle} tout procédé de transformation de matières premières (végétales, animales) en produits utiles (aliments, boissons), mis en œuvre de manière empirique depuis des siècles ou des millénaires, qui repose sur l'activité métabolique de \cle{microorganismes} (levures, bactéries) sans que leurs auteurs n'aient nécessairement conscience de leur existence. Ces savoir-faire sont transmis oralement et s'adaptent aux ressources locales.
\end{definition}

Dans la section précédente, nous avons découvert le principe général d'un procédé biotechnologique à travers la fermentation : un \cle{microorganisme} (levure ou bactérie) transforme une matière première organique en un produit utile à l'Homme, dans des conditions contrôlées. Ce principe, loin d'être une invention récente des laboratoires, est appliqué par les sociétés humaines depuis des millénaires, bien avant que les microscopes ne révèlent l'existence des microbes. Au Cameroun comme ailleurs, nos grands-parents fabriquaient déjà de la bière de maïs, du vin de palme, du bobolo, sans connaître le mot « biotechnologie ». Cette section passe en revue les grandes \cle{biotechnologies traditionnelles}, en identifiant pour chacune le produit, le microorganisme responsable, le type de fermentation et l'intérêt pour les populations.

\courssub{La bière : fermentation alcoolique des sucres de céréales}

\textbf{Observation.} Dans de nombreuses localités du Nord et de l'Extrême-Nord camerounais, on prépare le « \cle{bil-bil} », une bière traditionnelle obtenue à partir de maïs, de sorgho ou de mil. Les grains sont d'abord trempés et germés, puis bouillis avec de l'eau ; après refroidissement, le liquide est laissé à l'air libre pendant un ou deux jours. Il se met alors à mousser et acquiert un goût légèrement alcoolisé. D'où viennent la mousse et l'alcool ?

\textbf{Interprétation.} La germination des grains produit des \cle{enzymes} (amylases) qui découpent l'amidon des céréales en sucres simples (maltose, glucose). Les \cle{levures} (champignons microscopiques du genre \emph{Saccharomyces}), présentes naturellement sur les grains ou ajoutées via un levain, réalisent ensuite la \cle{fermentation alcoolique} : en l'absence d'oxygène (milieu anaérobie), elles transforment les sucres en éthanol et en dioxyde de carbone selon l'équation :
\[ \ce{C6H12O6 -> 2 C2H5OH + 2 CO2} \]
Le \ce{CO2} dégagé forme la mousse ; l'éthanol donne à la boisson son caractère alcoolisé.

\begin{exemplebox}[La bière au Cameroun : du bil-bil à la production industrielle]
Le \textbf{bil-bil} (bière de maïs, de sorgho ou de mil) est fabriqué artisanalement, notamment dans les régions septentrionales, et joue un rôle social et culturel important (cérémonies, hospitalité, apport énergétique). À l'échelle industrielle, les \textbf{Brasseries du Cameroun (SABC)} appliquent exactement le même principe biologique — fermentation alcoolique de sucres de céréales par des levures — mais avec des souches de levures sélectionnées (\emph{Saccharomyces cerevisiae} souches de brasserie), des cuves fermées en inox, une température contrôlée, un brassage mécanisé et un contrôle qualité permanent (taux d'alcool, turbidité, contaminants). C'est la même biotechnologie, traditionnelle d'un côté, industrialisée de l'autre.
\end{exemplebox}

\courssub{Le vin de palme : fermentation naturelle de la sève}

\textbf{Observation.} Dans les zones forestières du Centre, du Littoral, du Sud et de l'Ouest, le palmier à vin (\emph{Raphia hookeri} ou palmier à huile \emph{Elaeis guineensis}) est incisé et sa sève sucrée est recueillie dans une calebasse ou un bidon. Quelques heures après la récolte, la sève douce et sucrée devient pétillante et alcoolisée : c'est le \cle{vin de palme}. Pourtant, personne n'a ajouté de levure !

\textbf{Hypothèse et interprétation.} La fermentation est dite \cle{spontanée} ou naturelle : les levures présentes dans l'air, sur l'écorce du palmier et sur les parois du récipient (contamination naturelle) colonisent la sève, riche en sucres (saccharose, glucose, fructose), et réalisent la fermentation alcoolique. Aucun ensemencement artificiel n'est nécessaire : le milieu naturel fournit les microorganismes.

\begin{savaistu}[Pourquoi le vin de palme devient-il acide après quelques jours ?]
Le vin de palme se fermente spontanément dès la récolte de la sève. Mais si on le conserve plusieurs jours, d'autres microorganismes — des bactéries acétiques du genre \emph{Acetobacter} — transforment l'éthanol en \textbf{acide acétique} en présence d'oxygène (milieu aérobie) : c'est la \cle{fermentation acétique}. L'équation simplifiée est :
\[ \ce{C2H5OH + O2 -> CH3COOH + H2O} \]
Le vin de palme « tourne » alors en vinaigre, d'où son goût de plus en plus acide. C'est pourquoi il doit être consommé frais, dans les heures qui suivent la récolte, ou pasteurisé pour stopper l'activité microbienne.
\end{savaistu}

\courssub{Le yaourt : fermentation lactique du lait}

\textbf{Observation.} Le lait frais se conserve mal : laissé à la chaleur ambiante, il « tourne » rapidement (coagulation désordonnée, mauvaise odeur). En revanche, le yaourt se conserve plusieurs jours et présente une consistance épaisse, un goût acidulé agréable. Comment des bactéries peuvent-elles à la fois « transformer » et « conserver » le lait ?

\textbf{Mécanisme.} Le yaourt résulte de la \cle{fermentation lactique} du lait par deux espèces de bactéries thermophiles qui travaillent en symbiose (protocoopération) : \emph{Lactobacillus delbrueckii} subsp. \emph{bulgaricus} et \emph{Streptococcus thermophilus}. Ces bactéries transforment le \cle{lactose} (sucre du lait, disaccharide) en \cle{acide lactique} (acide fort) après hydrolyse en glucose et galactose :
\[ \ce{C12H22O11 + H2O -> 4 CH3-CHOH-COOH} \]
L'acide lactique abaisse le pH du lait (pH passe de ~6,6 à ~4,2) : les protéines (caséines, sous forme de micelles) perdent leur stabilité colloïdale et coagulent (précipitation isoelectrique), ce qui donne au yaourt sa consistance épaisse (« le lait caille »). Cette acidité empêche le développement des microbes de décomposition (putréfaction) et de nombreux pathogènes, ce qui \textbf{conserve} le lait plus longtemps.

\begin{exemplebox}[Le yaourt : un exemple type de biotechnologie traditionnelle]
Le yaourt résulte de la fermentation lactique du lait : les bactéries (\emph{Lactobacillus bulgaricus} et \emph{Streptococcus thermophilus}) transforment le lactose en acide lactique, ce qui fait cailler le lait et le conserve plus longtemps. On retiendra la chaîne complète : \textbf{matière première} (lait) $\rightarrow$ \textbf{microorganismes} (bactéries lactiques thermophiles) $\rightarrow$ \textbf{transformation} (lactose $\rightarrow$ acide lactique + abaissement pH) $\rightarrow$ \textbf{produit} (yaourt, aliment conservé, plus digeste, enrichi en vitamines B).
\end{exemplebox}

\courssub{Le pain : la levure qui fait lever la pâte}

\textbf{Observation.} Une pâte de farine, d'eau et de sel reste compacte. Mais si on y ajoute de la \cle{levure de boulanger} (\emph{Saccharomyces cerevisiae}) et qu'on la laisse reposer au chaud (vers \SI{30}{\celsius}), elle gonfle et double de volume. Après cuisson (environ \SI{200}{\celsius}), le pain est léger et alvéolé. Quel gaz a produit ces alvéoles ?

\textbf{Mécanisme.} La levure réalise la \cle{fermentation alcoolique} des sucres de la farine (maltose issu de l'amidon hydrolysé par les amylases de la farine, et sucres ajoutés) : elle produit de l'éthanol et du dioxyde de carbone. C'est le \ce{CO2}, piégé dans le réseau viscoélastique de \cle{gluten} (protéines de la farine : gliadine + gluténine) de la pâte, qui forme des bulles et fait \textbf{lever} la pâte (poussée). Pendant la cuisson, l'alcool s'évapore (point d'ébullition \SI{78}{\celsius}) et le gaz dilaté par la chaleur laisse les alvéoles fixes qui donnent au pain sa légèreté et sa mie aérée.

\begin{attention}[Erreur fréquente : ce n'est pas l'alcool qui fait lever la pâte]
Dans le pain, c'est le \textbf{dioxyde de carbone (\ce{CO2})} — et non l'alcool — qui fait lever la pâte. L'éthanol produit par la fermentation \textbf{s'évapore presque totalement pendant la cuisson} : le pain cuit ne contient pratiquement plus d'alcool (traces résiduelles < 0,5 \%). Ne confonds donc jamais les deux produits de la fermentation alcoolique : dans la bière ou le vin, c'est l'\emph{éthanol} qui est le produit recherché ; dans le pain, c'est le \emph{\ce{CO2}} (agent levant).
\end{attention}

\courssub{Autres exemples locaux camerounais d'aliments fermentés}

Le Cameroun possède un riche patrimoine d'aliments fermentés, véritables biotechnologies traditionnelles transmises de génération en génération, adaptées à l'absence de chaîne du froid et aux ressources locales :

\begin{itemize}
\item Le \cle{bobolo} (bâtons) et le \cle{miondo} (galettes fines) : pâte de \textbf{manioc fermenté} (\emph{Manihot esculenta}). Les racines de manioc (variétés amères riches en composés cyanogènes) sont épluchées, trempées dans l'eau (retrempe) pendant 2 à 4 jours : des bactéries lactiques (\emph{Lactobacillus} spp.) et des levures fermentent la pulpe, l'attendrissent, développent son arôme caractéristique et \textbf{contribuent à éliminer les composés cyanogènes} (linamarine, lotaustraline) par lessivage et hydrolyse enzymatique (linamarase microbienne), réduisant la toxicité. La pâte est ensuite broyée, façonnée, enveloppée dans des feuilles de \emph{Megaphrynium macrostachyum} (ngongo) et cuite à la vapeur.
\item Le \cle{koki} : gâteau de haricots (niébé, \emph{Vigna unguiculata}) cuit à la vapeur dans des feuilles de bananier ; le trempage préalable des graines et la fermentation spontanée (bactéries lactiques) améliorent la digestibilité (réduction des oligosaccharides responsables des flatulences, des facteurs antinutritionnels comme les tanins et l'acide phytique) et la texture.
\item Le \cle{nkui} : sauce gluante traditionnelle de l'Ouest (Bamileke), préparée à partir d'écorces de \emph{Triumfetta pentandra} battues longuement dans l'eau. La préparation traditionnelle fait intervenir une \textbf{fermentation microbienne spontanée} (levures et bactéries lactiques) qui développe la viscosité (exopolysaccharides), la saveur et la conservation.
\item Le \cle{foléré} (oseille de Guinée, \emph{Hibiscus sabdariffa}) : les feuilles (et parfois les calices) peuvent être fermentées (silage lactique) pour être conservées et utilisées hors saison dans les sauces, assurant une disponibilité en légumes toute l'année.
\item Le \textbf{poisson fumé et fermenté} (ex. \emph{bonga} fumé, \emph{ndolé} de poisson) : la fermentation (autolyse enzymatique + flore microbienne de surface) associée au fumage (chaleur, composés phénoliques antibactériens, déshydratation) permet de conserver le poisson plusieurs mois dans les régions où la chaîne du froid est absente.
\end{itemize}

\begin{popfigure}
\begin{center}
\begin{tikzpicture}[
  etape/.style={rectangle, rounded corners=3pt, draw=popblue, fill=popblueL, thick, minimum width=2.8cm, minimum height=1.2cm, align=center, font=\small},
  fleche/.style={-{Stealth[length=2.5mm]}, thick, popdark},
  annot/.style={below=2pt, font=\scriptsize\itshape, text=popmuted, align=center}
]
\node[etape, align=center] (a) at (0,0){Racines de\\ \textbf{manioc}\\ (variétés amères)};
\node[etape, align=center] (b) at (3.6,0){\textbf{Trempage / Retrempe}\\ (fermentation 2--4 jours)};
\node[etape, align=center] (c) at (7.2,0){\textbf{Broyage}\\ de la pulpe fermentée};
\node[etape, align=center] (d) at (10.8,0){\textbf{Façonnage}\\ en bâtons (bobolo)\\ ou galettes (miondo)};
\node[etape, align=center] (e) at (14.4,0){\textbf{Cuisson}\\ à la vapeur (feuilles ngongo)};
\draw[fleche] (a) -- (b);
\draw[fleche] (b) -- (c);
\draw[fleche] (c) -- (d);
\draw[fleche] (d) -- (e);
\node[annot, xshift=0cm, align=center] at (b.south){Bactéries lactiques \& levures\\ Hydrolyse cyanogènes + Acidification};
\end{tikzpicture}
\end{center}
\legende{Étapes simplifiées de la fabrication du bobolo / miondo : l'étape clé est le trempage (retrempe), pendant lequel les microorganismes fermentent la pulpe de manioc, l'attendrissent, lui donnent son goût et réduisent sa toxicité (détoxication).}
\end{popfigure}

\courssub{Bilan comparatif et intérêts des fermentations traditionnelles}

Le tableau suivant résume les principaux exemples étudiés. Il constitue un modèle de réponse attendu au Probatoire pour la question « Citer un exemple de biotechnologie traditionnelle et en expliquer le principe ».

\begin{popfigure}
\begin{center}
\begin{tikzpicture}
\node[inner sep=0pt, align=center]{
\begin{tabular}{|p{2.8cm}|p{4.5cm}|p{4.2cm}|}
\hline
\rowcolor{popblueL} \textbf{Produit (Exemple local)} & \textbf{Microorganisme(s) responsable(s)} & \textbf{Type de fermentation} \\
\hline
Bière traditionnelle (bil-bil) / Bière industrielle (SABC) & Levures (\emph{Saccharomyces cerevisiae}) & Fermentation alcoolique \\
\hline
Vin de palme (sève de raphia / palmier à huile) & Levures naturelles de l'environnement (\emph{Saccharomyces}, \emph{Candida}) & Fermentation alcoolique (spontanée) \\
\hline
Yaourt (lait caillé artisanal / industriel) & \emph{Lactobacillus delbrueckii} subsp. \emph{bulgaricus}, \emph{Streptococcus thermophilus} & Fermentation lactique \\
\hline
Pain (levure de boulanger) & Levure de boulanger (\emph{Saccharomyces cerevisiae}) & Fermentation alcoolique \\
\hline
Bobolo / Miondo (manioc) & Bactéries lactiques (\emph{Lactobacillus} spp.) et levures & Fermentation mixte (majoritairement lactique) \\
\hline
Koki (niébé) & Bactéries lactiques (flore spontanée) & Fermentation lactique \\
\hline
\end{tabular}
};
\end{tikzpicture}
\end{center}
\legende{Tableau récapitulatif : pour chaque produit traditionnel, le(s) microorganisme(s) responsable(s) et le type de fermentation mis en jeu.}
\end{popfigure}

\begin{aretenir}[Les trois grands intérêts des fermentations traditionnelles]
\begin{enumerate}
\item \textbf{Conservation des aliments} : l'acidité (yaourt, bobolo, koki, foléré fermenté, pH < 4,5) ou l'alcool (bière, vin de palme, \SIrange{3}{10}{\% vol}) empêchent le développement des microbes de décomposition (putréfaction) et de nombreux microbes pathogènes. Les aliments fermentés se conservent plus longtemps sans réfrigérateur — un atout majeur dans les zones rurales et en période de soudure.
\item \textbf{Amélioration du goût, de la texture et de la digestibilité} : la fermentation produit des métabolites secondaires (arômes : esters, aldéhydes, acides ; texture : exopolysaccharides, coagulation protéique) et prédigère partiellement les macromolécules (amidon $\rightarrow$ sucres, protéines $\rightarrow$ peptides/acides aminés, oligosaccharides indigestes $\rightarrow$ fermentescibles), ce qui rend les aliments plus faciles à digérer et plus appétents.
\item \textbf{Enrichissement nutritionnel et sécurité sanitaire (détoxication)} : certaines fermentations augmentent la teneur en vitamines (notamment du groupe B : B2, B9, B12 par synthèse microbienne) et en biodisponibilité des minéraux (réduction de l'acide phytique chélateur). Dans le cas crucial du manioc amer, la fermentation \textbf{élimine les composés cyanogènes} (linamarine), rendant l'aliment consommable sans danger (taux d'acide cyanhydrique résiduel < 10 ppm selon normes CODEX).
\end{enumerate}
\end{aretenir}

\begin{methode}[Citer correctement un exemple local de biotechnologie (Probatoire)]
Pour citer un exemple local de biotechnologie traditionnelle (au Probatoire comme dans la vie courante), associe \textbf{toujours trois éléments} :
\begin{enumerate}
\item le \textbf{produit} final (ex. : le bobolo, le vin de palme, le koki) ;
\item le \textbf{microorganisme} ou le \textbf{type de fermentation} (ex. : bactéries lactiques et levures ; fermentation lactique du manioc ; fermentation alcoolique spontanée) ;
\item l'\textbf{intérêt pour les populations} (ex. : conservation du manioc, détoxication, aliment de base énergétique ; boisson sociale ; source de protéines végétales digestes).
\end{enumerate}
Une réponse qui ne cite que le produit (« le bobolo est une biotechnologie ») est incomplète : elle ne montre ni le mécanisme biologique (savoir) ni l'intérêt socio-économique (savoir-être / contexte).
\end{methode}

\begin{exoresolu}[Identifier la biotechnologie traditionnelle : le lait caillé artisanal]
\textbf{Énoncé.} À Bafoussam, une ménagère prépare du lait caillé (yaourt artisanal) en laissant du lait tiède en contact avec un peu de yaourt de la veille. (a) Nomme le type de fermentation mis en jeu. (b) Cite les microorganismes responsables. (c) Explique pourquoi le yaourt se conserve plus longtemps que le lait frais. (d) Pourquoi ajoute-t-elle un peu de yaourt de la veille ?
\tcblower
\textbf{Solution détaillée.}

(a) Il s'agit de la \textbf{fermentation lactique} : le lactose du lait est transformé en acide lactique par des bactéries lactiques.

(b) Les microorganismes responsables sont des bactéries lactiques thermophiles : \emph{Lactobacillus delbrueckii} subsp. \emph{bulgaricus} et \emph{Streptococcus thermophilus} (apportés par le ferment).

(c) L'acide lactique produit abaisse le pH du lait (vers 4,2). Cette acidité fait coaguler les protéines (caséines) : le lait caille (formation du gel). Ce pH acide crée un milieu défavorable au développement des microbes de décomposition (bactéries putréfiantes, levures de'altération) et de nombreux pathogènes : le yaourt se conserve donc plus longtemps que le lait frais (pH neutre, milieu favorable à toute flore).

(d) Le yaourt de la veille sert de \textbf{ferment} (ou levain, \emph{starter}) : il apporte une population dense et active de bactéries lactiques vivantes qui vont ensemencer le lait frais et démarrer rapidement la fermentation (inoculation), assurant la dominance de la flore d'intérêt sur la flore de contamination. C'est le même principe que l'ajout de levure dans la pâte à pain ou de « mère » de vinaigre.
\end{exoresolu}

\begin{exercice}[À toi de jouer : L'évolution du vin de palme]
Le vin de palme récolté le matin est doux et sucré ; le soir, il est pétillant et légèrement alcoolisé ; trois jours plus tard, il est devenu acide et vineux.
\begin{enumerate}
\item Explique l'évolution du goût entre le matin et le soir (nomme le microorganisme et le type de fermentation).
\item Explique l'apparition du goût acide au bout de trois jours (nomme le type de fermentation, le microorganisme et le produit formé).
\item Propose une raison technique pour laquelle le vin de palme doit être consommé frais ou pasteurisé rapidement.
\end{enumerate}
\end{exercice}

\begin{corrige}[Exercice : L'évolution du vin de palme]
\begin{enumerate}
\item \textbf{Matin $\rightarrow$ Soir :} Des levures naturellement présentes dans l'environnement (air, écorce du palmier, parois du récipient) réalisent la \textbf{fermentation alcoolique} des sucres de la sève (saccharose, glucose, fructose) : elles produisent de l'éthanol (goût alcoolisé) et du \ce{CO2} (pétillant, mousse). Milieu anaérobie.
\item \textbf{3 jours plus tard :} En présence d'oxygène (récipient ouvert ou perméable), des bactéries acétiques du genre \emph{Acetobacter} réalisent la \textbf{fermentation acétique} (oxydation de l'éthanol) : elles transforment l'éthanol en \textbf{acide acétique} (vinaigre), d'où le goût acide prononcé. Équation : $\ce{C2H5OH + O2 -> CH3COOH + H2O}$.
\item Le vin de palme doit être consommé frais (ou pasteurisé/conditionné sous atmosphère inerte) car sa fermentation se poursuit spontanément et de manière non contrôlée : l'alcool (produit recherché) se transforme progressivement en vinaigre (produit de dégradation), ce qui dégrade ses qualités organoleptiques et sa valeur marchande. La pasteurisation (\SI{70}{\celsius}, 15-20 min) détruit les levures et bactéries acétiques, stabilisant le produit.
\end{enumerate}
\end{corrige}

En résumé, les biotechnologies traditionnelles — bière, vin de palme, yaourt, pain, bobolo, koki, nkui, poisson fermenté et bien d'autres — reposent toutes sur le même principe fondamental : exploiter l'activité métabolique de microorganismes (levures, bactéries) pour transformer et conserver les ressources alimentaires locales. Longtemps empiriques, ces savoir-faire ancestraux constituent aujourd'hui un patrimoine culturel et scientifique que la biologie moderne permet de comprendre, d'optimiser (sélection de souches, contrôle des paramètres) et de valoriser économiquement, comme nous le verrons avec les biotechnologies modernes dans la section suivante.

\coursec{Biotechnologies modernes : agriculture et médecine}

Dans la section précédente, nous avons vu que l'humanité utilise depuis des millénaires des microorganismes pour transformer des aliments : la fermentation du pain, du vin, de la bière de maïs (\emph{bil-bil}) ou du \emph{nkui} relève de la biotechnologie \emph{traditionnelle}, pratiquée de manière empirique, sans comprendre les mécanismes en jeu. Depuis le milieu du XX\textsuperscript{e} siècle, la découverte de la structure de l'ADN (1953) et la maîtrise de la culture cellulaire ont tout changé : on peut désormais \cle{manipuler directement les gènes et les cellules}. C'est la naissance des \cle{biotechnologies modernes}. Cette section te montre comment elles transforment aujourd'hui l'agriculture, la médecine et même la protection de l'environnement.

\courssub{Le génie génétique : principe général}

\textbf{Situation-problème.} Un diabétique de type 1 ne produit plus d'insuline, l'hormone qui régule la glycémie. Pendant des décennies, on lui injectait de l'insuline extraite du pancréas de porcs ou de bovins abattus : coûteuse, parfois mal tolérée (réactions allergiques), et disponible en quantité limitée. Comment produire de l'insuline \emph{humaine} en grande quantité, sans prélever de pancréas humain ? La réponse, trouvée en 1982, illustre le principe du génie génétique.

\begin{definition}[Génie génétique]
Le \cle{génie génétique} est l'ensemble des techniques permettant d'\textbf{isoler}, de \textbf{modifier} et de \textbf{transférer} des \cle{gènes} d'un organisme à un autre, afin de conférer à l'organisme receveur un nouveau caractère. Un organisme ayant reçu un gène étranger est dit \cle{transgénique} ; c'est un \cle{organisme génétiquement modifié} (OGM).
\end{definition}

Le principe, que tu dois savoir décrire simplement, comporte quatre étapes principales :
\begin{enumerate}
\item \textbf{Identifier et isoler} le gène d'intérêt (par exemple le gène humain codant l'insuline) dans l'ADN de l'organisme donneur ;
\item \textbf{Insérer} ce gène dans un \cle{vecteur} (souvent un \cle{plasmide}, petit anneau d'ADN bactérien non chromosomique), puis introduire ce vecteur recombinant dans une cellule hôte (bactérie \emph{Escherichia coli}, levure, ou cellule mammalienne) : c'est la \cle{transformation} (ou transfection) ;
\item \textbf{Multiplier} la cellule modifiée en \cle{bioréacteur} (cuve de culture contrôlée : température, pH, nutriments, oxygène) : chaque cellule fille porte le gène étranger et l'exprime ;
\item \textbf{Récolter et purifier} la substance produite (ici, l'insuline humaine).
\end{enumerate}

\begin{propriete}[Universalité du code génétique — admise]
Une bactérie peut exprimer un gène humain et fabriquer la protéine humaine correspondante, car le \cle{code génétique est universel} : les mêmes codons (triplets de nucléotides) correspondent aux mêmes acides aminés chez (presque) tous les êtres vivants. Ce résultat fondamental est \textbf{admis sans démonstration} en Première ; il justifie la possibilité même du génie génétique.
\end{propriete}

\begin{exemplebox}[L'insuline humaine produite par des bactéries — 1\textsuperscript{er} médicament OGM]
Depuis \textbf{1982}, l'insuline injectée aux diabétiques est produite industriellement par la bactérie \emph{Escherichia coli} (ou la levure \emph{Saccharomyces cerevisiae}) dans laquelle on a inséré le \textbf{gène humain codant la proinsuline} (précurseur de l'insuline). Les bactéries se multiplient dans d'immenses bioréacteurs (jusqu'à 500\,000\,L), synthétisent la proinsuline, qui est ensuite extraite, purifiée et clivée enzymatiquement pour donner l'insuline mature. C'est la \textbf{première application majeure du génie génétique en médecine} : production illimitée, hormone strictement humaine (donc mieux tolérée, pas de risque de transmission de pathogènes animaux), coût réduit.
\end{exemplebox}

\begin{popfigure}
\begin{tikzpicture}[font=\small, >=Stealth]
% Étape 1
\draw[fill=popblueL, rounded corners=3pt] (0,0) rectangle (2.8,1.5);
\node[align=center] at (1.4,0.75) {Cellule humaine\\ \textbf{Gène de l'insuline}\\ (ADN chromosomique)};
% Étape 2
\draw[fill=poppurpleL, rounded corners=3pt] (3.8,0) rectangle (6.6,1.5);
\node[align=center] at (5.2,0.75) {Insertion du gène\\ dans un plasmide\\ (vecteur)};
\draw[fill=poppurple!40] (5.2,-1.2) ellipse (0.8 and 0.45);
\draw[thick, poppurple] (5.2,-1.2) ellipse (0.3 and 0.15);
\node[below, popdark] at (5.2,-1.8) {\emph{E. coli} transgénique};
\draw[->, thin, popmuted] (5.2,0.0) -- (5.2,-0.6);
% Étape 3
\draw[fill=popgreenL, rounded corners=3pt] (7.6,0) rectangle (10.4,1.5);
\node[align=center] at (9.0,0.75) {Bioréacteur :\\ multiplication\\ des bactéries};
% Étape 4
\draw[fill=poporangeL, rounded corners=3pt] (11.4,0) rectangle (14.2,1.5);
\node[align=center] at (12.8,0.75) {Extraction,\\ purification,\\ maturation\\ \textbf{Insuline humaine}};
% Flèches
\draw[->, very thick, popdark] (2.8,0.75) -- (3.8,0.75);
\draw[->, very thick, popdark] (6.6,0.75) -- (7.6,0.75);
\draw[->, very thick, popdark] (10.4,0.75) -- (11.4,0.75);
\end{tikzpicture}
\legende{Principe simplifié de la production d'insuline humaine par génie génétique : le gène humain est transféré dans une bactérie via un plasmide, qui le fabrique en grande quantité en bioréacteur.}
\end{popfigure}

\begin{attention}[Un OGM n'est pas un organisme « artificiel »]
Erreur fréquente à l'examen : écrire qu'un OGM est un être « créé de toutes pièces » en laboratoire. C'est faux. Un OGM est un \textbf{organisme vivant ordinaire} (maïs, bactérie, levure, saumon...) dont on a \textbf{modifié un ou quelques gènes} parmi des dizaines de milliers. L'essentiel de son patrimoine génétique reste inchangé, et il continue à vivre, se nourrir et se reproduire comme ses congénères non modifiés.
\end{attention}

\courssub{Biotechnologies agricoles}

\textbf{Situation-problème.} Au Cameroun, la \emph{Maladie de Panama} (fusariose, champignon \emph{Fusarium oxysporum} f.sp. \emph{cubense} race tropicale 4) menace les plantations de bananiers plantain, et les plants de palmier à huile de qualité sont longs et coûteux à obtenir par semis classique (hétérogénéité, juvénilité longue). Comment produire rapidement, et en grand nombre, des plants \textbf{sains} (exempts de pathogènes) et \textbf{performants} (haut rendement) ?

\courssub{La sélection assistée par marqueurs et la culture de tissus végétaux}

La \cle{sélection assistée par marqueurs moléculaires} (SAM) consiste à croiser des variétés choisies en s'aidant de marqueurs génétiques (ADN) pour repérer \emph{très tôt}, au stade plantule, les individus porteurs des caractères souhaités (rendement, résistance à une maladie, qualité nutritionnelle). Elle accélère considérablement la sélection traditionnelle pratiquée depuis toujours par les paysans (qui phénotypait au champ, sur plante adulte).

La \cle{culture de tissus végétaux} (ou \cle{micropropagation \emph{in vitro}}) exploite une propriété remarquable des cellules végétales : la \cle{totipotence}, c'est-à-dire la capacité d'une seule cellule végétale (ou d'un petit fragment de tissu) à régénérer une plante entière, génétiquement identique à la plante mère (clone). Le protocole type est le suivant :
\begin{enumerate}
\item \textbf{Prélèvement} d'un explant (méristème apical, segment nodal, bout de feuille) sur une plante-mère saine et performante (dite « élite ») ;
\item \textbf{Stérilisation} de surface de l'explant (eau de Javel, alcool) pour éviter toute contamination microbienne ;
\item \textbf{Mise en culture} sur un \textbf{milieu de culture stérile} (gélose, sels minéraux, vitamines, sucres, régulateurs de croissance : auxines, cytokinines) en conditions aseptiques (boîte de Pétri, salle de transfert) ;
\item \textbf{Multiplication} : bouturage \emph{in vitro} répété pour obtenir des centaines de plantules ;
\item \textbf{Enracinement} puis \textbf{acclimatation} progressive en pépinière (humidité élevée, ombre) avant mise en champ.
\end{enumerate}

\begin{exemplebox}[Applications concrètes au Cameroun]
\begin{itemize}
\item \textbf{Bananiers et plantains} : la micropropagation à partir de méristèmes (souvent indemnes de virus) fournit des plants \textbf{exempts de virus} (BBTV, BSV) et de champignons, ce qui sécurise la production contre la fusariose et les viroses. Le Centre Africain de Recherche sur Bananiers et Plantains (CARBAP) à Njombé est un acteur majeur.
\item \textbf{Palmier à huile} : le clonage par culture de tissus (embryogenèse somatique) permet de multiplier les palmiers \emph{Tenera} à très haut rendement (variétés Deli $\times$ La Mé), au lieu d'attendre les semis, longs (2 ans en pépinière) et aléatoires (ségrégation). La SOCAPALM et la CDC utilisent massivement cette technique.
\item \textbf{Cacaoyers} : la sélection assistée par marqueurs (SNP, SSR) aide à obtenir précocement des variétés tolérantes à la pourriture brune des cabosses (\emph{Phytophthora megakarya}), fléau majeur des cacaoculteurs du Centre et du Sud.
\item \textbf{Manioc} : production de plants sains (exempts de la mosaïque africaine du manioc - CMD) par micropropagation.
\end{itemize}
\end{exemplebox}

\courssub{Les organismes génétiquement modifiés (OGM) en agriculture}

Le génie génétique permet de créer des plantes cultivées dotées de caractères nouveaux, impossibles à obtenir par croisement classique (barrière d'espèce) :
\begin{itemize}
\item \textbf{Maïs Bt} : un gène \emph{cry} de la bactérie \emph{Bacillus thuringiensis} confère au maïs une \textbf{résistance aux insectes} foreurs de tiges (\emph{Busseola fusca}, \emph{Chilo partellus}) et à la pyrale, réduisant l'usage d'insecticides chimiques. Cultivé en Afrique du Sud, au Kenya, au Nigeria (essais), pas encore commercialisé au Cameroun (cadre réglementaire en cours).
\item \textbf{Cultures tolérantes à la sécheresse} : des gènes de résistance au stress hydrique (ex. \emph{DREB}, \emph{CspB}) permettent à certaines variétés de maïs (projet WEMA/TELA) de maintenir leur rendement en zone soudano-sahélienne (Extrême-Nord, Nord Cameroun).
\item \textbf{Riz « doré » (Golden Rice)} : enrichi en \textbf{provitamine A} (bêta-carotène) grâce à l'introduction de gènes du pathway des caroténoïdes (phytoène synthase du maïs, désaturase de \emph{Erwinia}) pour lutter contre la cécité nutritionnelle infantile. Approuvé aux Philippines (2021), pas encore au Cameroun.
\item \textbf{Coton Bt} : résistance aux lépidoptères ravageurs, cultivé au Burkina Faso (arrêté puis relancé), au Soudan, en Éthiopie.
\end{itemize}

\begin{attention}[Ne pas confondre les trois procédés — Piège classique du Probatoire]
\begin{itemize}
\item \textbf{Sélection (classique ou assistée)} : ne fait que croiser des organismes de la \textbf{même espèce} (ou espèces très proches compatibles sexuellement). Aucun gène étranger n'est introduit. On recombine l'existant.
\item \textbf{Transgenèse (OGM)} : introduit un gène provenant d'une \textbf{autre espèce} (bactérie $\to$ plante, homme $\to$ bactérie), franchissant la barrière d'espèce.
\item \textbf{Culture de tissus (micropropagation)} : ne modifie \textbf{aucun gène}. Elle ne fait que multiplier végétativement (clonage) des plants identiques au génotype de départ.
\end{itemize}
Au Probatoire, ces trois procédés doivent être distingués avec précision dans un tableau ou un texte argumenté.
\end{attention}

\courssub{Biotechnologies médicales}

\textbf{Situation-problème.} Dans un centre de santé de district au Cameroun, comment confirmer en quinze minutes un paludisme avant de prescrire (éviter l'automédication et la résistance) ? Comment disposer de vaccins sûrs, d'hormones humaines et d'antibiotiques efficaces ? Toutes ces avancées reposent sur les biotechnologies médicales.

\courssub{Vaccins, antibiotiques et protéines thérapeutiques}

\begin{itemize}
\item \textbf{Vaccins à sous-unités / recombinants} : produits par culture de cellules (levures, cellules d'insectes, cellules mammaliennes CHO) transgéniques. Exemple majeur : le vaccin contre l'\textbf{hépatite B} (Engerix-B, GenHevac B) est fabriqué par des levures \emph{Saccharomyces cerevisiae} produisant l'antigène de surface HBsAg. Le vaccin ne contient \textbf{aucun virus vivant}, donc \textbf{aucun risque d'infection}. De même pour le vaccin VPH (Gardasil, Cervarix : protéines L1 en VLPs).
\item \textbf{Antibiotiques} : la \cle{pénicilline} et ses dérivés sont produits industriellement par la culture en bioréacteurs (fermentation aérée, 100\,000--200\,000\,L) de la moisissure \emph{Penicillium chrysogenum}, souches mutées et sélectionnées pour des rendements élevés (passage de quelques UI/mL à plus de 50\,000 UI/mL).
\item \textbf{Hormones et protéines thérapeutiques} : insuline (bactérie/levure), \textbf{hormone de croissance humaine} (somatropine, bactérie/levure), \textbf{facteurs de coagulation} VIII et IX (cellules CHO) pour les hémophiles, \textbf{érythropoïétine} (EPO, cellules CHO) pour l'anémie de l'insuffisant rénal, \textbf{interférons} (bactérie/levure) pour hépatites/sclérose en plaques, \textbf{anticorps monoclonaux} (cellules CHO) pour cancers/arthrite (ex. rituximab, trastuzumab).
\end{itemize}

\begin{savaistu}[La découverte de la pénicilline — Sérendipité et rigueur]
En \textbf{1928}, le médecin écossais \textbf{Alexander Fleming} remarque qu'une moisissure, \emph{Penicillium notatum} (aujourd'hui \emph{P. chrysogenum}), a contaminé une boîte de Pétri de \emph{Staphylococcus aureus}... et que les bactéries ont disparu tout autour de la moisissure (zone d'inhibition). Il en déduit que celle-ci sécrète une substance antibactérienne : la \textbf{pénicilline}. Mais il ne parviendra pas à la purifier. Ce n'est qu'en \textbf{1939--1941} que l'équipe de \textbf{Florey et Chain} (Oxford) réussit l'extraction, la purification et les premiers essais cliniques humains spectaculaires. Prix Nobel 1945 partagé Fleming-Florey-Chain. Une observation attentive, une hypothèse, une vérification collective : la démarche scientifique dans toute sa splendeur !
\end{savaistu}

\courssub{Tests de diagnostic rapide (TDR)}

Les \cle{tests de diagnostic rapide} (TDR) utilisent des \textbf{anticorps monoclonaux} produits par biotechnologie (cellules hybrides \emph{hybridomes} ou recombinant en cellules CHO), fixés sur une bandelette de nitrocellulose (immunochromatographie). Déposée sur une goutte de sang total, la bandelette révèle en \textbf{5--15 minutes} la présence d'un antigène cible :
\begin{itemize}
\item \textbf{Paludisme} : détection de l'antigène HRP2 (\emph{P. falciparum}) ou pLDH (pan-paludique). Révolutionné la prise en charge au Cameroun (stratégie « Test-Traiter » OMS) : plus de prescription à l'aveugle, économie d'ACT, surveillance de la résistance.
\item \textbf{VIH} : tests de dépistage rapide (Determine, SD Bioline, etc.) sur sang total/sérum/plasma, utilisés en PEC (Prise En Charge) et PTME (Prévention Transmission Mère-Enfant).
\item \textbf{Covid-19} : tests antigéniques nasaux (détection protéine N), déployés massivement en 2020--2022.
\end{itemize}
\begin{attention}[TDR $\neq$ PCR]
Le TDR détecte des \textbf{protéines} (antigènes ou anticorps) : rapide, peu coûteux, faisable au chevet, mais moins sensible que la PCR. La \textbf{PCR} (Réaction de Polymérisation en Chaîne) détecte l'\textbf{ADN/ARN} du pathogène : très sensible, spécifique, mais nécessite laboratoire, thermocycleur, personnel formé, délai de plusieurs heures. Au Probatoire, ne pas confondre le principe (immunologique vs amplification nucléique).
\end{attention}

\courssub{Cultures de cellules animales et thérapie cellulaire}

On sait cultiver \emph{in vitro} des cellules animales et humaines (lignées cellulaires continues ou cellules primaires), ce qui sert à :
\begin{itemize}
\item \textbf{Toxicologie / Pharmacologie} : tester l'efficacité et la toxicité de nouveaux médicaments (réduction expérimentation animale, éthique 3R : Remplacer, Réduire, Raffiner) ;
\item \textbf{Production de vaccins} : virus atténués/inactivés cultivés sur cellules Vero (rein de singe), cellules MDCK (rein de chien), ou cellules d'insectes (baculovirus) ;
\item \textbf{Ingénierie tissulaire} : cultures de kératinocytes (peau) pour greffe sur grands brûlés (ex. Epicel), cultures de chondrocytes (cartilage) pour réparation articulaire (MACI).
\end{itemize}
La \cle{thérapie cellulaire} consiste à administrer à un patient des cellules saines pour réparer un tissu défaillant. L'exemple historique et majeur est la greffe de \textbf{cellules souches hématopoïétiques} (CSH) de la moelle osseuse (ou sang périphérique mobilisé, ou sang de cordon), pratiquée pour soigner certaines \textbf{leucémies}, \textbf{lymphomes}, \textbf{aplasies médullaires} ou \textbf{drépanocytose} (greffe allogénique HLA-identique). À ton niveau, retiens simplement le principe : \emph{remplacer ou réparer des cellules malades par des cellules saines cultivées ou sélectionnées en laboratoire}. La thérapie génique (correction du gène défectueux \emph{dans} les cellules du patient, ex. \emph{Strimvelis} pour ADA-SCID, \emph{Zolgensma} pour amyotrophie spinale, \emph{Casgevy} pour drépanocytose/bêta-thalassémie par édition CRISPR/Cas9) en est l'évolution ultime, encore très coûteuse et peu accessible au Cameroun.

\courssub{Biotechnologies environnementales (complément)}

Les microorganismes peuvent aussi travailler pour l'environnement :
\begin{itemize}
\item \textbf{Traitement des eaux usées} : dans les stations d'épuration (ex. station de Bassa à Douala, ou futures stations Yaoundé), des bactéries aérobies dégradent la matière organique polluante (DBO, DCO) en \textbf{boues activées} (aération forcée), rendant l'eau moins dangereuse avant rejet en milieu naturel. Traitement anaérobie (méthanisation) pour boues et effluents très chargés (agro-industrie : brasseries, huileries, sucreries).
\item \textbf{Production de biogaz} : la \cle{digestion anaérobie} (méthanisation) des déchets organiques (fumier, résidus de marché, eaux-vannes, résidus agricoles) par un consortium bactérien (hydrolyse $\to$ acido-genèse $\to$ acéto-genèse $\to$ méthano-genèse) produit du \textbf{biogaz} (50--70\% \ce{CH4}, 30--50\% \ce{CO2}), utilisable pour la cuisson, l'éclairage ou la cogénération électricité/chaud. Le \textbf{digestat} résiduel est un excellent engrais organique. Des biodigesteurs domestiques (type chinois à dôme fixe, ou indien à tambour flottant) et institutionnels (lycées, prisons, fermes) se développent dans les zones rurales et périurbaines camerounaises : moins de bois de chauffe (déforestation), moins de déchets, assainissement, engrais pour les champs. Projet PNDP, GIZ, SNV, ONG locales.
\item \textbf{Bioremédiation} : utilisation de bactéries/champignons pour dégrader des polluants (hydrocarbures, métaux lourds, pesticides) sur sites contaminés (ex. dépôts pétroliers, sites miniers).
\end{itemize}

\courssub{Bilan : biotechnologies traditionnelles et modernes}

\begin{popfigure}
\begin{tikzpicture}[font=\small, grow=down, level distance=1.7cm,
  every node/.style={draw, rounded corners=3pt, align=center, fill=popcream, inner sep=3pt},
  edge from parent/.style={draw=popmuted, thick},
  level 1/.style={sibling distance=5.5cm},
  level 2/.style={sibling distance=3.2cm}]
\node[fill=popblueL, font=\small\bfseries, text width=3cm] {Biotechnologies}
  child { node[fill=popgreenL, text width=3.5cm] {\textbf{Traditionnelles}\\ (empiriques : fermentation,\\ pas de modification génétique)}
    child { node[align=center]{Aliments : pain, vin,\\ bière (bil-bil, nkui),\\ yaourt, fromage, ngoundja} }
    child { node[align=center]{Environnement :\\ biogaz domestique,\\ lagunage, compostage} }
  }
  child { node[fill=poporangeL, text width=3.5cm] {\textbf{Modernes}\\ (scientifiques : manipulation\\ gènes, cellules, enzymes)}
    child { node[align=center]{Génie génétique :\\ insuline, vaccins,\\ OGM (Bt, doré),\\ thérapie génique} }
    child { node[align=center]{Culture de tissus :\\ bananiers, palmier,\\ manioc, cacaoyer\\ (plants sains, clones)} }
    child { node[align=center]{Cultures cellulaires :\\ anticorps monoclonaux,\\ tests rapides (TDR),\\ thérapie cellulaire,\\ peau artificielle} }
    child { node[align=center]{Enzymologie :\\ lessives, cuir,\\ agroalimentaire,\\ bioéthanol 2G} }
  };
\end{tikzpicture}
\legende{Classification des biotechnologies. Les biotechnologies traditionnelles exploitent le métabolisme naturel de microorganismes sans les modifier ; les biotechnologies modernes manipulent gènes, cellules et enzymes de façon contrôlée et rationnelle.}
\end{popfigure}

\begin{center}
\begin{tabularx}{\linewidth}{|l|X|X|}\hline
\rowcolor{popblueL} \textbf{Critère} & \textbf{Biotechnologie traditionnelle} & \textbf{Biotechnologie moderne} \\ \hline
Connaissances de base & Empiriques, transmises oralement, sans comprendre le mécanisme moléculaire & Scientifiques rigoureuses (ADN, gènes, protéines, régulation métabolique, génomique) \\ \hline
Outils / Infrastructures & Jarres, fours, calebasses, fosses, cuves simples, conditions non stériles & Bioréacteurs pilotés, salles blanches (PSM), cultures stériles, enzymes de restriction, PCR, séquenceurs, CRISPR/Cas9 \\ \hline
Organismes utilisés & Microorganismes sauvages (levures, bactéries lactiques, moisissures) non modifiés & Microorganismes modifiés (OGM), cellules animales/végétales cultivées, enzymes purifiées, anticorps monoclonaux \\ \hline
Exemples emblématiques & Fermentation alimentaire : pain, vin, bière, yaourt, fromage, \emph{ngoundja}, \emph{kwem} ; biogaz artisanal, lagunage & Insuline recombinante, vaccins à sous-unités (HBsAg, VPH), OGM (maïs Bt, riz doré), micropropagation (bananier, palmier), TDR paludisme/VIH, anticorps monoclonaux, thérapie cellulaire/génique \\ \hline
Traçabilité / Contrôle & Qualité variable, difficile à standardiser, risque contamination pathogène & Qualité standardisée (BPF/Bonnes Pratiques de Fabrication), traçabilité lot par lot, sécurité accrue \\ \hline
\end{tabularx}
\end{center}

\begin{exoresolu}[Identifier un procédé biotechnologique et son classement]
\textbf{Énoncé.} Une coopérative de l'Extrême-Nord obtient, à partir d'un fragment de tige (méristème) d'un palmier à huile \emph{Tenera} à haut rendement, cinq mille plantules identiques cultivées sur milieu gélosé stérile enrichi en régulateurs de croissance. (a) De quel procédé biotechnologique s'agit-il ? (b) S'agit-il d'une biotechnologie traditionnelle ou moderne ? Justifier par deux arguments. (c) Les plantules obtenues sont-elles des OGM ? Justifier.
\tcblower
\textbf{Solution détaillée.}
\begin{enumerate}
\item[(a)] Il s'agit de la \textbf{culture de tissus végétaux} (ou \textbf{micropropagation \emph{in vitro}} par embryogenèse somatique ou organogenèse), fondée sur la \textbf{totipotence} des cellules végétales.
\item[(b)] C'est une biotechnologie \textbf{moderne}. Arguments : (1) Elle repose sur des connaissances scientifiques précises (composition du milieu de culture MS, rôle des auxines/cytokinines, stérilité stricte) et non sur un savoir empirique. (2) Elle nécessite des infrastructures de laboratoire (salle de transfert laminaire, autoclave, chambre de culture climatisée, éclairage contrôlé) inaccessibles à l'artisanat traditionnel.
\item[(c)] \textbf{Non}, les plantules ne sont \textbf{pas des OGM}. Justification : aucun gène étranger n'a été introduit dans le génome du palmier. Les plantules sont des \textbf{clones} (copies génétiques identiques) de la plante mère \emph{Tenera}. Un OGM suppose une modification du patrimoine génétique par transgenèse (introduction d'un gène d'une autre espèce) ou édition ciblée (CRISPR), ce qui n'est pas le cas ici. La culture de tissus est une technique de \textbf{multiplication végétative}, pas de modification génétique.
\end{enumerate}
\end{exoresolu}

\begin{aretenir}[L'essentiel de la section — Fiche de révision Probatoire]
\begin{itemize}
\item Le \textbf{génie génétique} (transgenèse) permet de transférer un gène d'un organisme à un autre ; c'est possible car le \textbf{code génétique est universel} (admis sans démonstration en 1ère). Étapes : isolement gène $\to$ insertion dans vecteur (plasmide) $\to$ transformation hôte $\to$ multiplication en bioréacteur $\to$ purification produit.
\item \textbf{Agriculture} :
      \begin{itemize}
      \item \textbf{Sélection assistée (marqueurs moléculaires)} : accélère la sélection classique, pas d'OGM.
      \item \textbf{Culture de tissus (micropropagation)} : exploite la \textbf{totipotence} cellulaire $\to$ plants sains (exempts virus), clones identiques, déploiement rapide (bananier CARBAP, palmier SOCAPALM/CDC, manioc, cacaoyer). \textbf{Pas d'OGM}.
      \item \textbf{OGM agricoles} : Maïs Bt (résistance insectes \emph{Bt}), Coton Bt, Maïs tolérant sécheresse (TELA), Riz doré (provitamine A). Distinguer sélection / transgenèse / culture de tissus.
      \end{itemize}
\item \textbf{Médecine} :
      \begin{itemize}
      \item \textbf{Protéines recombinantes} : Insuline (1982, 1\textsuperscript{er} médicament OGM), hormone de croissance, facteurs coagulation, EPO, interférons, anticorps monoclonaux (cellules CHO).
      \item \textbf{Vaccins recombinants} : Hépatite B (levure HBsAg), VPH (VLPs). Sécurité : pas de virus vivant.
      \item \textbf{Antibiotiques} : Pénicilline (\emph{Penicillium} en bioréacteur, souches améliorées par mutagenèse/sélection).
      \item \textbf{Diagnostic} : TDR (immunochromatographie, anticorps monoclonaux) : paludisme (HRP2/pLDH), VIH, Covid-19. Différencier TDR (protéines, rapide, terrain) et PCR (ADN/ARN, sensible, labo).
      \item \textbf{Thérapie cellulaire} : Greffe CSH (moelle osseuse) pour leucémies/drépanocytose. Thérapie génique (correction gène \emph{in vivo} ou \emph{ex vivo}, ex. CRISPR/Cas9 \emph{Casgevy}) : pointe avancée, coût très élevé.
      \end{itemize}
\item \textbf{Environnement} : Épuration eaux usées (boues activées aérobies, anaérobie), \textbf{Biogaz} (méthanisation déchets $\to$ \ce{CH4} + digestat engrais), Bioremédiation.
\item \textbf{Distinction cruciale} : OGM = gène étranger introduit (transgenèse). Clone (culture de tissus) = copie conforme, pas de gène étranger. Sélection = recombinaison naturelle/assistée intra-espèce.
\end{itemize}
\end{aretenir}

\coursec{Avantages, limites et débats}

Dans la section précédente, nous avons découvert les biotechnologies modernes : plantes génétiquement modifiées, production d'insuline humaine par des bactéries, vaccins, thérapies géniques. Ces applications suscitent un immense espoir... mais aussi de vives controverses. Au marché de Mokolo, une ménagère se réjouit que le ndolé soit produit toute l'année grâce aux semences améliorées ; au même moment, une émission de radio débat avec passion de la dangerosité supposée des OGM. Qui a raison ? Comme toujours en science, la réponse ne se trouve ni dans la peur ni dans l'enthousiasme aveugle, mais dans une \cle{analyse critique} : peser méthodiquement les \textbf{bénéfices} et les \textbf{risques} de chaque biotechnologie. C'est l'objet de cette section, qui te donnera aussi une méthode infaillible pour traiter ce type de question au Probatoire.

\courssub{Les avantages des biotechnologies}

\textbf{Production à grande échelle et à moindre coût}\par

Commençons par une observation qui a changé la vie de millions de malades. Avant 1982, l'insuline destinée aux diabétiques était extraite du pancréas de porcs et de bovins abattus : il fallait environ \SI{8000}{kg} de pancréas pour produire \SI{1}{kg} d'insuline, d'où un coût élevé, des pénuries régulières et des réactions allergiques chez certains patients (l'insuline animale diffère légèrement de l'insuline humaine par quelques acides aminés). Depuis 1982, le gène humain de l'insuline est inséré dans la bactérie \emph{Escherichia coli}, qui le fabrique en cuves de culture (fermenteurs) : la production devient quasi illimitée, le produit est identique à l'hormone humaine (insuline recombinante), et le coût a fortement diminué.

\begin{aretenir}[Avantages majeurs des biotechnologies]
\begin{itemize}
\item \textbf{Production à grande échelle et à moindre coût} : les micro-organismes se multiplient très vite dans des fermenteurs ; on obtient ainsi d'importantes quantités de substances utiles (insuline, vaccins, enzymes, antibiotiques) à un coût inférieur à celui de l'extraction animale ou végétale.
\item \textbf{Amélioration des rendements agricoles} : variétés améliorées (plus productives, résistantes aux maladies, à la sécheresse ou aux insectes), cultures in vitro pour multiplier rapidement des plants sains (bananiers, plantains, manioc au Cameroun).
\item \textbf{Conservation des aliments} : la fermentation (lactique notamment) acidifie le milieu et empêche le développement des micro-organismes de putréfaction : le lait caillé, le yaourt, le « nkui » fermenté ou le manioc transformé en bâton se conservent bien plus longtemps que les produits frais.
\item \textbf{Valorisation des déchets} : les micro-organismes transforment les déchets organiques en produits utiles : compost, biogaz (méthane produit par fermentation anaérobie des déchets, utilisable pour la cuisson), traitement biologique des eaux usées.
\end{itemize}
\end{aretenir}

\begin{savaistu}[Le biogaz, une biotechnologie d'avenir au Cameroun]
Dans certaines fermes et lycées techniques camerounais, des \textbf{biodigesteurs} transforment les déjections animales et les déchets de cuisine en \cle{biogaz} (méthane) grâce à des bactéries anaérobies. Le gaz alimente les cuisinières, et le résidu digéré (digestat) devient un excellent engrais organique. Une biotechnologie simple, locale, qui réduit la déforestation liée au bois de chauffe.
\end{savaistu}

\courssub{Les limites techniques}

Toute médaille a son revers. Reprenons notre démarche : observons les difficultés rencontrées concrètement par les utilisateurs de biotechnologies.

\begin{itemize}
\item \textbf{Risque de contamination des cultures} : une culture microbienne (levures pour la bière, bactéries productrices d'insuline) peut être envahie par des micro-organismes indésirables qui détruisent la production ou la rendent impropre à la consommation. Il faut donc travailler dans des conditions d'asepsie strictes, ce qui exige matériel et formation.
\item \textbf{Coût élevé des équipements modernes} : fermenteurs, incubateurs, appareils d'électrophorèse, enceintes stériles (PSM)... Les biotechnologies modernes demandent des investissements lourds, souvent hors de portée des petits producteurs et des laboratoires des pays en développement.
\item \textbf{Dépendance vis-à-vis des semences brevetées} : les semences OGM sont protégées par des brevets détenus par de grandes firmes. L'agriculteur doit les racheter chaque année (interdiction contractuelle de réutiliser la récolte comme semence, et pour les hybrides F1, perte de l'hétérosis en F2), ce qui crée une dépendance économique et menace les semences paysannes traditionnelles.
\end{itemize}

\begin{attention}[Ne pas confondre limite technique et danger sanitaire]
Dire « les équipements coûtent cher » ou « les semences sont brevetées » relève d'une limite \textbf{technique et économique}, pas d'un danger pour la santé. Au Probatoire, classe correctement chaque argument : technique, sanitaire, environnemental, économique ou éthique.
\end{attention}

\courssub{Les limites sanitaires et environnementales : le débat sur les OGM}

\textbf{Situation-problème.} En 2012, une étude française (équipe Séralini) affirme que des rats nourris au maïs OGM NK603 développent plus de tumeurs ; l'étude est ensuite retirée par la revue \emph{Food and Chemical Toxicology} pour insuffisance de données (taille d'échantillon trop faible, lignée de rats sujette aux tumeurs). Pendant ce temps, des centaines d'études indépendantes et des méta-analyses concluent que les OGM autorisés sont aussi sûrs que les aliments conventionnels. Comment s'y retrouver ? Par la \cle{démarche scientifique} : une affirmation n'a de valeur que si elle repose sur des expériences reproductibles, évaluées par la communauté scientifique (pair review).

\begin{itemize}
\item \textbf{Risque d'allergies} : l'introduction d'un gène étranger peut faire produire à la plante une nouvelle protéine, potentiellement allergène. C'est pourquoi tout nouvel OGM subit des tests d'allergénicité avant sa mise sur le marché (un OGM de soja contenant un gène de noix du Brésil a ainsi été abandonné dès les tests, avant toute commercialisation).
\item \textbf{Dissémination de gènes dans la nature} : le pollen d'une plante OGM peut féconder des plantes sauvages voisines ou des variétés conventionnelles et leur transmettre le gène introduit (ex. : gène de résistance à un herbicide), créant des « adventices » difficiles à éliminer (flux de gènes).
\item \textbf{Sélection d'organismes résistants} : l'usage massif d'une même toxine (ex. : toxine Cry du \emph{Bacillus thuringiensis} dans le maïs Bt ou le coton Bt) exerce une pression de sélection qui favorise l'apparition d'insectes résistants, rendant la technologie inefficace à terme.
\item \textbf{Nécessité d'évaluations rigoureuses} : chaque OGM doit être évalué au cas par cas, par des études scientifiques indépendantes, avant et après sa mise en culture (suivi environnemental post-commercialisation).
\end{itemize}

\begin{exemplebox}[Le maïs Bt : un cas d'école]
Le \textbf{maïs Bt} est un maïs génétiquement modifié qui produit une toxine insecticide (protéine Cry) issue de la bactérie \emph{Bacillus thuringiensis}.
\begin{itemize}
\item \textbf{Avantages} : la plante se protège elle-même contre la pyrale (insecte foreur) ; l'usage d'insecticides chimiques diminue fortement ; les rendements sont plus stables ; réduction de l'exposition des agriculteurs aux pesticides.
\item \textbf{Limites} : la pression de sélection peut faire apparaître des insectes résistants à la toxine Bt (nécessité de zones refuges non-Bt) ; les agriculteurs deviennent dépendants des semenciers qui détiennent les brevets ; la dissémination du gène vers des variétés locales ou sauvages est possible par le pollen.
\end{itemize}
Conclusion nuancée : le maïs Bt est utile s'il est cultivé avec des règles de biosécurité (zones refuges, rotation des cultures, surveillance des résistances), mais il pose de vraies questions économiques et environnementales.
\end{exemplebox}

\courssub{Les questions éthiques}

Au-delà de la science, les biotechnologies posent des questions de \textbf{société} :

\begin{itemize}
\item \textbf{Brevetabilité du vivant} : peut-on posséder un gène, une semence, un organisme vivant comme on possède une machine ? Les brevets récompensent l'investissement en recherche, mais beaucoup jugent choquant qu'une entreprise « possède » un vivant ou une séquence génétique naturelle.
\item \textbf{Accès équitable aux innovations} : les médicaments issus des biotechnologies (insuline, vaccins, anticorps monoclonaux) restent chers ; comment garantir que les populations pauvres, notamment en Afrique, en bénéficient (licences obligatoires, transferts de technologie) ?
\item \textbf{Respect des savoirs traditionnels} : les communautés locales détiennent des connaissances précieuses sur les plantes médicinales et les semences paysannes. Il arrive que des entreprises exploitent ces savoirs sans contrepartie ni consentement préalable : c'est la \textbf{biopiraterie}. Le droit international (Protocole de Nagoya sur l'accès aux ressources génétiques et le partage juste et équitable des avantages, 2010) cherche à protéger ces ressources et savoirs.
\end{itemize}

\courssub{La biosécurité}

Face à ces risques, la société ne renonce pas aux biotechnologies : elle les encadre.

\begin{definition}[Biosécurité]
La \cle{biosécurité} désigne l'ensemble des mesures visant à prévenir les risques que les biotechnologies (notamment l'usage d'organismes vivants modifiés -- OVM) font peser sur la santé humaine et sur l'environnement.
\end{definition}

Concrètement, la biosécurité comprend : le confinement des organismes modifiés en laboratoire (enceintes étanches, personnels formés, niveaux de confinement P1 à P4), l'évaluation des risques avant toute dissémination volontaire ou mise sur le marché, l'étiquetage des produits OGM (seuil de 0,9 \% en UE), et des accords internationaux comme le \textbf{Protocole de Carthagène} (2000, entré en vigueur 2003), que le Cameroun a ratifié, qui encadre les mouvements transfrontaliers d'organismes génétiquement modifiés (procédure d'accord préalable en connaissance de cause).

\courssub{Synthèse : tableau comparatif et balance bénéfices--risques}

\begin{popfigure}
\begin{center}
\begin{tikzpicture}[every node/.style={font=\small}]
% En-têtes
\fill[popblue] (0,0) rectangle (4.2,0.8);
\fill[popgreen] (4.2,0) rectangle (10.1,0.8);
\fill[poporange] (10.1,0) rectangle (16,0.8);
\node[text=white,font=\small\bfseries] at (2.1,0.4) {Biotechnologie};
\node[text=white,font=\small\bfseries] at (7.15,0.4) {Avantages};
\node[text=white,font=\small\bfseries] at (13.05,0.4) {Limites};
% Ligne 1
\fill[popblueL] (0,-2.2) rectangle (4.2,0);
\fill[popgreenL] (4.2,-2.2) rectangle (10.1,0);
\fill[poporangeL] (10.1,-2.2) rectangle (16,0);
\node[text width=3.8cm,align=center] at (2.1,-1.1) {\textbf{Fermentation alimentaire} (bière, yaourt, bâton de manioc)};
\node[text width=5.5cm,align=left] at (7.15,-1.1) {Conservation des aliments ; procédé simple et peu coûteux ; valeur ajoutée locale ; digestibilité améliorée};
\node[text width=5.5cm,align=left] at (13.05,-1.1) {Risque de contamination ; hygiène indispensable ; qualité variable selon les artisans ; chaîne du froid nécessaire};
% Ligne 2
\fill[popblueL!60] (0,-4.6) rectangle (4.2,-2.2);
\fill[popgreenL!60] (4.2,-4.6) rectangle (10.1,-2.2);
\fill[poporangeL!60] (10.1,-4.6) rectangle (16,-2.2);
\node[text width=3.8cm,align=center] at (2.1,-3.4) {\textbf{OGM agricoles} (maïs Bt, coton Bt)};
\node[text width=5.5cm,align=left] at (7.15,-3.4) {Rendements améliorés ; moins de pesticides ; résistance aux insectes/sécheresse ; revenus paysans potentiels};
\node[text width=5.5cm,align=left] at (13.05,-3.4) {Insectes résistants possibles ; dissémination de gènes ; dépendance aux semenciers brevetés ; coût semences};
% Ligne 3
\fill[popblueL] (0,-7) rectangle (4.2,-4.6);
\fill[popgreenL] (4.2,-7) rectangle (10.1,-4.6);
\fill[poporangeL] (10.1,-7) rectangle (16,-4.6);
\node[text width=3.8cm,align=center] at (2.1,-5.8) {\textbf{Production d'insuline} par bactéries modifiées};
\node[text width=5.5cm,align=left] at (7.15,-5.8) {Production massive et moins chère ; insuline identique à l'humaine ; moins d'allergies ; pureté constante};
\node[text width=5.5cm,align=left] at (13.05,-5.8) {Équipements coûteux ; asepsie stricte ; dépendance technologique vis-à-vis des firmes ; chaîne du froid};
% Contours
\draw[popdark,thick] (0,0.8) rectangle (16,-7);
\draw[popdark] (4.2,0.8) -- (4.2,-7);
\draw[popdark] (10.1,0.8) -- (10.1,-7);
\draw[popdark] (0,-2.2) -- (16,-2.2);
\draw[popdark] (0,-4.6) -- (16,-4.6);
\end{tikzpicture}
\end{center}
\legende{Tableau comparatif des avantages et limites de trois biotechnologies : chaque procédé présente des bénéfices réels et des contraintes qu'il faut connaître pour porter un jugement équilibré.}
\end{popfigure}

\begin{popfigure}
\begin{center}
\begin{tikzpicture}[scale=1]
% Pied de la balance
\fill[popdark] (-0.15,0) rectangle (0.15,3.2);
\fill[popdark] (-1.2,0) rectangle (1.2,0.25);
% Fléau
\draw[popdark,very thick] (-4.2,3.2) -- (4.2,3.2);
\fill[popgold] (0,3.2) circle (0.12);
% Fils et plateau gauche (bénéfices)
\draw[popdark] (-3.6,3.2) -- (-4.4,2.1);
\draw[popdark] (-3.6,3.2) -- (-2.8,2.1);
\draw[popgreen,very thick,fill=popgreenL] (-4.8,2.1) arc (180:360:1.2);
\node[text=popgreen!60!black,font=\bfseries] at (-3.6,1.4) {BÉNÉFICES};
\node[text width=4.6cm,align=center,font=\footnotesize] at (-3.6,0.6) {santé (insuline, vaccins)\\ rendements agricoles\\ conservation des aliments\\ valorisation des déchets};
% Fils et plateau droit (risques)
\draw[popdark] (3.6,3.2) -- (4.4,2.1);
\draw[popdark] (3.6,3.2) -- (2.8,2.1);
\draw[poporange,very thick,fill=poporangeL] (2.4,2.1) arc (180:360:1.2);
\node[text=poporange!70!black,font=\bfseries] at (3.6,1.4) {RISQUES};
\node[text width=4.6cm,align=center,font=\footnotesize] at (3.6,0.6) {contaminations\\ allergies possibles\\ dissémination de gènes\\ dépendance économique};
\end{tikzpicture}
\end{center}
\legende{La démarche d'analyse critique : comme une balance, on pèse les bénéfices et les risques d'une biotechnologie avant de formuler un avis argumenté. Les mesures de biosécurité jouent le rôle de régulateur.}
\end{popfigure}

\courssub{Méthode d'analyse d'une biotechnologie donnée}

\begin{methode}[Discuter les avantages et limites d'une biotechnologie (méthode Probatoire)]
Pour traiter la consigne « Discute des avantages et des limites de cette biotechnologie », structure ta réponse en trois temps :
\begin{enumerate}
\item \textbf{Décrire le procédé} : nomme la biotechnologie et explique brièvement son principe (organisme utilisé, produit obtenu, réaction clé si fermentation).
\item \textbf{Citer au moins deux avantages} : choisis-les dans des catégories différentes si possible (sanitaire, économique, agricole, environnemental).
\item \textbf{Citer au moins une limite} et \textbf{conclure par un avis argumenté} : ta conclusion doit mettre en balance bénéfices et risques, et proposer éventuellement une condition d'utilisation (biosécurité, encadrement, évaluation, zone refuge).
\end{enumerate}
\end{methode}

\begin{attention}[Erreur fréquente : « OGM = dangereux »]
Affirmer « les OGM sont dangereux » (ou, à l'inverse, « les OGM ne posent aucun problème ») est une \textbf{opinion}, pas un fait scientifique. Au Probatoire, on attend un raisonnement équilibré : des avantages \textbf{ET} des limites documentés, appuyés sur des exemples précis (maïs Bt, insuline, coton Bt...), sans parti pris. Une copie manichéenne perd des points, même si elle est bien écrite.
\end{attention}

\begin{exoresolu}[Discuter une biotechnologie : le yaourt produit industriellement]
\textbf{Énoncé.} Une laiterie de Douala produit du yaourt à grande échelle grâce à des ferments lactiques (\emph{Lactobacillus bulgaricus} et \emph{Streptococcus thermophilus}). Décris le procédé, cite deux avantages et une limite de cette biotechnologie, puis conclus.
\tcblower
\textbf{Solution détaillée.}
\begin{enumerate}
\item \textbf{Description du procédé} : il s'agit d'une \textbf{fermentation lactique}. Le lait est pasteurisé puis ensemencé avec des bactéries lactiques qui transforment le lactose en acide lactique selon la réaction : \ce{C6H12O6 -> 2 C3H6O3}. L'acidification (pH \textasciitilde 4,5) fait coaguler les protéines du lait (caséines) : on obtient le yaourt.
\item \textbf{Avantages} : (a) la fermentation \textbf{conserve} le lait, qui se gâterait en quelques heures sous nos climats tropicaux ; (b) le yaourt est plus \textbf{digeste} (le lactose est partiellement dégradé en acide lactique) et la production industrielle permet d'approvisionner les villes à \textbf{moindre coût} et en grande quantité.
\item \textbf{Limite} : la production exige une \textbf{asepsie rigoureuse} : toute contamination par des micro-organismes indésirables (levures, moisissures, bactéries pathogènes) rend le produit impropre à la consommation ; de plus, la chaîne du froid nécessaire à la conservation et au transport coûte cher en énergie.
\item \textbf{Conclusion argumentée} : la production industrielle de yaourt est une biotechnologie globalement bénéfique (sécurité alimentaire, emplois, produit nutritif accessible), à condition de respecter des règles d'hygiène strictes (bonnes pratiques de fabrication, HACCP) et de maintenir la chaîne du froid — c'est le rôle des mesures de biosécurité et des contrôles sanitaires officiels.
\end{enumerate}
\end{exoresolu}

\begin{exercice}[À toi de jouer]
Le gouvernement camerounais autorise des essais de coton Bt (coton génétiquement modifié résistant aux insectes) dans la région du Nord. En t'appuyant sur la méthode en trois temps, rédige un texte de dix à quinze lignes discutant les avantages et les limites de cette biotechnologie, puis donne ton avis argumenté en tenant compte des enjeux pour les producteurs locaux.
\end{exercice}

\begin{corrige}[Exercice]
\textbf{Éléments de correction attendus.}
\begin{enumerate}
\item \textbf{Procédé} : le coton Bt est un OGM (organisme génétiquement modifié) ; un gène \emph{cry} de la bactérie \emph{Bacillus thuringiensis} a été introduit dans son génome par transgenèse, lui permettant de produire une toxine insecticide (protéine Cry) qui tue les chenilles ravageuses (notamment \emph{Helicoverpa armigera}) en se liant à des récepteurs spécifiques de leur tube digestif.
\item \textbf{Avantages (au moins deux)} : réduction drastique des traitements insecticides chimiques (économie pour le producteur, moindre pollution des sols et eaux, moindre exposition des paysans aux produits toxiques) ; rendements plus élevés et plus réguliers grâce à une protection continue de la plante, donc revenus améliorés pour les producteurs du Nord Cameroun.
\item \textbf{Limites (au moins une)} : risque d'apparition d'insectes résistants à la toxine Bt (nécessité de zones refuges en coton non-Bt) ; dépendance vis-à-vis des semenciers (semences brevetées à racheter chaque année, interdiction de réutiliser la récolte) ; dissémination possible du gène \emph{cry} vers des plantes voisines par le pollen ; coût initial des semences plus élevé.
\item \textbf{Avis argumenté} : réponse nuancée, par exemple : le coton Bt peut améliorer significativement les revenus des cotonculteurs du Nord et réduire l'usage de pesticides dangereux, mais sa diffusion doit s'accompagner de mesures de biosécurité strictes (zones refuges obligatoires pour retarder les résistances, suivi environnemental post-commercialisation) et de garanties politiques pour l'accès équitable des petits producteurs aux semences (subventions, licences locales, maintien de variétés conventionnelles).
\end{enumerate}
Toute copie équilibrée, structurée en trois temps et sans affirmation manichéenne, obtient la totalité des points.
\end{corrige}

\begin{aretenir}[L'essentiel de la section]
\begin{itemize}
\item Les biotechnologies offrent des avantages considérables : production massive et moins chère (insuline, vaccins), rendements agricoles améliorés, conservation des aliments, valorisation des déchets (biogaz, compost).
\item Elles présentent aussi des limites : techniques (contamination, coût des équipements, semences brevetées), sanitaires et environnementales (allergies potentielles, dissémination de gènes, résistances aux toxines), et éthiques (brevet du vivant, équité d'accès, biopiraterie).
\item La \cle{biosécurité} regroupe les mesures de prévention de ces risques (confinement, évaluation, étiquetage, Protocole de Carthagène).
\item Pour discuter une biotechnologie : décris le procédé, cite deux avantages, cite au moins une limite, et conclus par un avis argumenté et équilibré.
\end{itemize}
\end{aretenir}

\coursec{Enjeux économiques et sociaux au Cameroun}

Dans la section précédente, nous avons pesé les avantages et les limites des biotechnologies, et nous avons vu qu'elles alimentent des débats légitimes. Mais ces débats ne sont pas abstraits : ils se jouent ici, chez nous, dans les brasseries de Douala, dans les champs de plantain de la région du Centre, dans les laboratoires de Yaoundé et sur les étals de nos marchés. Posons-nous alors une question très concrète : \emph{que rapportent — et que peuvent rapporter — les biotechnologies au Cameroun ?} C'est toute la question des \cle{enjeux} économiques, agricoles, sanitaires et sociaux.

\courssub{Les enjeux économiques : de la matière première à la richesse}

\textbf{Une observation de la vie quotidienne.} Sur le marché Mokolo à Yaoundé, un sac de manioc frais se vend rapidement... et doit l'être, car les tubercules pourrissent en quelques jours. Quelques étals plus loin, des bâtons de \emph{bobolo} (manioc fermenté) se conservent plusieurs semaines et se vendent plus cher au kilogramme. Même matière première, deux valeurs très différentes. D'où vient la différence ? D'un procédé biotechnologique : la \cle{fermentation}.

\begin{exemplebox}[Le manioc transformé : la biotechnologie crée de la richesse]
Ordres de grandeur pédagogiques : un tubercule de manioc frais se conserve \textbf{2 à 3 jours} après récolte ; transformé en bobolo par fermentation (action de bactéries lactiques et de levures naturellement présentes), il se conserve \textbf{plusieurs semaines}. Cette transformation :
\begin{itemize}
\item \textbf{multiplie la durée de conservation}, ce qui réduit les pertes après récolte (un problème majeur en Afrique centrale) ;
\item \textbf{augmente la valeur commerciale} du produit : le kilogramme de bobolo se vend nettement plus cher que le kilogramme de manioc brut ;
\item \textbf{crée des emplois} : épluchage, fermentation, enveloppage dans les feuilles, transport, vente.
\end{itemize}
C'est une illustration directe de la \cle{valeur ajoutée} : la biotechnologie, même traditionnelle, transforme une ressource périssable en richesse durable.
\end{exemplebox}

\begin{definition}[Valeur ajoutée et transformation locale]
La \cle{valeur ajoutée} est la richesse nouvelle créée lorsqu'une matière première est transformée en produit fini. \textbf{Transformer localement} (manioc en bobolo, maïs en bière, lait en yaourt) permet de garder cette richesse dans le pays au lieu d'exporter la matière brute et d'importer le produit fini, plus cher.
\end{definition}

\textbf{À l'échelle industrielle}, le raisonnement est identique. Les grandes brasseries installées au Cameroun (la SABC — Société Anonyme des Brasseries du Cameroun — et Guinness Cameroun) transforment le maïs, le sorgho et le sucre en bières et boissons par fermentation alcoolique industrielle, réalisée par des levures du genre \emph{Saccharomyces}. Elles emploient directement des milliers de salariés (brasseurs, techniciens de fermentation, contrôleurs qualité) et indirectement des dizaines de milliers de personnes (agriculteurs fournisseurs, transporteurs, détaillants). L'industrie agroalimentaire (laiteries, minoteries, huileries) repose sur le même principe : la transformation locale des matières premières agricoles.

Enfin, les biotechnologies médicales ouvrent une perspective stratégique : \textbf{réduire les importations}. Le Cameroun importe aujourd'hui la quasi-totalité de son insuline et de ses vaccins, ce qui coûte cher en devises et rend le pays dépendant de l'extérieur. Produire localement certains médicaments issus des biotechnologies (même sous licence) est un enjeu de souveraineté économique et sanitaire.

\begin{popfigure}
\begin{center}
\begin{tikzpicture}[
  mindmap,
  concept color=popblue,
  every node/.style={concept, text=white, font=\small\bfseries},
  root concept/.append style={font=\normalsize\bfseries},
  level 1/.append style={level distance=3.6cm, sibling angle=90},
  level 2/.append style={level distance=2.6cm, sibling angle=50, concept color=popblueL, text=popdark, font=\scriptsize},
  grow cyclic
]
\node[root concept]{Biotechnologies au Cameroun}
  child[concept color=poporange]{ node {Économie}
    child { node {Brasseries (SABC, Guinness) : emplois} }
  }
  child[concept color=popgreen]{ node {Agriculture}
    child { node {Plants sains de plantain (culture de tissus)} }
  }
  child[concept color=poppurple]{ node {Santé}
    child { node {Vaccins et insuline accessibles} }
  }
  child[concept color=poppink]{ node {Société}
    child { node {Savoir-faire : fermentation du manioc} }
  };
\end{tikzpicture}
\end{center}
\legende{Carte mentale des quatre grands domaines d'enjeux des biotechnologies au Cameroun, avec un exemple concret par branche.}
\end{popfigure}

\courssub{Les enjeux agricoles : nourrir une population croissante}

\textbf{Situation-problème.} Le Cameroun compte plus de 28 millions d'habitants, et sa population augmente vite. Or certaines cultures emblématiques, comme la banane plantain, sont menacées par des maladies (par exemple la cercosporiose noire, maladie fongique due au champignon \emph{Mycosphaerella fijiensis}) qui font chuter les rendements. Comment produire plus, sur des surfaces limitées, sans épuiser les sols ?

Une réponse biotechnologique existe déjà : la \cle{culture de tissus} (ou micropropagation). Le principe est simple : à partir d'un petit fragment d'une plante saine et performante (méristème), on produit en laboratoire, en milieu stérile, des milliers de \textbf{plants identiques et exempts de maladies}. Ces plants sains sont ensuite diffusés aux paysans.

\begin{exemplebox}[Plants sains de banane plantain et de palmier à huile]
Des structures de recherche comme l'\cle{IRAD} (Institut de Recherche Agricole pour le Développement) produisent et diffusent des plants améliorés :
\begin{itemize}
\item \textbf{banane plantain} : plants issus de culture de tissus, vigoureux, sans virus ni champignons, à rendement supérieur ;
\item \textbf{palmier à huile} : matériel végétal sélectionné à haute productivité en huile.
\end{itemize}
Conséquence directe : de meilleurs rendements à l'hectare, donc plus de nourriture disponible et plus de revenus pour le producteur — une arme concrète contre l'\cle{insécurité alimentaire}.
\end{exemplebox}

\begin{attention}[Piège fréquent : réduire la biotechnologie aux OGM]
Quand on entend « biotechnologie agricole », beaucoup pensent immédiatement aux OGM. C'est une erreur : la \textbf{culture de tissus}, la \textbf{sélection variétale} et la \textbf{fermentation des semences} sont des biotechnologies agricoles qui ne font intervenir \emph{aucune} modification génétique. Au Probatoire, cite de préférence un exemple précis (plants de plantain de l'IRAD) plutôt que le terme vague « OGM ».
\end{attention}

\courssub{Les enjeux sanitaires : vaccins, médicaments et lutte contre les maladies}

Les biotechnologies médicales touchent directement la santé des Camerounais :

\begin{itemize}
\item \textbf{Vaccins} : la vaccination (contre la rougeole, la poliomyélite, l'hépatite B, et plus récemment le paludisme avec les premiers vaccins antipaludiques déployés en Afrique, dont le Cameroun fait partie) repose sur des produits issus des biotechnologies modernes. L'enjeu est l'\textbf{accès} : disponibilité, prix, chaîne du froid jusqu'aux centres de santé ruraux.
\item \textbf{Insuline} : le diabète progresse au Cameroun. L'insuline humaine moderne est produite par des bactéries (\emph{Escherichia coli}) ou des levures génétiquement modifiées — un produit biotechnologique vital, mais coûteux car importé.
\item \textbf{Paludisme} : première cause de consultation et de mortalité infantile au Cameroun ; les tests de diagnostic rapide (TDR), produits par biotechnologie (anticorps spécifiques d'antigènes du parasite \emph{Plasmodium}), permettent un diagnostic en quelques minutes dans les centres de santé éloignés.
\end{itemize}

L'enjeu sanitaire est donc double : \textbf{disposer} de ces produits (production ou importation maîtrisée) et les \textbf{rendre accessibles} à toute la population, y compris rurale.

\courssub{Les enjeux sociaux : savoir-faire, formation et débat public}

\textbf{Valorisation des savoir-faire traditionnels.} La fermentation du manioc (bobolo, miondo), celle du maïs ou du mil (bil-bil), la fabrication du lait caillé : ces techniques ancestrales sont des biotechnologies à part entière. Les étudier scientifiquement permet de les \textbf{améliorer} (maîtrise de l'hygiène, standardisation, conservation) et de les transformer en activités économiques structurées, souvent portées par les femmes.

\textbf{Formation.} Le développement des biotechnologies exige des compétences : techniciens de laboratoire, ingénieurs agronomes, chercheurs en microbiologie. Les universités camerounaises et les centres de recherche comme l'\cle{IRAD} et l'\cle{IMPM} (Institut de Recherches Médicales et d'Études des Plantes Médicinales) forment ces cadres et mènent des recherches adaptées aux besoins locaux (plantes médicinales, fermentations, variétés améliorées).

\textbf{Débat public sur les OGM.} Faut-il autoriser la culture d'organismes génétiquement modifiés au Cameroun ? La question concerne tout le monde : consommateurs, agriculteurs, industriels. Un débat public éclairé suppose des citoyens capables de distinguer les faits scientifiques des rumeurs — c'est l'une des raisons d'être de cette leçon.

\begin{savaistu}[Le Cameroun dispose d'une loi sur la biosécurité]
Le Cameroun a ratifié le \textbf{Protocole de Cartagena} sur la prévention des risques biotechnologiques et s'est doté de la \textbf{loi n° 2003/006 du 21 avril 2003} portant régime de la biosécurité. Ce texte encadre l'utilisation, l'importation et la dissémination des organismes génétiquement modifiés sur le territoire national, afin de protéger la santé humaine et la biodiversité. Autrement dit : les biotechnologies modernes ne s'appliquent pas n'importe comment — elles sont réglementées.
\end{savaistu}

\courssub{Analyser un enjeu : une méthode et un exemple résolu}

\begin{methode}[Analyser un enjeu économique ou social d'une biotechnologie]
Face à une biotechnologie donnée, pose-toi toujours \textbf{trois questions} :
\begin{enumerate}
\item \textbf{Qui en bénéficie ?} (consommateurs, agriculteurs, industriels, malades, État...)
\item \textbf{Quelle richesse ou quel service crée-t-elle ?} (emplois, valeur ajoutée, produit plus sain ou mieux conservé, importations évitées...)
\item \textbf{Quels risques ou inégalités peut-elle engendrer ?} (dépendance, coût d'accès, risques sanitaires ou environnementaux, exclusion des petits producteurs...)
\end{enumerate}
Une analyse complète au Probatoire doit mentionner au moins un bénéficiaire, un gain et une limite.
\end{methode}

\begin{exoresolu}[Analyser l'enjeu économique de la transformation du manioc]
\textbf{Énoncé.} Une coopérative de femmes transforme chaque semaine 500 kg de manioc frais en bobolo. Le manioc brut se vend 100 FCFA le kilogramme ; le bobolo se vend 300 FCFA le kilogramme (on admet que 1 kg de manioc donne 1 kg de bobolo). (1) Calculer la valeur ajoutée hebdomadaire créée par la coopérative. (2) Citer un autre avantage de cette transformation. (3) Citer une limite possible.
\tcblower
\textbf{Solution.}
\begin{enumerate}
\item Valeur du manioc brut : $500 \times 100 = \num{50000}$ FCFA. Valeur du bobolo : $500 \times 300 = \num{150000}$ FCFA. Valeur ajoutée :
\[ \num{150000} - \num{50000} = \num{100000} \text{ FCFA par semaine.} \]
La fermentation a donc \textbf{triplé la valeur} de la matière première.
\item Autre avantage : la conservation passe de quelques jours à plusieurs semaines, ce qui \textbf{réduit les pertes} après récolte et permet de vendre sur des marchés éloignés ; de plus, l'activité \textbf{crée des emplois} (épluchage, enveloppage, vente).
\item Limite possible : la transformation demande du \textbf{travail, de l'eau potable et une hygiène rigoureuse} ; une fermentation mal maîtrisée peut altérer le produit. Le prix élevé peut aussi exclure les consommateurs les plus pauvres.
\end{enumerate}
\end{exoresolu}

\begin{popfigure}
\begin{center}
\begin{tikzpicture}
\begin{axis}[
  ybar,
  width=0.85\linewidth,
  height=6cm,
  bar width=1.2cm,
  ymin=0, ymax=400,
  ylabel={Valeur (FCFA / kg)},
  symbolic x coords={Manioc brut, Bobolo (fermenté)},
  xtick=data,
  nodes near coords,
  every node near coord/.append style={font=\small\bfseries},
  ymajorgrids=true,
  grid style={popline, dashed},
  axis line style={popdark},
]
\addplot[fill=poporange, draw=popdark] coordinates {(Manioc brut,100) (Bobolo (fermenté),300)};
\end{axis}
\end{tikzpicture}
\end{center}
\legende{Valeur commerciale comparée du manioc brut et du manioc transformé en bobolo (données fictives pédagogiques, en FCFA par kilogramme) : la fermentation triple la valeur du produit.}
\end{popfigure}

\courssub{Perspectives : vers une bioéconomie camerounaise}

L'ensemble de ces enjeux dessine une perspective d'avenir : construire une \cle{bioéconomie} locale, c'est-à-dire une économie fondée sur la \textbf{transformation biologique des ressources agricoles du pays} plutôt que sur leur exportation brute. Le Cameroun dispose d'atouts considérables : une biodiversité exceptionnelle, des savoir-faire traditionnels vivants, des institutions de recherche (IRAD, IMPM, universités) et un cadre juridique (loi de 2003 sur la biosécurité). Former des techniciens et des chercheurs, soutenir les unités de transformation rurales, maîtriser progressivement la production de certains médicaments : autant de chantiers où les biotechnologies peuvent devenir un levier de développement.

\begin{aretenir}[L'essentiel de la section]
\begin{itemize}
\item Les biotechnologies créent de la \textbf{valeur ajoutée} en transformant localement les matières premières (manioc, maïs, lait) : emplois, revenus, importations réduites.
\item \textbf{Agriculture} : plants sains issus de la culture de tissus (plantain, palmier à huile) pour améliorer les rendements et lutter contre l'insécurité alimentaire.
\item \textbf{Santé} : accès aux vaccins, à l'insuline et aux tests de diagnostic (paludisme, diabète).
\item \textbf{Société} : valorisation des savoir-faire traditionnels, formation de cadres camerounais, débat public encadré par la loi n° 2003/006 sur la biosécurité.
\item Pour analyser un enjeu : \emph{qui en bénéficie ? quelle richesse créée ? quels risques ?}
\end{itemize}
\end{aretenir}

\begin{attention}[Erreur fréquente : croire que la biotechnologie, c'est seulement les grandes industries]
Les biotechnologies ne concernent pas que les brasseries ou les laboratoires pharmaceutiques. La ménagère qui fermente le manioc pour faire le bobolo, le brasseur de bil-bil du village et l'éleveur qui vaccine ses poulets \textbf{pratiquent tous des biotechnologies} : ils utilisent des organismes vivants (levures, bactéries) ou des produits biologiques pour obtenir un bien ou un service. Dans une copie, montre que tu sais relier les biotechnologies traditionnelles et modernes.
\end{attention}

\begin{exercice}[À toi de jouer]
\begin{enumerate}
\item Définis la valeur ajoutée et explique en trois lignes pourquoi la transformation locale du manioc est un enjeu économique pour le Cameroun.
\item Cite deux institutions camerounaises de recherche impliquées dans les biotechnologies et précise le domaine de chacune.
\item En appliquant la méthode des trois questions, analyse l'enjeu social de la production locale de vaccins.
\end{enumerate}
\end{exercice}

\begin{corrige}[Exercice]
\begin{enumerate}
\item La valeur ajoutée est la richesse nouvelle créée par la transformation d'une matière première en produit fini. Transformer le manioc sur place (bobolo, farine fermentée) permet de conserver cette richesse au Cameroun, de créer des emplois et de réduire les pertes liées à la pourriture des tubercules.
\item L'\textbf{IRAD} (Institut de Recherche Agricole pour le Développement) : biotechnologies agricoles (plants sains de plantain, variétés améliorées). L'\textbf{IMPM} (Institut de Recherches Médicales et d'Études des Plantes Médicinales) : recherche médicale et valorisation des plantes médicinales.
\item \emph{Qui en bénéficie ?} Toute la population, surtout les enfants et les zones rurales. \emph{Quel service créé ?} Des vaccins disponibles, moins chers, en quantité maîtrisée ; des emplois qualifiés ; moins de devises dépensées en importations. \emph{Quels risques ?} Coût initial élevé des unités de production, nécessité d'un contrôle qualité strict, risque de dépendance technologique si la production se fait entièrement sous licence étrangère.
\end{enumerate}
\end{corrige}

\newpage

\coursec{Activité d'intégration}

\coursec{Application de la biotechnologie dans la production de biens et services}

\courssub{Notion de biotechnologie}

\begin{definition}[Biotechnologie]
Ensemble des techniques qui utilisent des organismes vivants (microorganismes, cellules animales ou végétales) ou leurs constituants (enzymes, ADN) pour produire des biens ou des services utiles à l'Homme. On distingue les \textbf{biotechnologies traditionnelles} (savoir-faire empiriques, fermentation spontanée) des \textbf{biotechnologies modernes} (génie génétique, culture cellulaire, enzymologie industrielle).
\end{definition}

\begin{exemplebox}[Domaines d'application majeurs]
\begin{itemize}
    \item \textbf{Agroalimentaire :} production de boissons, pain, fromage, yaourt, acides aminés, additifs.
    \item \textbf{Agriculture :} plantes améliorées (résistance aux maladies, tolérance à la sécheresse), biofertilisants, biopesticides, vitroplants.
    \item \textbf{Médecine / Pharmacie :} production d'insuline, vaccins, anticorps monoclonaux, thérapie génique, diagnostic moléculaire (PCR).
    \item \textbf{Environnement :} épuration des eaux usées, biodégradation des polluants (pétrole), valorisation des déchets (biogaz).
    \item \textbf{Énergie :} bioéthanol, biodiesel, biohydrogène à partir de biomasse.
\end{itemize}
\end{exemplebox}

\courssub{Principe d'un procédé biotechnologique : la fermentation}

\begin{propriete}[Fermentation : définition générale]
La fermentation est une voie de dégradation incomplète du glucose (ou d'autres substrats organiques) en l'absence d'oxygène (\textbf{anaérobiose stricte} ou relative), réalisée par des microorganismes (bactéries, levures) ou des cellules musculaires. Elle permet de régénérer le \cle{NAD$^+$} nécessaire à la glycolyse, produisant une faible quantité d'ATP (2 ATP par glucose) et des déchets métaboliques caractéristiques (acide lactique, éthanol + CO$_2$, acides divers).
\end{propriete}

\begin{experience}[Schéma général de la fermentation lactique]
\begin{popfigure}
\begin{tikzpicture}[
    node distance=1.2cm and 1.5cm,
    box/.style={draw=popdark, thick, rounded corners, fill=popcream, minimum width=2.5cm, minimum height=1cm, align=center},
    arrow/.style={-Stealth, thick, popdark},
    condition/.style={draw=popblue, thick, rounded corners, fill=popblueL, minimum width=2.5cm, minimum height=0.8cm, align=center, font=\small}
]
\node[box, align=center] (substrat){Substrat organique\\(Glucose, Amidon,\\Cellulose...)};
\node[condition, below=of substrat, align=center] (cond){Conditions : Anaérobiose\\Température optimale\\pH adapté\\Présence d'eau};
\node[box, right=of cond, xshift=1.5cm, align=center] (micro){Microorganismes\\(Bactéries lactiques,\\Levures, Moisissures)};
\node[box, below=of cond, align=center] (produits){Produits métaboliques\\(Acide lactique, Éthanol,\\CO$_2$, Acide acétique,\\Énergie ATP, Biomasse)};
\node[box, below=of produits, align=center] (effets){Effets sur la matière\\(Ramollissement, Détoxification,\\Conservation, Goût,\\Valeur nutritionnelle)};

\draw[arrow] (substrat) -- (cond);
\draw[arrow] (cond) -- (micro);
\draw[arrow] (micro) -- (produits);
\draw[arrow] (produits) -- (effets);
\draw[arrow] (cond) -| (micro);
\end{tikzpicture}
\legende{Schéma fonctionnel d'une fermentation : le substrat est transformé par les microorganismes sous conditions contrôlées en produits d'intérêt.}
\end{popfigure}
\end{experience}

\begin{propriete}[Équation bilan de la fermentation lactique (homolactique)]
\ce{C6H12O6 ->[$\text{bactéries lactiques}$] 2 CH3CH(OH)COOH + 2 ATP}
\emph{(Glucose $\rightarrow$ 2 Acide lactique + 2 ATP)}
\end{propriete}

\begin{propriete}[Équation bilan de la fermentation alcoolique]
\ce{C6H12O6 ->[levures] 2 C2H5OH + 2 CO2 + 2 ATP}
\emph{(Glucose $\rightarrow$ 2 Éthanol + 2 CO$_2$ + 2 ATP)}
\end{propriete}

\courssub{Biotechnologies traditionnelles : exemples}

\begin{aretenir}[Exemples camerounais de biotechnologies traditionnelles]
\begin{center}
\begin{tabularx}{\linewidth}{|l|X|X|X|}
\hline
\rowcolor{popblueL} \textbf{Produit} & \textbf{Substrat} & \textbf{Microorganismes principaux} & \textbf{Intérêt} \\
\hline
\textbf{Bobolo / Miondo / Bâton de manioc} & Tubercule de manioc (amidon) & Bactéries lactiques (\emph{Lactobacillus}, \emph{Leuconostoc}), Levures (\emph{Candida}, \emph{Saccharomyces}) & Détoxification (cyanure), conservation, goût acidulé, texture pâteuse \\
\hline
\textbf{Vin de palme / Raffia} & Sève de palmier (sucres) & Levures (\emph{Saccharomyces}), Bactéries acétiques & Boisson alcoolisée, source de vitamines B, valeur culturelle \\
\hline
\textbf{Bière traditionnelle (Bilibili, Nkang)} & Mil, maïs, sorgho (amidon) & Levures (\emph{Saccharomyces}), Amylases (maltage) & Boisson fermentée, apport énergétique, lien social \\
\hline
\textbf{Yaourt local (Lait caillé)} & Lait (lactose) & Bactéries lactiques (\emph{Streptococcus}, \emph{Lactobacillus}) & Conservation du lait, digestibilité, probiotiques \\
\hline
\textbf{Pain (levain)} & Farine de blé (gluten, amidon) & Levures (\emph{Saccharomyces}), Bactéries lactiques & Levée de la pâte, arômes, conservation \\
\hline
\end{tabularx}
\end{center}
\end{aretenir}

\courssub{Biotechnologies modernes : agriculture et médecine}

\begin{definition}[Biotechnologies modernes]
Ensemble des techniques récentes (depuis les années 1970-80) faisant appel au \textbf{génie génétique} (ADN recombinant, PCR, CRISPR-Cas9), à la \textbf{culture de cellules/tissus in vitro} et à l'\textbf{enzymologie de purification} pour obtenir des produits à haute valeur ajoutée ou des organismes aux caractéristiques nouvelles.
\end{definition}

\begin{exemplebox}[Applications modernes au Cameroun et dans le monde]
\begin{itemize}
    \item \textbf{Agriculture :}
    \begin{itemize}
        \item \textbf{Vitroplants (IRAD, CJARC) :} multiplication saine et rapide du bananier/plantain, du manioc, de l'ananas, du palmier à huile par microbouturage \textit{in vitro}. Garantie sans virus (ex. virus de la mosaïque du manioc).
        \item \textbf{Plantes transgéniques (recherche/import) :} coton Bt (résistant aux insectes), maïs tolérant à la sécheresse (projet TELA/WEMA), riz « Golden Rice » (vitamine A).
        \item \textbf{Biofertilisants / Biopesticides :} \emph{Bacillus thuringiensis} (Bt), \emph{Trichoderma}, \emph{Rhizobium} (fixation d'azote).
    \end{itemize}
    \item \textbf{Médecine / Santé :}
    \begin{itemize}
        \item \textbf{Protéines recombinantes :} Insuline humaine (levure/\emph{E. coli}), Hormone de croissance, Facteurs de coagulation, Vaccins à sous-unités (Hépatite B, VPH), Vaccins à ARNm (COVID-19).
        \item \textbf{Diagnostic moléculaire :} PCR (VIH, Tuberculose, COVID-19, HPV), séquençage (surveillance épidémique).
        \item \textbf{Thérapies innovantes :} Anticorps monoclonaux (cancers, maladies auto-immunes), Thérapie génique (drépanocytose, amaurose de Leber), Cellules CAR-T.
    \end{itemize}
    \item \textbf{Environnement / Industrie :}
    \begin{itemize}
        \item Épuration biologique des eaux usées (boues activées, lagunage).
        \item Bioremédiation sols pollués (hydrocarbures).
        \item Enzymes industrielles (amylases, protéases, cellulases) pour détergents, textile, papier, bioéthanol 2G.
    \end{itemize}
\end{itemize}
\end{exemplebox}

\begin{savaistu}[La biotechnologie au Cameroun : acteurs clés]
\begin{itemize}
    \item \textbf{IRAD (Institut de Recherche Agricole pour le Développement) :} programmes de vitroplants (bananier, manioc, palmier), amélioration variétale conventionnelle et biotechnologique.
    \item \textbf{IMPM (Institut de Recherche Médicale et d'Études des Plantes Médicinales) :} valorisation de la pharmacopée traditionnelle, recherche sur le VIH/Sida, paludisme, tuberculose ; diagnostic moléculaire.
    \item \textbf{Universités (Yaoundé I, Ngaoundéré, Dschang, Buea) :} laboratoires de biologie moléculaire, microbiologie, génétique.
    \item \textbf{Secteur privé :} brasseries (brassage contrôlé, levures sélectionnées), laiteries (yaourts industriels), unités de transformation artisanale modernisée (bobolo emballé, farines composites).
\end{itemize}
\end{savaistu}

\courssub{Avantages, limites et débats}

\begin{propriete}[Analyse comparative : Traditionnel vs Moderne]
\begin{center}
\begin{tabularx}{\linewidth}{|l|X|X|}
\hline
\rowcolor{popblueL} \textbf{Critère} & \textbf{Biotechnologies traditionnelles} & \textbf{Biotechnologies modernes} \\
\hline
\textbf{Contrôle du processus} & Faible (fermentation spontanée, variability) & Élevé (souches pures, bioréacteurs contrôlés, paramètres surveillés) \\
\hline
\textbf{Qualité / Sécurité} & Variable, risques de contamination (mycotoxines, pathogènes) & Standardisée, traçable, normes HACCP/ISO \\
\hline
\textbf{Rendement / Productivité} & Faible, lent & Élevé, rapide, optimisé (ingénierie métabolique) \\
\hline
\textbf{Coût / Accessibilité} & Faible, accessible aux ruraux (peu d'équipement) & Élevé (investissement, expertise, énergie, intrants) \\
\hline
\textbf{Impact environnemental} & Généralement faible (peu d'intrants chimiques) & Variable (consommation énergie/eau, OGM, déchets plastiques) \\
\hline
\textbf{Acceptabilité sociale} & Forte (patrimoine culinaire, confiance) & Débattue (OGM, éthique, brevetage du vivant, « naturel ») \\
\hline
\end{tabularx}
\end{center}
\end{propriete}

\begin{attention}[Débats éthiques et sociétaux (niveau Probatoire)]
\begin{itemize}
    \item \textbf{OGM :} Dissémination dans l'environnement, résistance des insectes/mauvaises herbes, dépendance aux semenciers (brevets), étiquetage.
    \item \textbf{Santé :} Accès aux médicaments innovants (coût), essais cliniques (consentement, populations vulnérables), modification du génome humain (lignée germinale interdite).
    \item \textbf{Biodiversité :} Biopiraterie (accès aux ressources génétiques et partage des avantages - Protocole de Nagoya), érosion génétique par monocultures.
    \item \textbf{Environnement :} Déchets plastiques (emballages biotech), consommation énergétique des bioréacteurs.
\end{itemize}
\end{attention}

\courssub{Enjeux économiques et sociaux au Cameroun}

\begin{aretenir}[Enjeux majeurs pour le Cameroun]
\begin{itemize}
    \item \textbf{Sécurité alimentaire :} Réduction des pertes post-récolte (30-40\% pour les tubercules) par transformation biotechnologique (farines, fécule, bobolo stable). Valorisation des cultures locales (manioc, plantain, maïs, patate douce).
    \item \textbf{Emploi et revenus :} Filières de transformation artisanale et semi-industrielle (coopératives féminines, PME agroalimentaires). Création de valeur ajoutée locale (export de produits transformés vs bruts).
    \item \textbf{Santé publique :} Production locale de réactifs de diagnostic, vaccins, médicaments essentiels (réduction facture d'importation, souveraineté sanitaire). Lutte contre les maladies négligées (onchocercose, bilharziose).
    \item \textbf{Environnement et énergie :} Biogaz à partir de déchets agricoles/urbains (Yaoundé, Douala), bioéthanol (canne à sucre, manioc), gestion durable des forêts (bioprospection).
    \item \textbf{Formation et recherche :} Nécessité de former techniciens/ingénieurs (biologistes moléculaires, bioinformaticiens, ingénieurs procédés) pour maîtriser les outils et ne pas rester simples consommateurs.
\end{itemize}
\end{aretenir}

\courssub{Activité d'intégration : La coopérative « Main Verte » de Mbalmayo}

\textbf{Objectifs de l'activité}

À l'issue de cette activité d'intégration, tu seras capable de :

\begin{itemize}
\item \cle{Décrire} de manière rigoureuse le principe d'un procédé biotechnologique simple : la \cle{fermentation} du manioc (substrat, microorganismes, conditions, produits) ;
\item \cle{Citer} et exploiter des exemples locaux d'application de la biotechnologie dans la vie quotidienne camerounaise ;
\item \cle{Discuter} les avantages et les limites de deux options d'amélioration d'un procédé traditionnel, en comparant leurs effets sur la qualité, la conservation et le coût ;
\item \cle{Calculer} un gain de revenu à partir de données chiffrées simples et en tirer une conclusion économique ;
\item \cle{Formuler} un avis argumenté, rédigé dans un français clair, à destination d'acteurs réels (une coopérative de productrices), c'est-à-dire mobiliser à la fois des savoirs scientifiques et des savoir-être (conseil, honnêteté, sens de l'intérêt collectif).
\end{itemize}

\textbf{Situation}

\begin{exemplebox}[La coopérative « Main Verte » de Mbalmayo]
À Mbalmayo, ville de la région du Centre située à une cinquantaine de kilomètres de Yaoundé, la coopérative de productrices « Main Verte » regroupe 40 femmes qui transforment le manioc en \textbf{bobolo} (bâtons de pâte de manioc fermentée, enveloppés dans des feuilles), un aliment de base très consommé au Cameroun, notamment avec le ndolé ou le poisson braisé.

Le procédé traditionnel de la coopérative est le suivant : les tubercules de manioc sont épluchés, lavés, puis mis à tremper dans des cuvées d'eau pendant 3 à 4 jours. Pendant ce trempage, des \textbf{microorganismes} (bactéries lactiques et levures naturellement présentes) fermentent les tubercules : les fibres se ramollissent, les composés toxiques (dérivés cyanés) sont détruits, et une légère acidité se développe. Les tubercules ramollis sont ensuite égouttés, écrasés en pâte, façonnés en bâtons, enveloppés dans des feuilles et cuits à la vapeur.

\textbf{Le problème :} la coopérative produit chaque mois \textbf{500 bâtons} de bobolo, vendus \textbf{100 FCFA l'unité} sur les marchés de Mbalmayo et de Yaoundé. Mais \textbf{20 \% de la production} est perdue : fermentation mal contrôlée (cuvées mal nettoyées, eau non renouvelée) donnant un goût désagréable, et surtout \textbf{mauvaise conservation} (bâtons stockés à l'air libre, moisissures après quelques jours). Ces pertes obligent les femmes à vendre le reste à bas prix, rapidement, par crainte de tout perdre.

Deux options d'amélioration sont envisagées :
\begin{itemize}
\item \textbf{Option A — Amélioration de l'hygiène de fermentation} : nettoyage systématique des cuvées, renouvellement de l'eau de trempage, tri des tubercules sains. Coût estimé : faible (environ \SI{5 000}{FCFA/mois}).
\item \textbf{Option B — Emballage moderne} : conditionnement des bâtons cuits sous sachets plastiques alimentaires fermés, ce qui prolonge la conservation à plusieurs semaines et permet de vendre en ville à meilleur prix. Coût estimé : plus élevé (environ \SI{15 000}{FCFA/mois}).
\end{itemize}

La présidente de la coopérative sollicite, par l'intermédiaire du proviseur, les élèves de la classe de Première A pour obtenir un \textbf{avis argumenté} : quelle option privilégier, et pourquoi ?
\end{exemplebox}

\textbf{Consignes}

Par groupes de 4 élèves, réponds aux tâches suivantes, puis prépare une restitution orale de 5 minutes.

\begin{enumerate}
\item \textbf{Décrire le procédé.} Rédige un texte de 8 à 10 lignes décrivant le principe de la fermentation du manioc en précisant : le \cle{substrat}, les \cle{microorganismes} impliqués, les \cle{conditions} nécessaires (anaérobiose, température, durée) et les \cle{produits} obtenus. Explique en quoi ce procédé est une \cle{biotechnologie traditionnelle}.
\item \textbf{Comparer les options.} Recopie et complète le tableau ci-dessous en indiquant au moins deux avantages et deux limites pour chaque option.
\begin{center}
\begin{tabularx}{\linewidth}{|l|X|X|}
\hline
\rowcolor{popblueL} \textbf{Critère} & \textbf{Option A : hygiène de fermentation} & \textbf{Option B : emballage moderne} \\
\hline
Avantages & & \\
\hline
Limites & & \\
\hline
\end{tabularx}
\end{center}
\item \textbf{Calculer le gain économique.} On suppose que, grâce aux améliorations, les pertes passent de 20 \% à 5 \% de la production mensuelle (500 bâtons à 100 FCFA l'unité).
\begin{enumerate}
\item Calcule le revenu mensuel actuel (avec 20 \% de pertes).
\item Calcule le revenu mensuel espéré (avec 5 \% de pertes).
\item Calcule le gain mensuel, puis déduis-en si le coût de l'option A (\SI{5 000}{FCFA/mois}) est rentable.
\end{enumerate}
\item \textbf{Rédiger l'avis.} Rédige un court avis (10 à 15 lignes) à remettre à la coopérative : il doit contenir un conseil clair (option A, B ou les deux), justifié par des arguments scientifiques (qualité et sécurité de l'aliment) et économiques (calculs de la question 3).
\end{enumerate}

\textbf{Ressources à mobiliser}

\begin{aretenir}[Ce qu'il faut avoir en tête]
\begin{itemize}
\item \textbf{Biotechnologie} : ensemble des techniques qui utilisent des organismes vivants (microorganismes, cellules) ou leurs constituants (enzymes) pour produire des biens ou des services utiles à l'Homme.
\item \textbf{Fermentation} : transformation d'une matière organique (le \cle{substrat}) par des microorganismes (bactéries, levures), généralement en l'absence d'oxygène, produisant des acides, de l'alcool ou du gaz carbonique. Schéma général de la fermentation lactique :
\[ \ce{C6H12O6 ->[$\text{bactéries lactiques}$] 2 CH3CH(OH)COOH + 2 ATP} \]
\item La fermentation du manioc permet : le \textbf{ramollissement} des tubercules, la \textbf{détoxification} (destruction des dérivés cyanés toxiques : linamarine $\rightarrow$ acide cyanhydrique volatil) et le développement du \textbf{goût caractéristique}.
\item \textbf{Pourcentages} : revenu $=$ quantité vendue $\times$ prix unitaire ; quantité vendue $=$ production $\times (1 - \text{taux de pertes})$.
\item \textbf{Analyse avantages/limites} : tout procédé biotechnologique se juge sur la qualité du produit, le coût, la durée de conservation, la sécurité sanitaire et l'accessibilité pour les producteurs.
\end{itemize}
\end{aretenir}

\textbf{Démarche et corrigé commenté}

\begin{exoresolu}[Conseil à la coopérative « Main Verte »]
\textbf{Énoncé.} Une coopérative de Mbalmayo produit 500 bâtons de bobolo par mois (100 FCFA l'unité) et en perd 20 \% par mauvaise conservation. (1) Décrire le procédé de fermentation du manioc. (2) Comparer les avantages et limites de l'amélioration de l'hygiène de fermentation (option A, 5 000 FCFA/mois) et de l'emballage moderne (option B, 15 000 FCFA/mois). (3) Calculer le gain de revenu si les pertes passent de 20 \% à 5 \%. (4) Rédiger un avis argumenté.

\tcblower

\textbf{Étape 1 — Description du procédé de fermentation du manioc.}

Le \textbf{substrat} est le tubercule de manioc, riche en amidon (polymère du glucose) et contenant des glucosides cyanogènes (linamarine, lotaustraline). Après épluchage et lavage, les tubercules sont trempés dans l'eau pendant 3 à 4 jours. Les \textbf{microorganismes} sont des bactéries lactiques (\emph{Lactobacillus} spp., \emph{Leuconostoc} spp.) et des levures (\emph{Candida} spp., \emph{Saccharomyces} spp.) naturellement présentes sur les tubercules et dans l'environnement : on parle de fermentation spontanée (ou « à l'inoculum naturel »). Les \textbf{conditions} sont : milieu aqueux, pauvre en oxygène (\textbf{anaérobiose}), température ambiante tropicale (\SIrange{25}{30}{\celsius}), pH s'acidifiant progressivement (vers 4,0-4,5), durée de \SIrange{3}{4}{jours}. Les microorganismes hydrolysent l'amidon en sucres fermentescibles, puis dégradent ces sucres principalement en \textbf{acide lactique} (fermentation lactique) et en éthanol/CO$_2$ (fermentation alcoolique). L'acidification du milieu active les enzymes endogènes (linamarases) qui hydrolysent les glucosides cyanogènes, libérant l'acide cyanhydrique (HCN) volatil qui s'échappe : c'est la \textbf{détoxification}. Simultanément, les parois cellulaires (pectines, cellulose) sont ramollies par l'action enzymatique (pectinases, cellulases microbiennes). La pâte obtenue est ensuite façonnée en bâtons, enveloppée dans des feuilles (bananier, marantacée) et cuite à la vapeur : c'est le bobolo.

Ce procédé est une \textbf{biotechnologie traditionnelle} car il utilise des organismes vivants pour transformer une denrée, selon un savoir-faire empirique transmis de génération en génération, sans équipement de laboratoire ni souches sélectionnées.

\textbf{Étape 2 — Tableau avantages/limites.}

\begin{center}
\begin{tabularx}{\linewidth}{|l|X|X|}
\hline
\rowcolor{popblueL} \textbf{Critère} & \textbf{Option A : hygiène de fermentation} & \textbf{Option B : emballage moderne} \\
\hline
Avantages & Coût très faible (\SI{5 000}{FCFA/mois}) ; améliore le goût et la régularité de la qualité organoleptique ; réduit fortement les risques sanitaires (élimination germes pathogènes, mycotoxines) ; valorise le savoir-faire local, aucune dépendance extérieure (intrants locaux). & Conservation prolongée (plusieurs semaines vs quelques jours) ; produit présentable, hygiénique, vendable en supermarché/ville à meilleur prix ; protection efficace contre les moisissures, insectes, poussière et manipulation. \\
\hline
Limites & Ne résout pas le problème de conservation \textbf{après cuisson} (le bobolo cuit nu moisit vite) ; exige une discipline quotidienne stricte (nettoyage, renouvellement eau) ; gains limités si le stockage final reste défaillant. & Coût élevé pour la trésorerie de la coopérative (\SI{15 000}{FCFA/mois}) ; nécessite un approvisionnement régulier en sachets alimentaires (dépendance fournisseurs) ; génère des déchets plastiques (problème environnemental non négligeable) ; efficace seulement si le produit entrant est déjà de bonne qualité (goût, sécurité). \\
\hline
\end{tabularx}
\end{center}

\textbf{Étape 3 — Calculs économiques.}

Revenu actuel (20 \% de pertes) : la coopérative vend $500 \times (1 - 0{,}20) = 500 \times 0{,}80 = 400$ bâtons.
\[ R_{\text{actuel}} = 400 \times \SI{100}{FCFA} = \SI{40 000}{FCFA/mois}. \]

Revenu espéré (5 \% de pertes) : la coopérative vend $500 \times (1 - 0{,}05) = 500 \times 0{,}95 = 475$ bâtons.
\[ R_{\text{espéré}} = 475 \times \SI{100}{FCFA} = \SI{47 500}{FCFA/mois}. \]

Gain mensuel brut (augmentation du chiffre d'affaires) :
\[ G_{\text{brut}} = \SI{47 500}{FCFA} - \SI{40 000}{FCFA} = \SI{7 500}{FCFA/mois}. \]

Analyse de rentabilité de l'Option A :
Bénéfice net mensuel Option A $= G_{\text{brut}} - \text{Coût Option A} = \SI{7 500}{FCFA} - \SI{5 000}{FCFA} = \SI{2 500}{FCFA/mois}$.
L'option A est donc \textbf{rentable dès le premier mois} (retour sur investissement immédiat).

Analyse de rentabilité de l'Option B :
Bénéfice net mensuel Option B $= G_{\text{brut}} - \text{Coût Option B} = \SI{7 500}{FCFA} - \SI{15 000}{FCFA} = -\SI{7 500}{FCFA/mois}$.
L'option B \textbf{n'est pas rentable} au prix de vente actuel (100 FCFA). Elle ne le deviendrait que si l'emballage permet d'augmenter le prix de vente (ex. \SI{120}{FCFA} en ville) : nouveau revenu $= 475 \times 120 = \SI{57 000}{FCFA}$ ; gain brut $= \SI{17 000}{FCFA}$ ; net $= \SI{2 000}{FCFA}$ (rentable mais risqué car dépend de l'acceptation du marché).

\textbf{Étape 4 — Avis rédigé (modèle).}

« À l'attention de la présidente de la coopérative Main Verte.

Objet : Avis sur les options d'amélioration de la production de bobolo.

Madame la Présidente,

Après analyse de votre situation, nous vous conseillons de \textbf{privilégier l'Option A (amélioration de l'hygiène de fermentation) comme priorité immédiate}.

\textbf{Argument scientifique :} Une fermentation propre et contrôlée (cuvées nettoyées, eau renouvelée, tri des tubercules) est la condition \textit{sine qua non} de la qualité et de la sécurité du bobolo. Elle garantit une acidification correcte (pH < 4,5), une détoxification efficace des dérivés cyanés et l'absence de germes pathogènes ou de mycotoxines. Sans cette base, tout emballage moderne ne ferait qu'enfermer un produit de mauvaise qualité, voire dangereux.

\textbf{Argument économique :} Nos calculs montrent qu'en réduisant les pertes de 20 \% à 5 \%, votre revenu mensuel passerait de \SI{40 000}{FCFA} à \SI{47 500}{FCFA}, soit un gain brut de \SI{7 500}{FCFA}. Ce gain couvre largement le coût de l'Option A (\SI{5 000}{FCFA}), dégageant un bénéfice net de \SI{2 500}{FCFA} dès le premier mois.

Dans un second temps, lorsque la qualité sera stable et que des économies auront été constituées, l'emballage moderne (Option B) pourra être envisagé pour conquérir les marchés urbains (supermarchés, Yaoundé) à un prix majoré, à condition de mettre en place une filière de collecte/recyclage des sachets usagés pour limiter l'impact environnemental.

Nous restons à votre disposition pour tout appui technique.

Les élèves de Première A, option SVTEEHB. »
\end{exoresolu}

\begin{attention}[Erreurs fréquentes à éviter]
\begin{itemize}
\item Calculer le revenu sur 500 bâtons sans retirer les pertes : le revenu ne porte que sur les bâtons \textbf{réellement vendus}.
\item Confondre \textbf{gain brut} (\SI{7 500}{FCFA}) et \textbf{bénéfice net} (gain moins le coût de l'option choisie).
\item Décrire la fermentation comme une « cuisson » ou une simple « macération » : c'est une transformation \textbf{biologique} active par des microorganismes, sans apport de chaleur externe pendant la phase de trempage.
\item Oublier de justifier le conseil final par \textbf{à la fois} un argument scientifique (qualité, sécurité, détoxification) et un argument économique chiffré (rentabilité).
\item Négliger l'aspect « déchets plastiques » dans l'évaluation de l'Option B : tout projet biotechnologique moderne doit intégrer la durabilité environnementale.
\end{itemize}
\end{attention}

\textbf{Bilan de l'activité}

Cette activité t'a permis de mettre en évidence plusieurs idées essentielles de la leçon :

\begin{itemize}
\item Une \textbf{biotechnologie traditionnelle} comme la fermentation du manioc repose sur des mécanismes biologiques précis (microorganismes, substrat, conditions d'anaérobiose, hydrolyse enzymatique, détoxification) que la science permet de comprendre et d'améliorer sans renier le savoir-faire local ;
\item L'amélioration d'un procédé biotechnologique se décide par une \textbf{analyse objective des avantages et des limites}, et non par intuition : l'option la plus « moderne » n'est pas toujours la plus rentable à court terme, ni la plus adaptée au contexte local ;
\item Les biotechnologies ont de véritables \textbf{enjeux économiques et sociaux au Cameroun} : quelques pourcents de pertes en moins représentent des milliers de FCFA de revenu supplémentaire pour des productrices, donc une amélioration concrète des conditions de vie et de la sécurité alimentaire ;
\item Enfin, tu as exercé un \textbf{savoir-être} : conseiller des acteurs réels exige rigueur scientifique, honnêteté dans les calculs, clarté dans l'expression et sens de l'intérêt collectif — des compétences utiles bien au-delà de la classe.
\end{itemize}

\newpage

\coursec{Méthodes et exercices résolus}

\coursec{Application de la biotechnologie dans la production de biens et services}

\courssub{Notion de biotechnologie}

\begin{definition}[Biotechnologie]
On appelle \cle{biotechnologie} toute technique qui utilise des \textbf{organismes vivants} (micro-organismes, cellules animales ou végétales) ou leurs \textbf{constituants} (enzymes, ADN, organites) pour \textbf{produire des biens ou des services} utiles à l'Homme (aliments, médicaments, énergie, dépollution, variétés améliorées...).
\end{definition}

\begin{propriete}[Classification des biotechnologies]
\begin{itemize}
\item \textbf{Biotechnologies traditionnelles (ou classiques) :} procédés empiriques, ancestraux, mis en œuvre sans connaissance moléculaire (fermentations alimentaires, sélection massale des plantes et animaux). Exemples : vin de palme, bière de mil, yaourt, fromage, pain, compostage.
\item \textbf{Biotechnologies modernes :} procédés fondés sur la biologie moléculaire et la génétique (génie génétique, culture de tissus \emph{in vitro}, enzymologie industrielle, bio-informatique). Exemples : insuline humaine recombinante, vaccins à ARN messager, plantes transgéniques (OGM), diagnostic PCR, bioremédiation.
\end{itemize}
\end{propriete}

\begin{methode}[Reconnaître et définir une biotechnologie (QCM, définitions)]
\begin{enumerate}
\item \cle{Définir} toujours la biotechnologie avec ses trois piliers : \textbf{organisme vivant (ou constituant)} + \textbf{procédé technique maîtrisé} + \textbf{bien ou service utile}.
\item \cle{Distinguer} traditionnel (fermentation empirique, sélection phénotypique) et moderne (ADN recombinant, culture cellulaire, génie enzymatique).
\item \cle{Éliminer les fausses pistes} : cuisson, séchage, salage, mécanisation, engrais chimiques, pesticides de synthèse ne sont \emph{pas} des biotechnologies (aucun organisme vivant ne réalise la transformation).
\item Justifier chaque réponse par \cle{un exemple concret}.
\end{enumerate}
\end{methode}

\begin{exoresolu}[QCM : qu'est-ce qu'une biotechnologie ?]
Énoncé. Pour chaque affirmation, réponds par Vrai ou Faux en justifiant.
\begin{enumerate}
\item La fabrication du vin de palme à partir de la sève de raphia est une biotechnologie.
\item Le séchage du manioc au soleil est une biotechnologie.
\item La production d'insuline par une bactérie génétiquement modifiée est une biotechnologie moderne.
\item Les biotechnologies n'existent que dans les laboratoires des pays industrialisés.
\end{enumerate}
\tcblower
\textbf{Solution détaillée.}
\begin{enumerate}
\item \textbf{Vrai.} La sève de raphia, riche en sucres, est transformée en vin de palme par des \textbf{levures} naturellement présentes qui réalisent une fermentation alcoolique : utilisation d'organismes vivants pour obtenir un bien utile. C'est une biotechnologie traditionnelle.
\item \textbf{Faux.} Le séchage est une simple déshydratation physique par la chaleur solaire : aucun micro-organisme ne transforme le produit. Ce n'est pas une biotechnologie (c'est une technique de conservation physique).
\item \textbf{Vrai.} Le gène humain de l'insuline est inséré dans une bactérie (génie génétique / ADN recombinant) qui produit la protéine en grande quantité : exemple type de biotechnologie moderne (biotechnologie médicale ou « rouge »).
\item \textbf{Faux.} Les biotechnologies traditionnelles (fermentation du manioc, du lait, du mil, vin de palme) sont pratiquées depuis des siècles dans les villages camerounais, sans laboratoire. Les biotechnologies sont aussi bien artisanales qu'industrielles.
\end{enumerate}
\textbf{Conclusion :} le critère décisif est l'intervention d'un organisme vivant (ou de ses enzymes) dans un procédé de production, qu'il soit ancestral ou moderne.
\end{exoresolu}

\courssub{Principe d'un procédé biotechnologique : la fermentation}

\begin{definition}[Fermentation]
La \cle{fermentation} est une dégradation incomplète de la matière organique (glucides) par des micro-organismes (levures, bactéries) \emph{en l'absence d'oxygène} (milieu anaérobie), libérant une faible quantité d'énergie (ATP) et produisant des métabolites caractéristiques (éthanol, \ce{CO2}, acide lactique, acide acétique...).
\end{definition}

\begin{methode}[Décrire le principe de la fermentation en 5 étapes (standard Probatoire)]
\begin{enumerate}
\item \cle{Identifier le micro-organisme} responsable (levure \emph{Saccharomyces}, bactéries lactiques \emph{Lactobacillus}, \emph{Streptococcus}...) et le \cle{substrat} (glucose, lactose, saccharose, amidon hydrolysé, jus de fruit, pâte de manioc...).
\item \cle{Préciser les conditions} : milieu \textbf{anaérobie} (absence d'\ce{O2}), température optimale (souvent \SIrange{25}{45}{\celsius}), pH favorable.
\item \cle{Écrire l'équation-bilan} de la transformation : substrat $\rightarrow$ produits + énergie (ATP).
\item \cle{Citer les produits obtenus} et leur utilité (éthanol, \ce{CO2}, acide lactique, biomasse...).
\item \cle{Conclure} en reliant le procédé à son application concrète (aliment, boisson, conservation, énergie).
\end{enumerate}
\textbf{Équations-bilan à connaître par cœur :}
\[ \text{Fermentation alcoolique : } \ce{C6H12O6 -> 2 C2H5OH + 2 CO2} + \text{énergie (2 ATP)} \]
\[ \text{Fermentation lactique : } \ce{C6H12O6 -> 2 C3H6O3} + \text{énergie (2 ATP)} \]
\end{methode}

\begin{experience}[Fabrication artisanale du yaourt (TP réalisable en classe)]
\textbf{Matériel :} lait entier, yaourt nature du commerce (ferment), thermomètre, casserole, pots en verre, couverture isotherme ou yaourtière.\\
\textbf{Protocole :}
\begin{enumerate}
\item Porter le lait à ébullition (\SI{90}{\celsius} pendant 5 min) pour le pasteuriser (détruire les microbes indésirables et dénaturer les protéines pour une meilleure texture).
\item Refroidir à \SI{45}{\celsius} (température optimale des bactéries lactiques).
\item Ensemencer : ajouter \SI{2}{\%} de yaourt du commerce (apport de \emph{Streptococcus thermophilus} et \emph{Lactobacillus bulgaricus} vivants), mélanger.
\item Répartir en pots, fermer et maintenir au chaud (\SI{42}{\celsius}--\SI{45}{\celsius}) pendant 4 à 6 h (yaourtière ou couverture épaisse).
\item Réfrigérer \SI{4}{\celsius} pour stopper la fermentation.
\end{enumerate}
\textbf{Observations :} le lait liquide se coagule (consistance gélifiée), pH passe de $\sim 6,6$ à $\sim 4,2$, goût acidulé apparaît.\\
\textbf{Interprétation :} les bactéries lactiques hydrolysent le lactose (glucose + galactose) puis le fermentent en acide lactique (fermentation lactique). L'acidification provoque la précipitation des caséines (protéines du lait) : coagulation. L'acidité inhibe les germes de contamination : conservation.
\end{experience}

\begin{exoresolu}[La fabrication du yaourt : analyse complète]
Énoncé. Au lycée de Ngaoundéré, un club scientifique fabrique du yaourt : du lait est chauffé, refroidi vers \SI{45}{\celsius}, ensemencé avec un peu de yaourt du commerce, puis maintenu au chaud pendant quelques heures. Le lait se coagule et prend un goût acidulé.
\begin{enumerate}
\item Nomme le procédé biotechnologique mis en jeu et les micro-organismes responsables.
\item Écris l'équation-bilan de la transformation du lactose.
\item Explique le rôle du yaourt du commerce ajouté au lait.
\item Pourquoi le lait coagule-t-il ?
\end{enumerate}
\tcblower
\textbf{Solution détaillée.}
\begin{enumerate}
\item \textbf{Procédé :} \cle{fermentation lactique}. \textbf{Micro-organismes :} bactéries lactiques thermophiles : \emph{Streptococcus thermophilus} et \emph{Lactobacillus delbrueckii} subsp. \emph{bulgaricus}.
\item \textbf{Équation-bilan :} Le lactose (\ce{C12H22O11}) est hydrolysé en glucose et galactose (chacun \ce{C6H12O6}), qui subissent la fermentation lactique :
\[ \ce{C12H22O11 + H2O -> 4 C3H6O3} + \text{énergie (ATP)} \]
(Globalement : 1 lactose $\rightarrow$ 4 acides lactiques).
\item \textbf{Rôle du yaourt du commerce :} il sert de \cle{ferment} (ou levain) : il apporte les souches pures, vivantes et en nombre suffisant des deux bactéries lactiques symbiotiques. Sans ensemencement, le lait pasteurisé est stérile et ne fermente pas (ou fermente mal par contamination).
\item \textbf{Coagulation :} l'acide lactique accumulé abaisse le pH jusqu'au point isoélectrique des caséines (pH $\approx 4,6$). Les micelles de caséines perdent leur charge négative, s'agrègent et forment un réseau gélifié tridimensionnel piégeant l'eau et les globules gras : c'est le \textbf{caillé} (yaourt).
\end{enumerate}
\textbf{Conclusion :} la fabrication du yaourt est une fermentation lactique réalisée par un consortium bactérien ensemencé, transformant le lactose en acide lactique en anaérobiose, ce qui acidifie, coagule et conserve le lait.
\end{exoresolu}

\courssub{Biotechnologies traditionnelles : exemples locaux}

\begin{methode}[Construire une réponse « exemples locaux » bien structurée]
\begin{enumerate}
\item \cle{Classer} en : biotechnologies \textbf{traditionnelles} (fermentations artisanales) et \textbf{modernes} (agricoles, médicales, industrielles).
\item Pour chaque exemple, donner le \cle{triplet} : \textbf{produit obtenu} -- \textbf{micro-organisme ou technique} -- \textbf{matière première locale}.
\item \cle{Ancrer dans le contexte camerounais} : produits réels fabriqués localement (boissons, aliments, variétés, vaccins).
\item Terminer par une phrase de \cle{synthèse} sur l'intérêt économique ou social.
\end{enumerate}
\textbf{Rappel :} un exemple « local » = produit \emph{fabriqué} au Cameroun à partir de ressources locales (maïs, manioc, mil, sorgho, lait, sève de raphia, banane plantain...) par un procédé biologique.
\end{methode}

\begin{exoresolu}[Biotechnologies traditionnelles au Cameroun]
Énoncé. Cite trois exemples de biotechnologies traditionnelles pratiquées au Cameroun, en précisant pour chacun la matière première, le micro-organisme principal et le type de fermentation.
\tcblower
\textbf{Solution détaillée.}
\begin{center}
\begin{tabularx}{\linewidth}{|l|X|X|X|}\hline
\rowcolor{popblueL} \textbf{Produit (région)} & \textbf{Matière première} & \textbf{Micro-organismes clés} & \textbf{Type de fermentation} \\ \hline
\textbf{Bière de mil / « bil-bil »} (Grand Nord, Adamaoua) & Mil (penicillaria), sorgho, maïs maltés & Levures (\emph{Saccharomyces cerevisiae}), bactéries lactiques & Alcoolique + lactique (acidification) \\ \hline
\textbf{Vin de palme / « matango »} (Sud, Centre, Littoral, Ouest) & Sève de raphia (\emph{Raphia hookeri}) ou palmier à huile & Levures (\emph{Saccharomyces}, \emph{Candida}), bactéries acétiques & Alcoolique $\rightarrow$ acétique (si aération) \\ \hline
\textbf{Kwatinkwang / Bâton de manioc} (Centre, Sud, Est, Littoral) & Racines de manioc (variétés amères) trempées & Bactéries lactiques (\emph{Lactobacillus} spp.), levures & Lactique (attendrissement, détoxication) \\ \hline
\textbf{Ngoundja / Feuille de manioc fermentée} (Ouest, Nord-Ouest) & Jeunes feuilles de manioc pilées & Bactéries lactiques, \emph{Bacillus} spp. & Lactique (conservation, goût) \\ \hline
\textbf{Lait caillé / « kindirmo »} (Nord, Adamaoua, Extrême-Nord) & Lait de vache cru & Bactéries lactiques mésophiles naturelles & Lactique \\ \hline
\end{tabularx}
\end{center}
\textbf{Note :} la fermentation du manioc (kwatinkwang, miondo, bobolo) élimine partiellement les glycosides cyanogènes (linamarine) toxiques, rendant le produit comestible. C'est un enjeu majeur de sécurité alimentaire.
\end{exoresolu}

\courssub{Biotechnologies modernes : agriculture et médecine}

\begin{aretenir}[Biotechnologies modernes majeures]
\begin{itemize}
\item \textbf{Agriculture (Biotechnologie verte) :}
  \begin{itemize}
  \item \cle{Culture de tissus \emph{in vitro} (microbouturage) :} multiplication rapide de plants sains, exempts de virus (manioc, bananier, plantain, ananas, pomme de terre) à partir de méristèmes. Utilisée à l'IRAD (Institut de Recherche Agricole pour le Développement) et dans des pépinières privées.
  \item \cle{Sélection assistée par marqueurs moléculaires (SAM) :} sélection précoce de caractères d'intérêt (résistance maladies, rendement) sans attendre la phénotype.
  \item \cle{Plantes transgéniques (OGM) :} introduction d'un gène étranger (ex. gène \emph{Bt} de \emph{Bacillus thuringiensis} pour résistance aux insectes ; gène \emph{CP4 EPSPS} pour tolérance au glyphosate). \textbf{Au Cameroun :} essais en champ confiné pour coton Bt et maïs Bt (IRAD), pas encore de culture commerciale autorisée à grande échelle (cadre loi n°2003/006 sur la biosécurité).
  \end{itemize}
\item \textbf{Médecine (Biotechnologie rouge) :}
  \begin{itemize}
  \item \cle{Protéines recombinantes :} insuline humaine, hormone de croissance, facteurs de coagulation, vaccins à sous-unités (Hépatite B, VPH) produits par bactéries (\emph{E. coli}) ou levures (\emph{Pichia pastoris}) ou cellules CHO (ovaire de hamster chinois).
  \item \cle{Vaccins à ARN messager :} technologie plateforme (COVID-19) : ARN codant l'antigène encapsulé dans nanoparticules lipidiques.
  \item \cle{Diagnostic moléculaire :} PCR (amplification d'ADN) pour détection du VIH (charge virale), paludisme (\emph{Plasmodium} spp.), tuberculose, HPV, COVID-19. Tests rapides antigéniques (immunochromatographie).
  \item \cle{Thérapie génique et cellulaire :} encore au stade recherche/essais cliniques (drépanocytose, certains cancers).
  \end{itemize}
\item \textbf{Industrie \& Environnement (Biotechnologie blanche / grise) :}
  \begin{itemize}
  \item Enzymes industrielles (amylases, protéases, cellulases) pour lessives, textile, papier, bioéthanol 2G (biomasse lignocellulosique).
  \item Bioremédiation : bactéries dégradant hydrocarbures (pollution pétrolière), métaux lourds.
  \item Bioplastiques (PHA, PLA) produits par fermentation bactérienne.
  \end{itemize}
\end{itemize}
\end{aretenir}

\begin{exoresolu}[Biotechnologies modernes au Cameroun]
Énoncé. Cite deux exemples de biotechnologies modernes appliquées ou en développement au Cameroun, en précisant l'organisme/technique utilisé et le produit/service visé.
\tcblower
\textbf{Solution détaillée.}
\begin{enumerate}
\item \textbf{Microbouturage \emph{in vitro} de plants de bananier/plantain et manioc (IRAD, centres de multiplication) :} technique de culture de tissus sur milieu nutritif stérile à partir de méristèmes apicaux. \textbf{Produit :} plants sains, exempts de virus (BBTV, CMV) et de nématodes, à haut potentiel de rendement. \textbf{Enjeu :} relance des filières banane/plantain et manioc, sécurité alimentaire.
\item \textbf{Production et utilisation de vaccins recombinants et tests PCR :} dans le Programme Élargi de Vaccination (PEV) et les laboratoires de santé publique (Centre Pasteur du Cameroun, labouratoires hospitaliers). \textbf{Exemples :} vaccin anti-hépatite B (protéine HBsAg recombinante levure), vaccin anti-HPV, tests PCR charge virale VIH, diagnostic paludisme par PCR. \textbf{Enjeu :} lutte contre les maladies endémiques, santé publique.
\end{enumerate}
\end{exoresolu}

\courssub{Avantages, limites et débats}

\begin{methode}[Rédiger une discussion argumentée « avantages / limites »]
\begin{enumerate}
\item \cle{Rappeler en une phrase} le principe de la biotechnologie concernée.
\item Dresser les \cle{avantages} classés : économiques (rendement, revenu, coût), sociaux (santé, alimentation, emploi), environnementaux (moins d'intrants chimiques).
\item Dresser les \cle{limites et risques} : techniques (efficacité partielle), sanitaires (allergénicité, toxicité long terme), écologiques (dissémination gènes, résistance ravageurs, effet non-cible), éthiques (brevetage du vivant, consentement), économiques (dépendance semences, coût intrants).
\item \cle{Nuancer} : un avantage peut avoir une contrepartie (ex. rendement élevé mais dépendance aux semences brevetées).
\item Conclure par un \cle{avis argumenté} ou conditions d'utilisation responsable (réglementation biosécurité, évaluation risques, étiquetage, formation, contrôle qualité).
\end{enumerate}
\textbf{Règle d'or Probatoire :} une discussion se note sur l'\emph{argumentation}. Chaque point doit être justifié, pas simplement énuméré. Terminer par une conclusion nuancée.
\end{methode}

\begin{exoresolu}[Les organismes génétiquement modifiés (OGM) en agriculture]
Énoncé. Certains agriculteurs souhaitent introduire au Cameroun des semences de maïs génétiquement modifié résistant aux insectes ravageurs (foreurs). Présente deux avantages et deux limites de cette biotechnologie, puis donne ton avis argumenté.
\tcblower
\textbf{Solution détaillée.}
\textbf{Principe :} le maïs OGM a reçu, par génie génétique (transgénèse), un gène \emph{cry} issu de la bactérie \emph{Bacillus thuringiensis} (gène Bt) codant une protéine insecticide (Cry1Ab) active sur les lépidoptères foreurs (\emph{Busseola fusca}, \emph{Chilo partellus}).
\textbf{Avantages :}
\begin{enumerate}
\item \textbf{Rendements améliorés et sécurité alimentaire :} la plante exprime sa propre toxine, protégeant l'épi et la tige. Moins de pertes à la récolte (jusqu'à 30--40\% gagnés en zone de forte pression), production plus stable.
\item \textbf{Réduction des insecticides chimiques :} moins de pulvérisations $\rightarrow$ économie pour l'agriculteur, moins de pollution sols/eaux, moins d'intoxications aiguës/chroniques des applicateurs, préservation de la faune auxiliaire (si sélectivité respectée).
\end{enumerate}
\textbf{Limites :}
\begin{enumerate}
\item \textbf{Risques environnementaux :} (a) Dissémination du gène Bt par pollinisation vers maïs conventionnel ou parents sauvages (teosinte absent au Cameroun, mais flux de gènes possible) ; (b) Apparition d'insectes résistants (sélection de mutants) $\rightarrow$ nécessité de \cle{refuges} (zones non-Bt) pour gérer la résistance ; (c) Effets possibles sur insectes non-cibles (papillons, abeilles) bien que la protéine Cry soit généralement spécifique.
\item \textbf{Dépendance économique et questions sanitaires/sociétales :} (a) Semences brevetées : interdiction de réutiliser la récolte comme semence (droit de propriété intellectuelle), achat annuel obligatoire $\rightarrow$ dépendance vis-à-vis de semenciers multinationaux ; (b) Innocuité long terme : évaluée cas par cas (substantial equivalence), mais surveillance post-commercialisation nécessaire ; (c) Acceptabilité sociétale : demande d'étiquetage, principe de précaution, souveraineté alimentaire.
\end{enumerate}
\textbf{Avis argumenté :} l'introduction du maïs Bt peut être bénéfique pour la productivité et la réduction des pesticides au Cameroun, \emph{à condition} d'un cadre strict : (1) évaluation environnementale et sanitaire préalable indépendante (CNAB - Comité National de Biosécurité), (2) plan de gestion de la résistance (refuges obligatoires, rotation cultures), (3) information/formation des producteurs, (4) étiquetage pour le choix du consommateur, (5) préservation de variétés locales non-OGM (banques de gènes).
\textbf{Conclusion :} comme toute biotechnologie, l'OGM présente un rapport avantages/risques qu'il faut évaluer scientifiquement et encadrer juridiquement avant diffusion.
\end{exoresolu}

\courssub{Enjeux économiques et sociaux au Cameroun}

\begin{methode}[Analyser les enjeux économiques et sociaux d'une biotechnologie]
\begin{enumerate}
\item \cle{Enjeu économique :} création de revenus/emplois (filières transformation artisanale/industrielle), réduction importations (blé, lait en poudre, médicaments), valorisation matières premières locales (maïs, sorgho, manioc, lait, palme), attractivité investissements (biotechpharma).
\item \cle{Enjeu social :} sécurité alimentaire (conservation par fermentation sans chaîne de froid : yaourt, manioc, lait caillé), santé publique (vaccins, diagnostics rapides, médicaments biotech), autonomisation (micro-entreprises femmes/jeunes : yaourts, jus, vin de palme, plants \emph{in vitro}).
\item \cle{Défis à lever :} équipements/laboratoires vétustes, formation insuffisante (techniciens, ingénieurs), financement recherche (budget \textless 0,5\% PIB), contrôle qualité/hygiène produits artisanaux (risques toxi-infections), cadre juridique biosécurité (loi 2003, décrets d'application), propriété intellectuelle.
\item \cle{Conditions du développement :} recherche publique (IRAD, Universités, Centre Pasteur) + secteur privé (brasseries, laiteries, semenciers, start-ups biotech) + État (politiques incitatives, infrastructures, formation).
\end{enumerate}
\end{methode}

\begin{exoresolu}[Enjeux de la biotechnologie au Cameroun]
Énoncé. En t'appuyant sur deux exemples, montre que le développement des biotechnologies représente un enjeu économique et social majeur pour le Cameroun.
\tcblower
\textbf{Solution détaillée.}
\textbf{Exemple 1 -- Filières brassicoles et fermentations alimentaires (traditionnelles/industrielles).} La production de bière (industrielle : maïs/sorgho local chez SABC/UCB ; artisanale : mil « bil-bil ») et d'aliments fermentés (kwatinkwang, yaourt, kindirmo) mobilise des milliers d'acteurs (producteurs, maltateurs, brasseurs, transformatrices, vendeuses). \textbf{Économique :} emplois directs/indirects, revenus ruraux, substitution importations (malt d'orge, lait en poudre). \textbf{Social :} la fermentation conserve les aliments sans réfrigération (acide lactique, éthanol), sécurise l'alimentation en période de soudure, améliore la digestibilité et la valeur nutritionnelle (vitamines B, biodisponibilité minéraux).
\textbf{Exemple 2 -- Biotechnologies médicales (modernes).} Vaccins (PEV : rougeole, fièvre jaune, HPV, hépatite B, COVID-19), tests rapides (paludisme, VIH, COVID), traitements (ART VIH, insuline). \textbf{Social :} réduction mortalité infanto-juvénile/maternelle, contrôle épidémies, espérance de vie. \textbf{Économique :} population en bonne santé = main-d'œuvre productive ; production locale future de vaccins/biomédicaments (projet usine vaccins) réduirait la facture pharmaceutique (centaines de milliards FCFA/an) et assurerait la souveraineté sanitaire.
\textbf{Conclusion :} des fermentations villageoises aux biomédicaments, les biotechnologies sont un levier de développement économique (filières, emplois, import-substitution) et social (alimentation, santé) pour le Cameroun. Leur essor requiert investissement dans la R\&D, la formation technique/universitaire, le contrôle qualité et un cadre réglementaire incitatif (biosécurité, propriété intellectuelle, partenariats public-privé).
\end{exoresolu}

\courssub{Compétence transversale : Exploiter des données expérimentales sur une fermentation}

\begin{methode}[Analyser un tableau ou une courbe de fermentation]
\begin{enumerate}
\item \cle{Lire les grandeurs} : variables mesurées (temps, température, sucre résiduel, gaz \ce{CO2}, éthanol, pH, biomasse...).
\item \cle{Décrire les variations} : quelle grandeur augmente/diminue, cinétique (rapide/lent), plateau/arrêt.
\item \cle{Relier les variations entre elles} : disparition substrat $\leftrightarrow$ apparition produits (stœchiométrie).
\item \cle{Interpréter biologiquement} : phase de latence, exponentielle, stationnaire, mort ; rôle du micro-organisme ; facteurs limitants (substrat, pH, toxicité produits, température).
\item \cle{Conclure} en reliant au procédé (artisanal/industriel) et aux paramètres de pilotage.
\end{enumerate}
\end{methode}

\begin{exoresolu}[Suivi d'une fermentation alcoolique]
Énoncé. On place du jus de raisin sucré avec des levures (\emph{Saccharomyces cerevisiae}) dans un récipient muni d'un dispositif recueillant le gaz dégagé. Résultats :

\begin{center}
\begin{tabularx}{\linewidth}{|l|X|X|X|X|X|}\hline
\rowcolor{popblueL} Temps (h) & 0 & 6 & 12 & 24 & 48 \\ \hline
Volume de \ce{CO2} dégagé (mL) & 0 & 15 & 40 & 75 & 80 \\ \hline
Masse de glucose restant (g) & 100 & 88 & 62 & 20 & 15 \\ \hline
\end{tabularx}
\end{center}
\begin{enumerate}
\item Décris l'évolution des deux grandeurs mesurées.
\item Interprète ces résultats en les reliant à l'équation-bilan.
\item Pourquoi le dégagement de gaz ralentit-il fortement après 24 h ?
\end{enumerate}
\tcblower
\textbf{Solution détaillée.}
\begin{enumerate}
\item \textbf{Description.} Le volume de \ce{CO2} dégagé augmente avec le temps : 0 mL (0 h), 75 mL (24 h), puis seulement 80 mL (48 h) — net ralentissement après 24 h. Simultanément, la masse de glucose diminue : 100 g (0 h), 20 g (24 h), 15 g (48 h).
\item \textbf{Interprétation.} Les levures réalisent la \textbf{fermentation alcoolique} en anaérobiose : $\ce{C6H12O6 -> 2 C2H5OH + 2 CO2}$ + énergie. La disparition du glucose (80 g consommés en 24 h) s'accompagne logiquement du dégagement de \ce{CO2} (75 mL) et de la production d'éthanol (non mesuré). La concomitance confirme l'origine fermentaire du gaz.
\item \textbf{Ralentissement après 24 h :} deux facteurs limitants combinés : (1) \textbf{Épuisement du substrat} : glucose résiduel faible (20 g $\rightarrow$ 15 g), vitesse de réaction limitée par la concentration en substrat (cinétique de Michaelis-Menten) ; (2) \textbf{Toxicité de l'éthanol accumulé} : concentration éthylique devenue inhibitrice pour les levures (dénaturation protéines, perturbation membranes), entraînant arrêt de l'activité fermentaire. C'est pourquoi les vins naturels titrent rarement au-delà de 12--15\% vol.
\end{enumerate}
\textbf{Conclusion :} le dégagement de \ce{CO2} est la signature cinétique de l'activité fermentaire ; la fermentation s'auto-limite par épuisement du sucre et/ou toxicité de l'alcool produit.
\end{exoresolu}

\begin{attention}[Les 5 erreurs qui coûtent des points au Probatoire]
\begin{enumerate}
\item \textbf{Confondre fermentation et respiration cellulaire.} Fermentation = \emph{anaérobie} (sans \ce{O2}), rendement énergétique faible (2 ATP/glucose). Respiration = \emph{aérobie} (avec \ce{O2}), rendement élevé ($\sim$30--32 ATP/glucose). Écrire « les levures respirent le glucose pour faire du vin » est \textbf{faux} : elles le \emph{fermentent}.
\item \textbf{Équation-bilan fausse ou non équilibrée.} Fermentation alcoolique : $\ce{C6H12O6 -> 2 C2H5OH + 2 CO2}$ (éthanol + gaz carbonique). Pas d'\ce{O2} produit, pas de vinaigre (acide acétique = fermentation acétique, aérobie). Fermentation lactique : $\ce{C6H12O6 -> 2 C3H6O3}$ (acide lactique \textbf{sans gaz}).
\item \textbf{Oublier de nommer le micro-organisme.} Une description sans agent biologique (levures pour alcool/pain, bactéries lactiques pour yaourt/manioc, \emph{Acetobacter} pour vinaigre) est incomplète : l'organisme vivant \emph{est} la biotechnologie.
\item \textbf{Donner des exemples non biotechnologiques.} Séchage, salage, fumage, cuisson, engrais chimiques, pesticides de synthèse, labour mécanique ne sont \emph{pas} des biotechnologies. Vérifier : \emph{organisme vivant + procédé + produit utile}.
\item \textbf{Discussion « avantages/limites » en liste brute sans justification.} Chaque point doit être \emph{argumenté} (mécanisme, conséquence, chiffre si possible) et la copie doit se terminer par une \textbf{conclusion nuancée} (conditions d'usage responsable). Une simple énumération perd la moitié des points d'argumentation.
\end{enumerate}
\end{attention}

\begin{savaistu}[Biotechnologies et patrimoine camerounais]
Le Cameroun possède une diversité exceptionnelle de fermentations traditionnelles : plus de 50 produits fermentés répertoriés (boissons : bil-bil, amgba, foléré, vin de palme, nkang ; aliments : kwatinkwang, miondo, bobolo, ngoundja, kindirmo, koki, soja fermenté dawadawa). Ces savoir-faire, transmis oralement, constituent un patrimoine immatériel et un réservoir de souches microbiennes d'intérêt biotechnologique (probiotiques, enzymes, métabolites). L'IRAD et les universités (Yaoundé I, Ngaoundéré, Dschang, Buea) mènent des programmes de \textbf{caractérisation moléculaire} de ces souches (levures \emph{Saccharomyces}, \emph{Candida}, \emph{Kluyveromyces} ; bactéries \emph{Lactobacillus}, \emph{Weissella}, \emph{Leuconostoc}, \emph{Bacillus}) pour les sélectionner comme \textbf{ferments starters} standardisés, améliorant l'hygiène, la régularité et la qualité nutritionnelle des produits artisanaux. C'est un exemple concret de \textbf{valorisation moderne des biotechnologies traditionnelles}.
\end{savaistu}

\newpage

\coursec{Exercices}

\courssub{Je vérifie mes connaissances}

\begin{exercice}[QCM : microorganismes et fermentations]
Pour chaque question, une seule réponse est exacte. Entoure la bonne lettre.

\textbf{1.} La fabrication de la bière de maïs (\og bili-bili \fg{}) met en jeu principalement :
\begin{itemize}
\item[a)] des bactéries lactiques et une fermentation lactique ;
\item[b)] des levures (\emph{Saccharomyces}) et une fermentation alcoolique ;
\item[c)] des moisissures et une fermentation acétique ;
\item[d)] des virus et une fermentation butyrique.
\end{itemize}

\textbf{2.} Le yaourt est obtenu grâce à :
\begin{itemize}
\item[a)] des levures qui transforment le lait en alcool ;
\item[b)] des bactéries lactiques (\emph{Lactobacillus}, \emph{Streptococcus thermophilus}) qui transforment le lactose en acide lactique ;
\item[c)] des champignons filamenteux qui coagulent le lait ;
\item[d)] une simple ébullition prolongée du lait.
\end{itemize}

\textbf{3.} Dans la panification, le gaz qui fait lever la pâte est :
\begin{itemize}
\item[a)] le dioxygène ; \item[b)] le diazote ; \item[c)] le dioxyde de carbone ; \item[d)] le méthane.
\end{itemize}

\textbf{4.} Le vin de palme résulte de la fermentation :
\begin{itemize}
\item[a)] lactique de la sève de palmier ;
\item[b)] alcoolique des sucres de la sève de palmier par des levures ;
\item[c)] acétique du glucose par des bactéries ;
\item[d)] butyrique de l'amidon.
\end{itemize}
\end{exercice}

\begin{exercice}[Vrai ou faux justifié]
Pour chaque affirmation, indique si elle est vraie ou fausse, puis justifie en une ou deux phrases.
\begin{enumerate}
\item La biotechnologie est une science née au XXI\up{e} siècle.
\item La fermentation alcoolique se déroule en présence abondante de dioxygène.
\item La production d'insuline par des bactéries génétiquement modifiées est un exemple de biotechnologie moderne.
\item Le \og miondo \fg{} et le \og bobolo \fg{} (bâtons de manioc fermenté) relèvent d'une biotechnologie traditionnelle camerounaise.
\end{enumerate}
\end{exercice}

\begin{exercice}[Définitions]
Définis en deux ou trois phrases chacun des termes suivants :
\begin{enumerate}
\item Biotechnologie ;
\item Fermentation ;
\item Organisme génétiquement modifié (OGM).
\end{enumerate}
\end{exercice}

\begin{exercice}[Équation bilan à compléter]
Recopie et complète l'équation bilan de la fermentation alcoolique du glucose par la levure :
\[ \ce{C6H12O6 -> ... C2H5OH + ... CO2 + energie} \]
Cite ensuite deux boissons produites au Cameroun qui résultent de cette fermentation.
\end{exercice}

\courssub{Je m'entraîne}

\begin{exercice}[Question de cours structurée]
\begin{enumerate}
\item Définis la biotechnologie.
\item À l'aide de deux exemples chacune, distingue biotechnologie traditionnelle et biotechnologie moderne.
\item Cite trois domaines d'application de la biotechnologie dans la vie quotidienne.
\end{enumerate}
\end{exercice}

\begin{exercice}[Étude de document : l'insuline \og fabriquée \fg{} par des bactéries]
\textbf{Document.} Avant 1982, l'insuline destinée aux diabétiques était extraite du pancréas de porcs ou de bœufs abattus : il fallait environ \SI{8000}{kg} de pancréas pour obtenir \SI{1}{kg} d'insuline, et certains patients développaient des allergies. Depuis 1982, on insère le gène humain de l'insuline dans le patrimoine génétique de la bactérie \emph{Escherichia coli}. Cultivées en grand nombre dans des cuves appelées fermenteurs, ces bactéries produisent de l'insuline humaine, que l'on extrait et purifie ensuite.

\begin{enumerate}
\item Identifie, dans ce procédé, le microorganisme utilisé et la molécule d'intérêt produite.
\item Décris en trois étapes le principe de ce procédé biotechnologique.
\item Cite un avantage de cette biotechnologie par rapport à l'ancienne méthode d'extraction.
\item Cite une limite ou un risque associé à l'utilisation d'organismes génétiquement modifiés.
\end{enumerate}
\end{exercice}

\begin{exercice}[Calcul : rendement théorique d'une fermentation]
Une brasserie de Douala fait fermenter $m = \SI{1000}{kg}$ (1 tonne) de glucose selon la réaction :
\[ \ce{C6H12O6 -> 2 C2H5OH + 2 CO2} \]
On donne les masses molaires : $M(\ce{C}) = \SI{12}{g.mol^{-1}}$ ; $M(\ce{H}) = \SI{1}{g.mol^{-1}}$ ; $M(\ce{O}) = \SI{16}{g.mol^{-1}}$.
\begin{enumerate}
\item Calcule les masses molaires du glucose et de l'éthanol.
\item Détermine la quantité de matière (en mol) de glucose mise en fermentation.
\item En utilisant les coefficients stœchiométriques, calcule la masse théorique d'éthanol produite.
\item En réalité, la brasserie ne récolte que \SI{400}{kg} d'éthanol. Calcule le rendement de la fermentation en pourcentage et propose une explication à cet écart.
\end{enumerate}
\end{exercice}

\begin{exercice}[Avantages et limites d'une biotechnologie locale]
Au Cameroun, la fermentation du manioc permet de produire le \og bâton de manioc \fg{} (bobolo, miondo), aliment de base de nombreuses familles.
\begin{enumerate}
\item Explique pourquoi la fermentation du manioc est indispensable avant sa consommation (indice : le manioc amer contient des composés cyanés toxiques).
\item Présente deux avantages de cette biotechnologie traditionnelle pour les populations.
\item Présente une limite de ce procédé artisanal et propose une amélioration possible.
\end{enumerate}
\end{exercice}

\courssub{Je me prépare au Probatoire}

\begin{exercice}[Sujet type Probatoire : la brasserie et le vin de palme]
\textbf{Situation-problème.} Dans la région du Littoral, une coopérative souhaite valoriser la production locale de maïs en fabriquant de la bière artisanale. Un membre de la coopérative affirme : \og La fabrication de la bière et celle du vin de palme reposent sur le même phénomène biologique. \fg{}

\begin{enumerate}
\item Définis la biotechnologie et montre que la fabrication de la bière en relève. \hfill (2 pts)
\item Nomme le microorganisme responsable de la fabrication de la bière et précise le type de fermentation mis en jeu. \hfill (2 pts)
\item Écris l'équation bilan de cette fermentation. \hfill (2 pts)
\item En comparant la bière et le vin de palme, dis si l'affirmation du membre de la coopérative est correcte ; justifie ta réponse. \hfill (2 pts)
\item La coopérative fait fermenter \SI{360}{kg} de glucose. Sachant que $M(\text{glucose}) = \SI{180}{g.mol^{-1}}$ et $M(\text{éthanol}) = \SI{46}{g.mol^{-1}}$, calcule la masse théorique d'éthanol obtenue. \hfill (3 pts)
\item Dégage deux retombées économiques de ce projet pour la localité et une précaution d'hygiène à respecter pendant la fabrication. \hfill (3 pts)
\end{enumerate}
\end{exercice}

\begin{exercice}[Sujet type Probatoire : biotechnologies médicales et agricoles]
\textbf{Document 1.} Au Centre Pasteur du Cameroun, des vaccins et des tests de diagnostic sont produits ou conditionnés grâce à des techniques issues de la biologie moléculaire.

\textbf{Document 2.} Des variétés améliorées de maïs et de manioc, plus résistantes aux maladies et à la sécheresse, sont diffusées auprès des agriculteurs camerounais par l'IRAD (Institut de Recherche Agricole pour le Développement).

\begin{enumerate}
\item Définis la biotechnologie moderne et distingue-la de la biotechnologie traditionnelle. \hfill (3 pts)
\item Pour chacun des deux documents, précise le domaine d'application (médical ou agricole) et l'intérêt pour les populations camerounaises. \hfill (4 pts)
\item Explique le principe général de l'obtention d'un organisme génétiquement modifié (OGM). \hfill (3 pts)
\item Certains consommateurs se méfient des OGM. Présente un argument en faveur de leur utilisation et un argument contre, puis donne ton avis argumenté. \hfill (4 pts)
\end{enumerate}
\end{exercice}

\begin{exercice}[Essai argumenté type épreuve]
\textbf{Sujet.} \og Les biotechnologies sont-elles une chance pour le développement du Cameroun ? \fg{}

À l'aide de tes connaissances et d'exemples camerounais (brasseries, transformation du manioc, production de médicaments, semences améliorées, etc.), rédige un développement structuré d'environ vingt-cinq lignes comportant :
\begin{enumerate}
\item une introduction définissant la biotechnologie et posant la problématique ;
\item au moins deux arguments économiques en faveur des biotechnologies ;
\item au moins un argument social ou sanitaire ;
\item au moins une limite ou un risque (environnemental, éthique ou économique) ;
\item une conclusion dans laquelle tu prends clairement position.
\end{enumerate}
\hfill (Barème indicatif : introduction 2 pts ; arguments 8 pts ; limite 2 pts ; conclusion 2 pts ; cohérence et qualité de la langue 2 pts.)
\end{exercice}

\newpage

\coursec{Corrigés des exercices}

\coursec{Corrigés des exercices d'application — Leçon NA08}

\begin{corrige}[Exercice 1 — QCM : microorganismes et fermentations]
\textbf{1. b)} La fabrication de la bière de maïs (bili-bili ou kwatcha) repose sur la fermentation alcoolique réalisée par des levures du genre \emph{Saccharomyces}, qui transforment les sucres (glucose, maltose) en éthanol et dioxyde de carbone.
\textbf{2. b)} Le yaourt résulte de la fermentation lactique : des bactéries lactiques (\emph{Lactobacillus bulgaricus} et \emph{Streptococcus thermophilus}) transforment le lactose (sucre du lait) en acide lactique, ce qui acidifie le milieu, coagule les protéines (caséines) et donne sa texture et son goût caractéristique.
\textbf{3. c)} En panification, les levures (\emph{Saccharomyces cerevisiae}) effectuent une fermentation alcoolique ; le gaz produit est le dioxyde de carbone ($\ce{CO2}$), qui forme des bulles dans le réseau de gluten et fait lever la pâte.
\textbf{4. b)} Le vin de palme est obtenu par fermentation alcoolique spontanée des sucres (saccharose, glucose, fructose) présents dans la sève de palmier (rafia ou palmier à huile), sous l'action de levures naturelles.
\end{corrige}

\begin{corrige}[Exercice 2 — Vrai ou faux justifié]
\begin{enumerate}
\item \textbf{Faux.} La biotechnologie \emph{traditionnelle} (fermentation pour le pain, le vin, la bière, le fromage) est pratiquée depuis des millénaires (Égypte antique, Mésopotamie). Seule la biotechnologie \emph{moderne} (génie génétique, culture cellulaire, PCR) est récente (développement majeur à partir des années 1970-1980).
\item \textbf{Faux.} La fermentation alcoolique est un métabolisme \emph{anaérobie} strict : elle se déroule \emph{en l'absence de dioxygène}. En présence d'$\ce{O2}$, les levures privilégient la respiration cellulaire (beaucoup plus efficace énergétiquement).
\item \textbf{Vrai.} L'insuline humaine recombinante (Humulin\textsuperscript{\textregistered}, 1982) est le premier médicament issu du génie génétique : le gène humain de l'insuline est inséré dans \emph{E. coli} (ou levure), qui la produit en grande quantité dans des fermenteurs. C'est un exemple emblématique de biotechnologie moderne (biotechnologie médicale / « biotechnologie rouge »).
\item \textbf{Vrai.} Le \og miondo \fg{} (bâton de manioc fermenté mou, région du Centre/Sud/Est) et le \og bobolo \fg{} (bâton de manioc fermenté plus ferme, région du Littoral/Sud-Ouest) sont obtenus par fermentation naturelle (trempage, détoxication, fermentation lactique et levurienne) de racines de manioc. C'est un patrimoine biotechnologique traditionnel camerounais majeur.
\end{enumerate}
\end{corrige}

\begin{corrige}[Exercice 3 — Définitions]
\begin{enumerate}
\item \textbf{Biotechnologie :} Ensemble des techniques qui utilisent des organismes vivants (microorganismes, cellules animales ou végétales, enzymes) ou des parties d'organismes (ADN, organites) pour produire ou modifier des produits, améliorer des plantes ou des animaux, ou développer des micro-organismes à des usages spécifiques (définition OMS/Convention sur la diversité biologique). On distingue biotechnologies traditionnelles (fermentations empiriques) et modernes (génie génétique).
\item \textbf{Fermentation :} Processus métabolique anaérobie (sans dioxygène) au cours duquel un microorganisme (levure, bactérie) dégrade un substrat organique (glucose, lactose, etc.) en produits finis plus simples (éthanol, acide lactique, $\ce{CO2}$, etc.) en récupérant de l'énergie (ATP) par phosphorylation au niveau du substrat. Exemples : fermentation alcoolique, fermentation lactique.
\item \textbf{Organisme génétiquement modifié (OGM) :} Organisme (bactérie, plante, animal) dont le matériel génétique (ADN) a été modifié par une technique de génie génétique (transgenèse, édition de génome type CRISPR-Cas9) d'une manière qui ne s'obtient pas naturellement par multiplication et/ou recombinaison génétique naturelle. L'organisme exprime alors un ou plusieurs gènes étrangers (transgènes) ou modifiés.
\end{enumerate}
\end{corrige}

\begin{corrige}[Exercice 4 — Équation bilan à compléter]
L'équation bilan de la fermentation alcoolique du glucose par la levure (\emph{Saccharomyces cerevisiae}) est :
\[ \ce{C6H12O6 -> 2 C2H5OH + 2 CO2 + energie (ATP)} \]
\textbf{Deux boissons camerounaises issues de cette fermentation :}
\begin{itemize}
\item La bière de maïs (\og bili-bili \fg{} ou \og kwatcha \fg{}) ;
\item Le vin de palme (sève de raphia ou de palmier à huile fermentée).
\end{itemize}
\emph{Autres exemples :} bière de mil/sorgho (\og bili-bili \fg{} au Nord), vin de banane, \og mbuh \fg{} (boisson fermentée à base de maïs germé).
\end{corrige}

\begin{corrige}[Exercice 5 — Question de cours structurée]
\begin{enumerate}
\item \textbf{Définition :} Voir Exercice 3, réponse 1.
\item \textbf{Distinction :}
\begin{center}
\begin{tabularx}{\linewidth}{|l|X|X|}\hline
 & \textbf{Biotechnologie traditionnelle} & \textbf{Biotechnologie moderne} \\ \hline
\textbf{Principe} & Utilisation empirique d'organismes entiers (levures, bactéries) sans modification génétique directe. & Manipulation directe du génome (ADN recombinant, édition génomique, culture de cellules/tissus \emph{in vitro}). \\ \hline
\textbf{Exemple 1} & Fabrication du pain, de la bière, du vin, du yaourt, du fromage, du miondo/bobolo. & Production d'insuline humaine par \emph{E. coli} transgénique ; vaccins à ARN messager. \\ \hline
\textbf{Exemple 2} & Compostage, ensilage, épuration biologique des eaux usées (boues activées). & Maïs Bt résistant aux insectes (gène \emph{cry} de \emph{Bacillus thuringiensis}) ; riz doré (provitamine A) ; thérapie génique. \\ \hline
\end{tabularx}
\end{center}
\item \textbf{Trois domaines d'application dans la vie quotidienne :}
\begin{enumerate}
\item \textbf{Alimentaire :} Production de boissons (bière, vin), produits laitiers (yaourt, fromage), pain, condiments, probiotiques.
\item \textbf{Médical / Santé :} Production de médicaments (insuline, hormones de croissance, vaccins, anticorps monoclonaux), diagnostic (tests PCR, tests rapides VIH/paludisme), thérapie génique.
\item \textbf{Agricole / Environnement :} Semences améliorées (OGM ou sélection assistée par marqueurs), biofertilisants (bactéries fixatrices d'azote \emph{Rhizobium}), biopesticides (\emph{Bacillus thuringiensis}), dépollution (biorémédiation sols/eaux pollués par hydrocarbures).
\end{enumerate}
\end{enumerate}
\end{corrige}

\begin{corrige}[Exercice 6 — Étude de document : l'insuline « fabriquée » par des bactéries]
\begin{enumerate}
\item \textbf{Microorganisme :} La bactérie \emph{Escherichia coli} (souche modifiée, non pathogène, ex. K12). \\
\textbf{Molécule d'intérêt :} L'insuline humaine (protéine hormonale de 51 acides aminés, deux chaînes A et B liées par ponts disulfure).
\item \textbf{Principales étapes du procédé (biotechnologie moderne / ADN recombinant) :}
\begin{enumerate}
\item \textbf{Construction du vecteur recombinant :} Isolement du gène humain codant l'insuline (ou synthèse chimique) et insertion dans un plasmide (vecteur d'expression) contenant un promoteur fort et un gène de résistance (ex. ampicilline). Introduction de ce plasmide dans \emph{E. coli} compétentes (transformation).
\item \textbf{Culture et fermentation (amplification) :} Culture des bactéries transformées en grand volume dans des fermenteurs (cuves stériles, contrôlées : température \SI{37}{\celsius}, pH, agitation, apport en nutriments et $\ce{O2}$ pour la croissance, puis conditions d'induction de l'expression du gène). Les bactéries se multiplient et expriment le gène $\rightarrow$ production massive d'insuline (souvent sous forme de « corps d'inclusion » ou sécrétée).
\item \textbf{Extraction et purification (aval) :} Récolte de la biomasse (centrifugation), lyse bactérienne, séparation de l'insuline des protéines bactériennes (chromatographies successives : échange d'ions, exclusion stérique, phase inverse), repliage oxydatif (formation des ponts disulfure si chaînes séparées), formulation finale (contrôle qualité : pureté, stérilité, activité biologique).
\end{enumerate}
\item \textbf{Avantage majeur :} L'insuline produite est \textbf{strictement identique à l'insuline humaine} (séquence d'acides aminés), ce qui élimine les réactions allergiques et immunologiques observées avec l'insuline porcine/bovine (différences de 1 acide aminé pour le porc, 3 pour le bœuf). Autre avantage : production illimitée, indépendante de l'abattage animal, coût réduit à grande échelle.
\item \textbf{Limite / Risque associé aux OGM (ici bactéries GM) :}
\begin{itemize}
\item \textbf{Risque de dissémination :} Fuite accidentelle de bactéries transgéniques dans l'environnement (eaux usées, air) avec transfert horizontal possible du gène de résistance aux antibiotiques (marqueur de sélection) vers des bactéries pathogènes.
\item \textbf{Contamination du produit :} Présence résiduelle d'ADN recombinant, d'endotoxines (LPS de \emph{E. coli}), ou de protéines hôtes dans le produit final $\rightarrow$ exigence de purification drastique.
\item \textbf{Aspect éthique / sociétal :} Brevetage du vivant, dépendance des pays du Sud vis-à-vis des firmes biotechnologiques du Nord, acceptabilité sociale des OGM.
\end{itemize}
\end{enumerate}
\end{corrige}

\begin{corrige}[Exercice 7 — Calcul : rendement théorique d'une fermentation]
\textbf{Données :} $m_{\text{glucose}} = \SI{1000}{kg} = 10^6 \text{ g}$.
$M(\ce{C}) = 12$, $M(\ce{H}) = 1$, $M(\ce{O}) = 16 \text{ g.mol}^{-1}$.
Réaction : $\ce{C6H12O6 -> 2 C2H5OH + 2 CO2}$.

\begin{enumerate}
\item \textbf{Masses molaires :}
\begin{align*}
M(\text{glucose } \ce{C6H12O6}) &= 6\times12 + 12\times1 + 6\times16 \\
&= 72 + 12 + 96 = \textbf{\SI{180}{g.mol^{-1}}} \\
M(\text{éthanol } \ce{C2H5OH} \equiv \ce{C2H6O}) &= 2\times12 + 6\times1 + 16 \\
&= 24 + 6 + 16 = \textbf{\SI{46}{g.mol^{-1}}}
\end{align*}

\item \textbf{Quantité de matière de glucose :}
\[ n_{\text{glucose}} = \frac{m}{M} = \frac{10^6 \text{ g}}{180 \text{ g.mol}^{-1}} = \frac{10000}{{1,8}} = \textbf{5555,56 mol} \quad (\approx \SI{5,56e3}{mol}) \]

\item \textbf{Masse théorique d'éthanol :}
D'après l'équation : $1 \text{ mol glucose} \rightarrow 2 \text{ mol éthanol}$.
\[ n_{\text{éthanol, théorique}} = 2 \times n_{\text{glucose}} = 2 \times 5555,56 = \SI{11111,11}{mol} \]
\[ m_{\text{éthanol, théorique}} = n \times M = 11111,11 \times 46 = \SI{511111,11}{g} = \textbf{\SI{511,11}{kg}} \]

\item \textbf{Rendement réel :}
Masse réelle récoltée $m_{\text{réel}} = \SI{400}{kg}$.
\[ \text{Rendement } \eta = \frac{m_{\text{réel}}}{m_{\text{théorique}}} \times 100 = \frac{400}{511,11} \times 100 \approx \textbf{78,3 \%} \]

\textbf{Explications de l'écart (rendement < 100 \%) :}
\begin{itemize}
\item Une partie du glucose est utilisée pour la \textbf{biomasse} (croissance et multiplication des levures) et non pour la production d'éthanol.
\item \textbf{Produits secondaires :} Formation de glycérol, d'acide acétique, d'acide succinique, d'esters, de composés aromatiques (métabolites secondaires).
\item \textbf{Perte par évaporation :} L'éthanol est volatil (point d'ébullition \SI{78,4}{\celsius}) ; pertes lors de la fermentation (non étanche) et surtout lors de la distillation.
\item \textbf{Inhibition :} L'éthanol inhibe l'activité enzymatique des levures à haute concentration (généralement > 12-15 \% vol), arrêtant la fermentation avant épuisement total du sucre.
\item \textbf{Conditions non optimales :} Température, pH, pression osmotique, carence en nutriments (azote, phosphates, vitamines).
\end{itemize}
\end{enumerate}
\end{corrige}

\begin{corrige}[Exercice 8 — Avantages et limites d'une biotechnologie locale : le bâton de manioc]
\begin{enumerate}
\item \textbf{Indispensabilité de la fermentation (détoxication) :} Le manioc amer (\emph{Manihot esculenta} Crantz, variétés à haut rendement) contient des \textbf{glucosides cyanogènes} (principalement la \textbf{linamarine} et la lotaustraline). Lors de la mastication ou du broyage, l'enzyme endogène \textbf{linamarase} (β-glucosidase) hydrolyse ces glucosides en libérant de l'\textbf{acide cyanhydrique (HCN)}, un poison puissant bloquant la cytochrome c oxydase (arrêt respiration cellulaire). La fermentation (trempage prolongé 3-5 jours) permet :
\begin{itemize}
\item La lixiviation (diffusion dans l'eau de trempage) des glucosides cyanogènes hydrosolubles.
\item L'activité de linamarases microbiennes (bactéries, levures, moisissures) qui hydrolysent la linamarine résiduelle $\rightarrow$ évacuation de l'HCN gazeux.
\item \textbf{Résultat :} Teneur en cyanure résiduel ramenée sous le seuil de sécurité OMS/FAO (< 10 mg/kg équivalent HCN), rendant le produit comestible sans danger aigu (konzo, goître, neuropathies).
\end{itemize}
\item \textbf{Deux avantages pour les populations :}
\begin{enumerate}
\item \textbf{Sécurité alimentaire et conservation :} Transformation d'une racine très périssable (2-3 jours après récolte) en un produit stable (bâton de manioc) conservable plusieurs semaines/mois à température ambiante. Disponibilité d'un aliment de base énergétique (glucides complexes) toute l'année, y compris en saison de soudure.
\item \textbf{Valorisation économique et sociale :} Activité génératrice de revenus pour les femmes (productrices, transformatrices, vendeuses sur les marchés). Maintien des savoir-faire traditionnels, identité culinaire (accompagnement du \og ndolé \fg{}, \og eru \fg{}, \og koki \fg{}, poisson braisé). Faible coût technologique (accessible au niveau villageois).
\end{enumerate}
\item \textbf{Limite et amélioration :}
\begin{itemize}
\item \textbf{Limite :} Procédé \textbf{artisanal, non standardisé} : durée de trempage variable, qualité de l'eau non contrôlée, hygiène précaire (risque de contamination fongique $\rightarrow$ mycotoxines : aflatoxines ; contamination bactérienne pathogène). Teneur résiduelle en cyanure imprévisible (parfois > normes). Pénibilité du travail (portage, pilage, tamisage, cuisson au feu de bois).
\item \textbf{Amélioration possible :} \textbf{Technologie « fécule / farine de manioc améliorée »} : râpage mécanique $\rightarrow$ pressage hydraulique (détoxication rapide en quelques heures) $\rightarrow$ séchage solaire ou artificiel contrôlé $\rightarrow$ farine stable, standardisée, faible teneur en cyanure, utilisable pour pains, biscuits, pâtes (projets IRAD/IITA/FAO au Cameroun). Utilisation de souches de levures/bactéries sélectionnées (starters) pour fermentation maîtrisée (goût, sécurité, rapidité).
\end{itemize}
\end{enumerate}
\end{corrige}

\begin{corrige}[Exercice 9 — Sujet type Probatoire : la brasserie et le vin de palme]
\textbf{Barème indicatif : 14 points}

\begin{enumerate}
\item \textbf{Définition et justification (2 pts) :}
\begin{itemize}
\item \textbf{Biotechnologie :} Utilisation d'organismes vivants (microorganismes) ou de leurs enzymes pour transformer des matières premières en produits utiles (définition acceptée).
\item \textbf{Justification :} La fabrication de la bière utilise des levures (\emph{Saccharomyces}) pour transformer les sucres du moût (maïs malté) en alcool et $\ce{CO2}$ $\rightarrow$ c'est un procédé biotechnologique (fermentation alcoolique), relevant de la biotechnologie traditionnelle.
\end{itemize}
\textit{Points de vigilance :} Citer le microorganisme (levure) et la transformation (sucre $\rightarrow$ alcool/$\ce{CO2}$). Ne pas confondre avec simple cuisine.

\item \textbf{Microorganisme et type de fermentation (2 pts) :}
\begin{itemize}
\item Microorganisme : \textbf{Levures} (genre \emph{Saccharomyces}, esp. \emph{S. cerevisiae}).
\item Type de fermentation : \textbf{Fermentation alcoolique} (anaérobie).
\end{itemize}

\item \textbf{Équation bilan (2 pts) :}
\[ \ce{C6H12O6 -> 2 C2H5OH + 2 CO2 + energie (ATP)} \]
\textit{Points de vigilance :} Coefficients stœchiométriques (2 devant éthanol et $\ce{CO2}$), mention « énergie » ou « ATP ». Formule de l'éthanol : $\ce{C2H5OH}$ ou $\ce{C2H6O}$.

\item \textbf{Comparaison bière / vin de palme (2 pts) :}
\begin{itemize}
\item \textbf{Affirmation CORRECTE.}
\item \textbf{Justification :} Les deux procédés reposent sur le \textbf{même phénomène biologique} : la \textbf{fermentation alcoolique} réalisée par des \textbf{levures} (\emph{Saccharomyces} spp.) qui transforment les sucres (glucose, fructose, saccharose, maltose) en \textbf{éthanol} et \textbf{dioxyde de carbone} en l'absence d'oxygène.
\item Différences : Substrat initial (moût de maïs malté vs sève de palmier), microflore spontanée vs levain sélectionné, teneur en alcool, goût, mais \textbf{même réaction biochimique fondamentale}.
\end{itemize}

\item \textbf{Calcul de la masse théorique d'éthanol (3 pts) :}
\begin{itemize}
\item Données : $m_{\text{glucose}} = \SI{360}{kg} = \SI{360000}{g}$ ; $M_{\text{glucose}} = \SI{180}{g.mol^{-1}}$ ; $M_{\text{éthanol}} = \SI{46}{g.mol^{-1}}$.
\item $n_{\text{glucose}} = \frac{360000}{180} = \SI{2000}{mol}$.
\item $n_{\text{éthanol}} = 2 \times n_{\text{glucose}} = \SI{4000}{mol}$.
\item $m_{\text{éthanol, théorique}} = 4000 \times 46 = \SI{184000}{g} = \textbf{\SI{184}{kg}}$.
\end{itemize}
\textit{Points de vigilance :} Conversion kg $\rightarrow$ g, utilisation des coefficients stœchiométriques (2), unité finale en kg.

\item \textbf{Retombées économiques et précaution d'hygiène (3 pts) :}
\begin{itemize}
\item \textbf{Retombées économiques (2 sur 3 attendues) :}
\begin{enumerate}
\item Création d'emplois locaux (brasseurs, transporteurs, vendeurs, cultivateurs de maïs).
\item Valorisation de la production locale de maïs (débouché stable, revenu pour agriculteurs, lutte contre gaspillage post-récolte).
\item Revenus fiscaux pour la commune/département (taxes, patentes).
\item Réduction des importations de boissons (devises économisées).
\item Développement de filières annexes (embouteillage, distribution, tourisme local).
\end{enumerate}
\item \textbf{Précaution d'hygiène (1 pt) :}
\begin{itemize}
\item Stérilisation / nettoyage rigoureux du matériel (cuves, cuillères, filtres, bouteilles) pour éviter contaminations (bactéries acétiques $\rightarrow$ vinaigre, moisissures, pathogènes).
\item Qualité de l'eau utilisée (potable, traitée).
\item Hygiène du personnel (lavage mains, tenues propres).
\item Contrôle de la fermentation (température, durée) pour éviter métabolites indésirables (méthanol, acide acétique en excès).
\item Pasteurisation ou conditionnement aseptique du produit fini.
\end{itemize}
\end{itemize}
\end{enumerate}
\end{corrige}

\begin{corrige}[Exercice 10 — Sujet type Probatoire : biotechnologies médicales et agricoles]
\textbf{Barème indicatif : 14 points}

\begin{enumerate}
\item \textbf{Définition et distinction (3 pts) :}
\begin{itemize}
\item \textbf{Biotechnologie moderne :} Ensemble des techniques de manipulation directe du génome (ADN recombinant, édition de génome, fusion cellulaire, culture \emph{in vitro} de cellules/tissus) pour modifier le patrimoine génétique d'un organisme ou produire des molécules d'intérêt, au-delà des barrières naturelles de reproduction.
\item \textbf{Distinction :} La biotechnologie \emph{traditionnelle} utilise des organismes entiers \emph{tels quels} (sélection naturelle/empirique, fermentation spontanée) sans modification génétique directe. La biotechnologie \emph{moderne} repose sur la \textbf{biologie moléculaire} (isolement, clonage, transfert de gènes, modification ciblée de l'ADN) $\rightarrow$ obtention d'OGM, production de protéines recombinantes, diagnostic moléculaire (PCR), sélection assistée par marqueurs (SAM).
\end{itemize}

\item \textbf{Domaines et intérêts (4 pts = 2 $\times$ 2 pts) :}
\begin{center}
\begin{tabularx}{\linewidth}{|l|X|X|}\hline
\textbf{Document} & \textbf{Domaine d'application} & \textbf{Intérêt pour les populations camerounaises} \\ \hline
\textbf{Doc 1} (Centre Pasteur : vaccins, tests diagnostic) & \textbf{Médical / Santé (Biotechnologie rouge)} & \begin{itemize}\item Production locale / conditionnement de vaccins (fièvre jaune, COVID-19, méningite, HPV) $\rightarrow$ souveraineté sanitaire, coûts réduits, chaîne du froid maîtrisée.\item Tests de diagnostic rapides (TDR paludisme, VIH, hépatites, PCR) $\rightarrow$ prise en charge précoce, surveillance épidémiologique, décentralisation (centres de santé ruraux).\item Formation de biologistes/techniciens camerounais.\end{itemize} \\ \hline
\textbf{Doc 2} (IRAD : variétés maïs/manioc améliorées) & \textbf{Agricole (Biotechnologie verte)} & \begin{itemize}\item Variétés \textbf{résistantes aux maladies} (streak du maïs, mosaïque du manioc, pourriture racinaire) $\rightarrow$ sécurité des récoltes, baisse pesticides.\item Variétés \textbf{tolérantes à la sécheresse} (adaptation changement climatique, zones soudano-sahéliennes Nord/Extrême-Nord).\item \textbf{Rendements accrus} + qualité nutritionnelle (manioc biofortifié provitamine A, maïs QPM lysine/tryptophane) $\rightarrow$ sécurité alimentaire, revenus agriculteurs, nutrition.\item Diffusion par IRAD/extensions $\rightarrow$ accessibilité semences certifiées.\end{itemize} \\ \hline
\end{tabularx}
\end{center}

\item \textbf{Principe général d'obtention d'un OGM (3 pts) :}
\begin{enumerate}
\item \textbf{Identification et isolement du gène d'intérêt} (ex: gène \emph{cry} toxine insecticide de \emph{Bt} ; gène \emph{EPSPS} tolérance glyphosate ; gène insuline humaine) $\rightarrow$ PCR, banque génomique, synthèse.
\item \textbf{Construction du cassette d'expression (vecteur) :} Assemblage \emph{in vitro} du gène + \textbf{promoteur} (constitutif ou spécifique) + \textbf{terminateur} + souvent \textbf{gène marqueur sélectionnable} (résistance antibiotique/herbicide) dans un plasmide / vecteur binaire (Agrobactère) / cassette CRISPR.
\item \textbf{Transfert (transformation) dans la cellule cible :}
\begin{itemize}
\item Plantes : \emph{Agrobacterium tumefaciens} (infection naturelle transfert T-DNA) \textbf{OU} biolistique (canon à particules d'or/ADN) \textbf{OU} électroporation de protoplastes.
\item Bactéries/Levures : Transformation (choc thermique, électroporation).
\item Animaux : Micro-injection dans pronucléus ovocyte, vecteurs viraux, CRISPR-Cas9 dans zygote.
\end{itemize}
\item \textbf{Sélection et régénération :} Culture \emph{in vitro} sur milieu sélectif (antibiotique/herbicide) $\rightarrow$ seules cellules transformées survivent $\rightarrow$ régénération plante entière (organogenèse/embryogenèse somatique) $\rightarrow$ acclimatation serre.
\item \textbf{Caractérisation moléculaire et phénotypique :} Vérification insertion (PCR, Southern blot), nombre de copies, expression (ARN, protéine), stabilité héréditaire, équivalence substantielle (sécurité).
\end{enumerate}
\textit{Points de vigilance :} Citer au moins 3 étapes logiques (gène $\rightarrow$ vecteur $\rightarrow$ transfert $\rightarrow$ sélection $\rightarrow$ vérification). Préciser \emph{Agrobacterium} ou biolistique pour les plantes.

\item \textbf{Arguments pour/contre OGM et avis argumenté (4 pts) :}
\begin{itemize}
\item \textbf{Argument POUR (exemple) :} \textbf{Sécurité alimentaire face au changement climatique.} Variétés tolérantes à la sécheresse (maïs WEMA/TELA au Kenya, Nigeria, Mozambique \emph{en essais}), résistantes aux ravageurs (maïs Bt, coton Bt au Burkina Faso, Cameroun \emph{essais confinés}) $\rightarrow$ rendements stabilisés, moins de pertes, moins d'insecticides chimiques (santé agriculteurs, environnement), revenus accrus pour petits exploitants.
\item \textbf{Argument CONTRE (exemple) :} \textbf{Risque de dépendance économique et perte de biodiversité cultivée.} Semences brevetées (droits de propriété intellectuelle) $\rightarrow$ interdiction de réutiliser grains de récolte comme semences (achat obligatoire chaque saison) $\rightarrow$ endettement paysans. Risque de contamination génétique variétés paysannes (flux de gènes) $\rightarrow$ érosion variétés locales adaptées. Résistance des insectes ravageurs au toxin Bt (évolution) $\rightarrow$ refuge non-Bt obligatoire difficile à appliquer petit champ.
\item \textbf{Avis argumenté (exemple attendu) :} \og Les biotechnologies modernes, dont les OGM, sont un \textbf{outil puissant} pour l'agriculture camerounaise face aux défis climatiques et démographiques, \textbf{à condition d'un encadrement strict} : évaluation rigoureuse cas par cas (sécurité sanitaire/environnementale), protection des variétés paysannes (droits des agriculteurs, partage des avantages), développement de semences \textbf{non brevetées / domaine public} par la recherche publique (IRAD, universités) pour l'accessibilité, et coexistence gérée avec l'agriculture conventionnelle et biologique. Une rejet systématique prive le pays d'innovations utiles ; une adoption sans régulation expose à des risques. \fg{}
\end{itemize}
\textit{Points de vigilance :} L'avis doit être \textbf{nuancé}, étayé par des faits (pas seulement « c'est bien » ou « c'est mal »), et contextualisé Cameroun (IRAD, loi sur la biosécurité 2003, décret 2012, essais confinés coton Bt/maïs Bt).
\end{enumerate}
\end{corrige}

\begin{corrige}[Exercice 11 — Essai argumenté type épreuve]
\textbf{Barème indicatif : 14 points (Intro 2, Arguments 8, Limite 2, Conclusion 2, Cohérence/Langue 2)}

\textbf{Proposition de développement structuré (environ 25 lignes) :}

\textbf{Introduction.}
La biotechnologie, définie comme l'ensemble des techniques utilisant des organismes vivants ou leurs composants pour produire des biens et services, traverse l'histoire de l'humanité : de la fermentation traditionnelle (bière, pain, miondo) au génie génétique moderne (insuline, OGM, vaccins). Au Cameroun, pays à forte biodiversité et en quête d'émergence à l'horizon 2035, la question se pose avec acuité : \og Les biotechnologies sont-elles une chance pour le développement du Cameroun ? \fg{} Nous soutiendrons qu'elles constituent un levier majeur, à condition d'une gouvernance rigoureuse.

\textbf{Arguments économiques.}
Premièrement, les biotechnologies traditionnelles et modernes dynamisent l'\textbf{agro-industrie et l'emploi}. Les brasseries (Société des Brasseries du Cameroun, brasseries artisanales bili-bili) transforment localement maïs, mil, sorgho, créant des milliers d'emplois directs et indirects (agriculteurs, transporteurs, distributeurs) et générant des recettes fiscales substantielles. La filière manioc (bobolo, miondo, farine, amidon) assure la sécurité alimentaire de millions de ménages et fait vivre une myriade de transformatrices rurales. Deuxièmement, la \textbf{biotechnologie médicale} réduit la facture sanitaire et favorise l'accès aux soins. La production locale de vaccins et tests de diagnostic au Centre Pasteur du Cameroun (fièvre jaune, COVID-19, paludisme, VIH) diminue la dépendance aux importations coûteuses, raccourcit les délais d'approvisionnement et forme des cadres de haut niveau. L'arrivée prochaine d'usines de production de médicaments essentiels (antirétroviraux, antipaludéens) par transfert de technologie renforcera cette souveraineté pharmaceutique.

\textbf{Argument social et sanitaire.}
Troisièmement, les biotechnologies améliorent directement la \textbf{santé et la nutrition des populations}. La diffusion par l'IRAD de variétés de manioc biofortifié en provitamine A et de maïs QPM (Quality Protein Maize, riche en lysine/tryptophane) combat la malnutrition infantile (avitaminose A, kwashiorkor) dans les zones vulnérables. Les tests de diagnostic rapide (TDR) au paludisme et au VIH, fruits de la biologie moléculaire, permettent une prise en charge précoce au niveau des centres de santé périphériques, sauvant des vies. La thérapie génique et les vaccins à ARN messager ouvrent des perspectives contre les maladies génétiques (drépanocytose, très prévalente au Cameroun) et les épidémies émergentes.

\textbf{Limite / Risque.}
Cependant, un \textbf{risque environnemental et socio-économique majeur} persiste : la dissémination non maîtrisée d'Organismes Génétiquement Modifiés (OGM) pourrait entraîner une contamination génétique des variétés paysannes (riz, maïs, coton), l'émergence de résistances chez les ravageurs (insectes Bt) ou les adventices (tolérance herbicide), et une dépendance des paysans aux semences brevetées des multinationales. La loi camerounaise sur la biosécurité (2003) et son décret d'application (2012) doivent être pleinement opérationnalisés (Comité National de Biosécurité, évaluations risques cas par cas, étiquetage, surveillance post-commercialisation) pour prévenir ces dérives.

\textbf{Conclusion.}
En conclusion, les biotechnologies \textbf{sont indéniablement une chance pour le développement du Cameroun} : elles valorisent nos ressources locales, créent de la richesse, soignent et nourrissent mieux. Mais cette chance ne se concrétisera que par une \textbf{appropriation nationale} : investissement massif dans la recherche publique (IRAD, universités, Centre Pasteur), formation d'experts en biosécurité et bioéthique, cadre réglementaire appliqué sans complaisance, et implication des communautés locales dans les choix technologiques. Le Cameroun a les atouts (biodiversité, jeunesse, savoir-faire) pour devenir un hub biotechnologique en Afrique centrale ; il lui faut la volonté politique de transformer l'essai.
\end{corrige}

\newpage

\coursec{Fiche bilan}

\begin{popfigure}
\begin{tikzpicture}[
  mindmap,
  concept color=popblue,
  text=white,
  font=\sffamily\small,
  level 1 concept/.append style={level distance=4.5cm, sibling angle=51.4, font=\sffamily\normalsize},
  level 2 concept/.append style={level distance=3cm, sibling angle=60, font=\sffamily\footnotesize},
  every node/.style={concept, execute at begin node=\hskip0pt},
  popblue/.style={concept color=popblue},
  popgreen/.style={concept color=popgreen},
  poporange/.style={concept color=poporange},
  poppurple/.style={concept color=poppurple},
  poppink/.style={concept color=poppink},
  popgold/.style={concept color=popgold},
  popdark/.style={concept color=popdark}
]
\node[concept, align=center]{Biotechnologie\\Production biens/services}
  child[popblue] { node[concept, align=center]{Définition\\\& Domaines}
    child { node[concept, align=center]{Utilisation\\êtres vivants} }
    child { node[concept] {Agro-alimentaire} }
    child { node[concept, align=center]{Santé\\Environnement} }
  }
  child[popgreen] { node[concept, align=center]{Fermentation\\Principes}
    child { node[concept, align=center]{Alcoolique\\$\ce{C6H12O6 -> 2C2H5OH + 2CO2}$} }
    child { node[concept, align=center]{Lactique\\$\ce{C6H12O6 -> 2CH3CHOHCOOH}$} }
    child { node[concept] {Anaérobie stricte} }
  }
  child[poporange] { node[concept, align=center]{Biotech\\Traditionnelles}
    child { node[concept] {Pain, Vin, Bière} }
    child { node[concept] {Yaourt, Fromage} }
    child { node[concept, align=center]{Produits locaux:\\foléré, bili-bili} }
  }
  child[poppurple] { node[concept, align=center]{Biotech\\Modernes}
    child[poppurple] { node[concept, align=center]{Agriculture\\OGM, Biofertilisants} }
    child[poppurple] { node[concept, align=center]{Médecine\\Insuline, Vaccins, OGM} }
    child[poppurple] { node[concept, align=center]{Environnement\\Bioremédiation} }
  }
  child[poppink] { node[concept, align=center]{Avantages\\\& Limites}
    child { node[concept] {Rendement, Qualité} }
    child { node[concept, align=center]{Risques éthiques\\\& écologiques} }
    child { node[concept] {Coût, Brevets} }
  }
  child[popgold] { node[concept, align=center]{Enjeux\\Cameroun}
    child { node[concept] {Sécurité alimentaire} }
    child { node[concept] {Emploi, Revenus} }
    child { node[concept, align=center]{Recherche locale\\IRAD, Universités} }
  }
  child[popdark] { node[concept, align=center]{Démarche\\Investigation}
    child { node[concept, align=center]{Observer\\Problématique} }
    child { node[concept, align=center]{Hypothèses\\Expérimenter} }
    child { node[concept, align=center]{Conclure\\Communiquer} }
  };
\end{tikzpicture}
\legende{Carte mentale récapitulative de la leçon NA08 : Biotechnologie et production de biens/services (1ère A).}
\end{popfigure}

\begin{bilan}
\courssub{Définitions clés}
\begin{itemize}
  \item \cle{Biotechnologie} : Ensemble des techniques utilisant des organismes vivants (micro-organismes, cellules, enzymes) ou des parties d'organismes pour produire des biens (aliments, médicaments, énergie) ou des services (dépollution, diagnostic).
  \item \cle{Fermentation} : Métabolisme anaérobie de dégradation incomplète du glucose assurant la régénération du \ce{NAD+} à partir de \ce{NADH} (pas de chaîne respiratoire, pas d'\ce{O2}, rendement énergétique faible : 2 ATP/glucose).
  \item \cle{Souche} : Population de micro-organismes d'une même espèce, génétiquement identiques, sélectionnés pour une aptitude technologique précise (vitesse, rendement, saveur).
  \item \cle{OGM (Organisme Génétiquement Modifié)} : Organisme dont le matériel génétique a été modifié par une technique de génie génétique (transgenèse, édition de génome) et non par reproduction/croisement naturel.
  \item \cle{Bioremédiation} : Utilisation de micro-organismes ou de plantes pour dégrader ou immobiliser des polluants dans l'environnement (sols, eaux).
\end{itemize}

\courssub{Principes \& Réactions à connaître par cœur}
\begin{center}
\begin{tabularx}{\linewidth}{|l|X|X|}\hline
\rowcolor{popblueL}
\textbf{Type} & \textbf{Équation bilan (mhchem)} & \textbf{Produits d'intérêt / Exemples camerounais} \\ \hline
Fermentation \textbf{alcoolique} (levures \emph{Saccharomyces}) &
$\ce{C6H12O6 -> 2 C2H5OH + 2 CO2 + 2 ATP}$ &
Éthanol (carburant, boissons : \emph{bili-bili}, \emph{kwatcha}, vin de palme), \ce{CO2} (levée pain) \\ \hline
Fermentation \textbf{lactique} (bactéries lactiques \emph{Lactobacillus, Streptococcus}) &
$\ce{C6H12O6 -> 2 CH3CHOHCOOH + 2 ATP}$ &
Acide lactique (conservation, texture) : yaourt, fromage, \emph{foléré} (feuilles d'hibiscus fermentées), \emph{soja} fermenté \\ \hline
Fermentation \textbf{acétique} (bactéries acétiques \emph{Acetobacter}, aérobie) &
$\ce{C2H5OH + O2 -> CH3COOH + H2O}$ &
Vinaigre (vin de palme, jus de fruits) \\ \hline
\end{tabularx}
\end{center}

\courssub{Méthodes express}
\begin{enumerate}
  \item \textbf{Décrire un procédé biotechnologique (fermentation) :} 1) Micro-organisme (souche) + 2) Substrat (milieu de culture : C, N, minéraux, vitamines) + 3) Conditions (T°, pH, anaérobie/aérobie, agitation) $\rightarrow$ 4) Biomasse + Métabolites d'intérêt (cible) + Déchets $\rightarrow$ 5) Séparation/Purification (aval).
  \item \textbf{Citer des exemples locaux :} Classer en \emph{traditionnels} (bili-bili, foléré, pain, yaourt artisanal, vin de palme, fromage du Nord) et \emph{modernes} (vaccins vétérinaires LANAVET, biofertilisants \emph{Rhizobium} IRAD, culture in vitro plantain/bananier, diagnostic VIH/Sida PCR).
  \item \textbf{Analyser avantages/limites :} Avantages = rendement, régularité, nouveaux produits (insuline humaine), valeur ajoutée. Limites = coût R\&D, risque dissémination OGM, dépendance semences brevetées, acceptabilité sociétale, cadre réglementaire.
\end{enumerate}
\end{bilan}

\begin{aretenir}[Checklist avant le Probatoire]
\begin{itemize}
  \item $\square$ Savoir définir \textbf{biotechnologie} et distinguer \textbf{traditionnelle} (empirique, fermentation spontanée) et \textbf{moderne} (souches sélectionnées, génie génétique, conditions contrôlées).
  \item $\square$ Écrire les \textbf{équations bilan} des fermentations alcoolique et lactique (réactifs, produits, ATP) et nommer les micro-organismes impliqués.
  \item $\square$ Expliquer le \textbf{rôle du \ce{NAD+}/\ce{NADH}} : la fermentation régénère le \ce{NAD+} pour permettre à la glycolyse de se poursuivre en l'absence d'\ce{O2}.
  \item $\square$ Citer \textbf{3 exemples traditionnels camerounais} (ex: bili-bili, foléré, vin de palme, fromage peul, pain) et identifier le micro-organisme et le type de fermentation.
  \item $\square$ Décrire le \textbf{principe d'un fermentateur industriel} (cuve stérile, agitation, aération/anaérobiose, contrôle T°/pH, capteurs, récolte).
  \item $\square$ Donner \textbf{2 applications agricoles modernes} (ex: plantes Bt résistantes insectes, tolérance herbicide, biofertilisants \emph{Rhizobium}, plants vitro sains bananier/manioc) et \textbf{2 médicales} (insuline recombinante, vaccins hépatite B/HPV, tests PCR, thérapie génique).
  \item $\square$ Différencier \textbf{OGM} (transgenèse/édition) et \textbf{non-OGM} (sélection assistée par marqueurs, mutagénèse, culture in vitro) — distinction juridique et technique.
  \item $\square$ Argumenter les \textbf{avantages} (sécurité alimentaire, santé, environnement, économie) et \textbf{limites/risques} (biodiversité, résistance antibiotiques/pesticides, éthique, brevet, dépendance paysannes) d'une biotech donnée.
  \item $\square$ Analyser les \textbf{enjeux au Cameroun} : souveraineté alimentaire (semences améliorées), santé (vaccins locaux LANAVET), emploi (filières brassicoles, transformation), recherche (IRAD, universités), cadre légal (Loi 2003 sur biotechnologie, Cartagena).
  \item $\square$ Appliquer la \textbf{démarche d'investigation} à un cas concret : ex: « Pourquoi mon yaourt maison ne prend pas ? » $\rightarrow$ Hypothèses (température, souche, lait antibiotiques) $\rightarrow$ Expérience (témoins) $\rightarrow$ Conclusion.
  \item $\square$ Maîtriser le \textbf{vocabulaire précis} : souche, inoculum, milieu de culture, amont/aval, métabolite primaire/secondaire, biorémédiation, transgenèse, CRISPR-Cas9 (notion).
  \item $\square$ Savoir \textbf{interpréter une courbe de croissance bactérienne} (phases latence, exponentielle, stationnaire, mort) et la relier à la production de biomasse vs métabolites.
\end{itemize}
\end{aretenir}

\findecours

\end{document}
